\pdfoutput=1
\documentclass[11pt]{article}
\usepackage[T1]{fontenc}
\usepackage{lmodern}
\usepackage{amsmath,amssymb,amsthm,mathtools}
\usepackage{booktabs}
% Generated by scripts/make_claims_tex.py from claims.json -- do not edit by hand.
\newcommand{\clGA}{0.1425}
\newcommand{\clqA}{1.384}
\newcommand{\clqBetaALo}{1.3835037091}
\newcommand{\clqBetaAHi}{1.3835037092}
\newcommand{\clgBetaA}{0.14195}
\newcommand{\clthmAstepOne}{2000}
\newcommand{\clthmAsub}{256}
\newcommand{\clthmAboxes}{344}
\newcommand{\clbetaClo}{1.1530164635}
\newcommand{\clbetaChi}{1.1530164636}
\newcommand{\clbetaCbound}{1.1531}
\newcommand{\clthmCcells}{6700}
\newcommand{\clthmCworst}{0.6013}
\newcommand{\clthmCtail}{2.11\cdot10^{-6}}
\newcommand{\clindepDps}{150}
\newcommand{\clindepMbig}{40}
\newcommand{\clindepCells}{938}
\newcommand{\clindepWorst}{0.6136}
\newcommand{\clindepTail}{1.1\cdot10^{-33}}
\newcommand{\clthrM}{1.103}
\newcommand{\clthrBeta}{0.930}
\newcommand{\clgLocalBeta}{1.63}
\newcommand{\clquadErr}{10^{-13}}
\newcommand{\cldirectErr}{10^{-13}}
\newcommand{\clchainSamples}{400}
\newcommand{\clrunChainTested}{3000}
\newcommand{\clrunChainW}{10^{-15}}
\newcommand{\cldivErr}{10^{-13}}
\newcommand{\clcornerSamples}{573}
\newcommand{\clnearFemSamples}{16}
\newcommand{\cllamFEMtwo}{-16.857416}
\newcommand{\cllamZeroTwo}{-16.857437}
\newcommand{\clgeomAsamples}{205\,912}
\newcommand{\clgeomAmin}{0.1445}
\newcommand{\clcellsLocal}{414}
\newcommand{\clcellsNear}{4140}
\newcommand{\clcellsFar}{109\,717}
\newcommand{\clcellsFarFiles}{24}
\newcommand{\clrLocMin}{1/4000}
\newcommand{\clrLocMax}{1/60}
\newcommand{\clRMin}{1/500}
\newcommand{\clRMax}{2/15}
\newcommand{\clmaBoxes}{1146}
\newcommand{\clpWeakBoxes}{445\,704}
\newcommand{\clpWeakLeaves}{335\,814}
\newcommand{\clpWeakW}{726}
\newcommand{\clpWeakS}{6606}
\newcommand{\clpWeakL}{328\,482}
\newcommand{\clpWeakDepth}{7}
\newcommand{\clpWeakLhs}{3.3098}
\newcommand{\clpWeakRhs}{0.1411}
\newcommand{\clpWeakFactor}{23}
\newcommand{\clpGtwo}{2.2878}
\newcommand{\clpGone}{1.1370}
\newcommand{\clpSmallThr}{0.0116}
\newcommand{\clpCell}{0.7361}
\newcommand{\clpRho}{0.2938}
\newcommand{\clpCellComp}{0.73612}
\newcommand{\clpRhoComp}{0.29376}
\newcommand{\clpCellIndep}{0.7404}
\newcommand{\clpRhoIndep}{0.2893}
\newcommand{\clpIneqC}{337}
\newcommand{\clpRingRec}{674}
\newcommand{\clpLamRec}{85}
\newcommand{\clpQmin}{1.1758}
\newcommand{\clpQminK}{123}
\newcommand{\clpQminAlpha}{2.06}
\newcommand{\clpQgeThirty}{1.48}
\newcommand{\clpQleSix}{1.5}
\newcommand{\clpQminpw}{1.1759}
\newcommand{\clpQminliu}{1.1757}
\newcommand{\clpRingN}{252}
\newcommand{\clpRingNN}{95\,634}
\newcommand{\clpLamN}{25}
\newcommand{\clpLamNN}{975}
\newcommand{\clpRingNpw}{295}
\newcommand{\clpRingNNpw}{130\,980}
\newcommand{\clpLamNpw}{30}
\newcommand{\clpLamNNpw}{1395}
\newcommand{\clpRingNliu}{160}
\newcommand{\clpRingNNliu}{38\,640}
\newcommand{\clpLamNliu}{16}
\newcommand{\clpLamNNliu}{408}
\newcommand{\clpRoots}{100}
\newcommand{\clpBsides}{36}
\newcommand{\clpCellsB}{99}
\newcommand{\clpCellsW}{1}
\newcommand{\clpThinBword}{four}
\newcommand{\clpPrisms}{1730}
\newcommand{\clpPrismspw}{1730}
\newcommand{\clpPrismsliu}{1735}
\newcommand{\clpNodes}{1779}
\newcommand{\clpNodespw}{1779}
\newcommand{\clpNodesliu}{1741}
\newcommand{\clpNmax}{184}
\newcommand{\clpNmaxpw}{216}
\newcommand{\clpNmaxliu}{117}
\newcommand{\clpNdof}{51\,060}
\newcommand{\clpNdofpw}{70\,308}
\newcommand{\clpNdofliu}{20\,709}
\newcommand{\clpKtop}{286}
\newcommand{\clpKtoppw}{286}
\newcommand{\clpKtopliu}{290}
\newcommand{\clpAtop}{2404}
\newcommand{\clpAtoppw}{2404}
\newcommand{\clpAtopliu}{2859}
\newcommand{\clpTimels}{12, 11, 3}
\newcommand{\clpTimepw}{17, 24, 5}
\newcommand{\clpTimeliu}{4, 1, 1}
\newcommand{\clpMargin}{8\cdot10^{-6}}
\newcommand{\clpEigCount}{2660}
\newcommand{\clpIndepRej}{11}
\newcommand{\clpSandLo}{4\cdot10^{-7}}
\newcommand{\clpSandHi}{1.3\cdot10^{-2}}
\newcommand{\clpSpare}{1.72}
\newcommand{\clpRaLow}{0.13}
\newcommand{\clpRaMax}{0.581}
\newcommand{\clpRaMaxAlpha}{1.9}
\newcommand{\clpRaLast}{0.484}
\newcommand{\clpHessMin}{0.87}
\newcommand{\clpHessMinAlpha}{1.3}
\newcommand{\clpHessMax}{2.07}
\newcommand{\clpNodeLo}{2}
\newcommand{\clpNodeHi}{4}
\newcommand{\clpIvNear}{7}
\newcommand{\clpIvMax}{14}
\newcommand{\clpGapLo}{0.14}
\newcommand{\clpGapHi}{0.34}
\newcommand{\clTCrowa}{1.2 & 1.3429 & 51.721 & $0.324$ & $0.832$ & $1.553$\\}
\newcommand{\clTCrowb}{1.5 & 1.6304 & 55.169 & $0.202$ & $0.500$ & $1.837$\\}
\newcommand{\clTCrowc}{2 & 2.1695 & 58.303 & $0.0751$ & $0.180$ & $2.120$\\}
\newcommand{\clTCrowd}{3 & 3.4090 & 59.855 & $0.00655$ & $0.0153$ & $2.246$\\}
\newcommand{\clTCrowe}{4 & 4.7159 & 59.989 & $4.81\cdot10^{-4}$ & $1.12\cdot10^{-3}$ & $2.250$\\}
\newcommand{\clTCrowf}{6 & 7.3471 & 59.9999 & $2.49\cdot10^{-6}$ & $5.81\cdot10^{-6}$ & $2.250$\\}

\usepackage[margin=1in]{geometry}
\usepackage[colorlinks=true,urlcolor=blue,citecolor=blue,linkcolor=blue]{hyperref}

\newtheorem{theorem}{Theorem}[section]
\newtheorem{proposition}[theorem]{Proposition}
\newtheorem{lemma}[theorem]{Lemma}
\newtheorem{corollary}[theorem]{Corollary}
\theoremstyle{remark}
\newtheorem{remark}[theorem]{Remark}

\emergencystretch=3em
\newcommand{\R}{\mathbb{R}}
\newcommand{\E}{E}
\newcommand{\one}{\mathbf 1}
\newcommand{\dd}{\,\mathrm{d}}
\newcommand{\Gp}{G_{\perp}}
\newcommand{\Bp}{B_{\perp}}
\newcommand{\Np}{N_{\perp}}
\newcommand{\cR}{\mathcal{R}}
\newcommand{\cP}{\mathcal{P}}
\DeclareMathOperator{\tr}{tr}
\DeclareMathOperator{\artanh}{artanh}
\DeclareMathOperator{\adj}{adj}
\DeclareMathOperator{\diam}{diam}

\title{The equilateral triangle is the extremal triangle for the lowest Robin eigenvalue, for every
boundary parameter}
\author{Jason Davies\thanks{\href{mailto:jason@jasondavies.com}{\nolinkurl{jason@jasondavies.com}}}}
\date{November 19, 2026}

\begin{document}

\maketitle
\begin{abstract}
Let $\lambda_1(T,\alpha)$ be the lowest eigenvalue of the Laplacian on a triangle $T$ with the Robin
condition $\partial_\nu u+\alpha u=0$. Lotoreichik and Lu conjectured that, among triangles of given
area, the equilateral triangle minimises $\lambda_1(T,\alpha)$ when $\alpha>0$ and maximises it when
$\alpha<0$. We prove both halves, with the equilateral triangle as the unique extremiser; the same holds
at fixed perimeter. The proof is computer-assisted.

For $\alpha>0$ the result is new; previously the inequality was known for all triangles only when
$\alpha\sqrt{|T|}$ is smaller than an absolute but non-explicit constant (Ferudun). Thin triangles are
handled by slicing, and $\alpha\sqrt{|T|}\le1/100$ by an explicit vector field and a Schur-complement
criterion. Near the equilateral triangle $E$, the second-order perturbation constant of $E$ is bounded
from above through certified \emph{lower} bounds for two perturbed eigenvalue problems, uniformly in
$\alpha$. The remaining compact set of shapes is covered by $\clpRoots$ cells; an exact identity in the
Klein model of the hyperbolic plane of shapes, together with concavity, reduces each cell to certified
Crouzeix--Raviart lower bounds at its vertices, which are checked in exact integer and ball arithmetic.

For $\alpha<0$ the first proof is due to Ma (computer-assisted, 2026); ours was found independently.
Write $\beta=-\alpha\sqrt{|T|}$. For $\beta\le6/5$ two trial functions suffice. For
$\beta\ge\clbetaCbound$ an exact corner bound disposes of all triangles with a small angle, and on the
remaining near-equilateral triangles three transplanted corner waves, with a one-step perturbation
estimate, reduce the problem to a one-variable inequality, which is verified in ball arithmetic.
\end{abstract}

\noindent\textbf{AI transparency statement.} The results, their proofs, the code for the certified
computations and this manuscript were produced by artificial intelligence. The human author initiated
the broader research and directed the AI systems, but did not write the text and has not checked the
arguments. Details are given at the end of the paper.

\section{Introduction}\label{sec:intro}

For a bounded Lipschitz domain $\Omega\subset\R^2$ and $\alpha\in\R$ let
\begin{equation}\label{eq:var}
  \lambda_1(\Omega,\alpha)=\inf_{0\neq u\in H^1(\Omega)}
  \frac{\int_\Omega|\nabla u|^2+\alpha\int_{\partial\Omega}u^2}{\int_\Omega u^2}
\end{equation}
be the lowest eigenvalue of $-\Delta$ with the Robin condition $\partial_\nu u+\alpha u=0$, where $\nu$
is the outer normal. For $\alpha>0$ the Bossel--Daners inequality~\cite{Bossel,Daners} states that the
disc minimises $\lambda_1(\cdot,\alpha)$ among domains of given area. For $\alpha<0$ we have
$\lambda_1<0$, and Bareket's conjecture~\cite{Bareket} that the disc maximises $\lambda_1(\cdot,\alpha)$
among simply connected planar domains of given area is open (see~\cite{AFK,KLV} for its history). The
polygonal analogue for triangles is the following.

\begin{quote}
\textbf{Conjecture} (Lotoreichik--Lu; AIM problem list~\cite{AimPL}, Conjecture~1.3; see
also~\cite{LaSi,KLV}). Among all triangles of a given area, $\lambda_1(T,\alpha)$ is minimised by the
equilateral triangle when $\alpha>0$ and maximised by it when $\alpha<0$.
\end{quote}

We prove the conjecture. By scaling, $\lambda_1(sT,\alpha)=s^{-2}\lambda_1(T,s\alpha)$ for $s>0$, so
only the dimensionless coupling $\alpha\sqrt{|T|}$ matters.

\begin{theorem}\label{thm:main}
Let $T$ be a triangle and let $\E_T$ be the equilateral triangle with $|\E_T|=|T|$.
\begin{enumerate}
\item[\textup{(1)}] If $\alpha>0$, then $\lambda_1(T,\alpha)\ge\lambda_1(\E_T,\alpha)$.
\item[\textup{(2)}] If $\alpha<0$, then $\lambda_1(T,\alpha)\le\lambda_1(\E_T,\alpha)$.
\end{enumerate}
In both cases equality holds only if $T$ is equilateral.
\end{theorem}

\begin{corollary}[fixed perimeter]\label{cor:perimeter}
Let $T$ be a triangle and let $\E^T$ be the equilateral triangle with the same perimeter as $T$. Then
$\lambda_1(T,\alpha)\ge\lambda_1(\E^T,\alpha)$ if $\alpha>0$ and $\lambda_1(T,\alpha)\le\lambda_1(\E^T,\alpha)$ if
$\alpha<0$. In both cases equality holds only if $T$ is equilateral.
\end{corollary}

As $\alpha\to+\infty$, Theorem~\ref{thm:main}(1) becomes the P\'olya--Szeg\H{o}
theorem for the Dirichlet eigenvalue~\cite{PS}, which is proved by Steiner symmetrisation; as
$\alpha\to0^\pm$ both parts reduce to the isoperimetric inequality for triangles, because
$\lambda_1(T,\alpha)=\alpha P/|T|+O(\alpha^2)$, with $P$ the perimeter. Steiner symmetrisation does not
control the boundary term in~\eqref{eq:var}.

\paragraph{Priority and related work.} The case $\alpha<0$ was first proved by Yicen Ma~\cite{Ma}, in a
computer-assisted preprint deposited on Zenodo on 6 October 2026, which has not been refereed. The code
and certificate files that the preprint describes are not part of the deposit; an AI referee working
for us re-implemented its middle-range criterion in Arb ball arithmetic and confirmed it on
$1\le\beta\le10$ ($\clmaBoxes$ boxes; the script is in our repository, labelled as a check of prior
work). We found the proof in Part~I independently and learned of~\cite{Ma} only after it was complete;
Part~I is therefore an independent second proof of Theorem~\ref{thm:main}(2). The two proofs share
several ingredients: the trial functions for small coupling and the identity of
Lemma~\ref{lem:transplant} (which is (5) of~\cite{Ma}), the exact corner bound of
Lemma~\ref{lem:corner} (which is (7)--(8) of~\cite{Ma}), the same threshold angle separating
near-equilateral from far triangles, and the convexity bound for the strain
(Lemma~\ref{lem:geomC}). They differ in the treatment of near-equilateral triangles. Ma uses the span of
$u$, $xu$, $yu$ ($u$ the equilateral ground state), the trial function $(1-\tau\,v\cdot x)u$ and Taylor
bounds in the strain, certifies $1\le\beta\le10$ by interval arithmetic and treats $\beta\le1$ and
$\beta\ge10$ analytically. We use the three transplanted corner waves, exact side-ratio identities and an
estimate of the tunnelling gap, certify a single one-variable inequality for all
$\beta\ge\clbetaCbound$, and use the two-function argument of Theorem~\ref{thm:A} for $\beta\le6/5$.
The perimeter version of the case $\alpha<0$ (Corollary~\ref{cor:perimeter}), a second,
independently written certificate and public code are specific to Part~I. Ma does not treat $\alpha>0$.

For $\alpha>0$, Ferudun~\cite{Ferudun} proved, in an unrefereed AI-assisted preprint, the comparison of
Theorem~\ref{thm:main} for $0<|\alpha|\sqrt{|T|}<\varepsilon$, with an absolute but non-explicit
$\varepsilon$; for $\alpha>0$ this is the only earlier result covering all triangles that we know of.
Krej\v{c}i\v{r}\'{i}k, Lotoreichik and Vu~\cite{KLV} proved, for $\alpha<0$, that the equilateral
triangle is a strict local maximiser for $\beta\le0.92$ (numerically up to about $1.63$), and that the
comparison holds for each fixed triangle when $\beta$ is sufficiently small or sufficiently large, with
thresholds depending on the triangle, and uniformly for small $\beta$ on compact sets of shapes.
Carbotti, Cito and Pallara~\cite{CCP} prove the existence of optimal generalised polygons without
identifying them. For the Robin \emph{torsional rigidity}, Gui, Hu, Li and Yang~\cite{GHLY} proved that
the equilateral triangle is the maximiser among triangles of given area for every positive boundary
parameter, by cutting the polygon from its incentre into cells with Neumann conditions on the cuts;
torsional rigidity is additive over the cells, whereas for the eigenvalue this relaxation gives only the
smallest cell eigenvalue, which is not enough. For the Dirichlet problem, Endo, Liu and
Mariano~\cite{ELM} proved sharp inequalities on triangles with certified finite elements; part of our
strategy for $\alpha>0$ (Crouzeix--Raviart lower bounds~\cite{CR,Liu,CG}) is of the same type.

\paragraph{The case $\alpha<0$ (Part~I).} Write $\beta=-\alpha\sqrt{|T|}>0$, let $\E$ be the equilateral
triangle of area~$1$ and let $\beta(M)$ be the increasing parametrisation of the coupling introduced in
\S\ref{sec:prelim}; put $\beta_C=\beta(13/10)$, so that $\clbetaClo\le\beta_C\le\clbetaChi$
(Lemma~\ref{lem:certC}). Theorem~\ref{thm:main}(2) follows from two overlapping results.

\begin{theorem}\label{thm:A}
If $0<\beta\le6/5$, then $\lambda_1(T,-\beta)<\lambda_1(\E,-\beta)$ for every non-equilateral triangle $T$
of area~$1$.
\end{theorem}

\begin{theorem}\label{thm:C}
If $\beta\ge\beta_C$, then $\lambda_1(T,-\beta)<\lambda_1(\E,-\beta)$ for every non-equilateral triangle
$T$ of area~$1$.
\end{theorem}

The proof of Theorem~\ref{thm:A} (Section~\ref{sec:A}) uses only the constant function and the affinely
transplanted ground state of $\E$, together with a two-variable inequality checked in interval
arithmetic. The transplanted ground state alone cannot work for large $\beta$: its second variation at
$\E$ changes sign at $\beta\approx\clgLocalBeta$ (Remark~\ref{rem:A}). The proof of
Theorem~\ref{thm:C} (Section~\ref{sec:C}) has two parts.
\begin{itemize}
\item \emph{A corner bound.} For every triangle, every angle $\theta$ and every $\beta>0$,
$\lambda_1(T,-\beta)<-\beta^2/\sin^2(\theta/2)$ (Lemma~\ref{lem:corner}). Since
$\lambda_1(\E,-\beta)=-4\beta^2/\tanh^2M$, this settles every triangle whose smallest angle $\theta$
satisfies $\sin(\theta/2)\le\frac12\tanh M$. What remains is a neighbourhood of $\E$ of size about
$1-\tanh M\approx6e^{-2\cdot3^{1/4}\beta}$.
\item \emph{Three corner waves.} The ground state of $\E$ is a sum of three functions $\varphi_i$, each a
plane wave concentrated at one corner plus its reflection in the opposite side. On a near-equilateral
triangle we use the span of their affine transplants. A deformation couples the transplanted ground
state to the two-dimensional complement, which lies only a tunnelling splitting of order $1-\tanh M$
above the ground state. A single step of second-order perturbation theory then gains more than the
deformation costs, by a factor of order $(1-\tanh M)^{-1}$. The resulting criterion reduces to a
one-variable inequality in $M$, which we verify in ball arithmetic for all $M\ge13/10$
(Lemma~\ref{lem:certC}).
\end{itemize}

\paragraph{The case $\alpha>0$ (Part~II).} Here we need \emph{lower} bounds for $\lambda_1(T,\alpha)$,
which no trial function provides. With $|T|=1$, the shapes of triangles form the hyperbolic plane, and we
use its Klein model: a point $w$ of the unit disc, with $w=0$ for $\E$. The parameter set
$\{(w,\alpha)\}$ is covered by four regions (Section~\ref{p-sec:overview}): thin triangles, for all
$\alpha$ (Theorem~\ref{p-thm:thin}, by slicing); weak coupling $\alpha\le1/100$, for all shapes
(Theorem~\ref{p-thm:uniform}); a neighbourhood $|w|\le0.1194$ of $\E$, for all $\alpha\ge1/100$
(Proposition~\ref{p-prop:local}); and the remaining compact middle region, for all $\alpha\ge1/100$
(Proposition~\ref{p-prop:middle}). Near $\E$ a Schur-complement criterion needs an \emph{upper} bound
for a second-order perturbation constant $R_1(\alpha)$ of $\E$, accurate to a few per cent and uniform
in $\alpha$; we obtain it from certified lower bounds for the eigenvalues of two perturbed problems at
distance $1/10$ from $\E$ in opposite directions, for which the first-order errors cancel
(Lemma~\ref{p-lem:ring}). In the middle region, the lowest eigenvalue is a concave function of the
coefficients of the quadratic form pulled back to $\E$; convex combinations of these coefficients are
exactly the Euclidean triangles of the Klein model, up to an explicit factor
(Lemma~\ref{p-lem:cell}). Certified Crouzeix--Raviart lower bounds at the vertices of $\clpRoots$ cells
and at grid points in $\alpha$ then cover the whole region.

\paragraph{Organisation.} Part~I (Sections~\ref{sec:prelim}--\ref{sec:cross}) proves
Theorem~\ref{thm:main}(2) and Corollary~\ref{cor:perimeter} for $\alpha<0$, and Part~II
(Sections~\ref{p-sec:overview}--\ref{p-sec:proof}) proves Theorem~\ref{thm:main}(1) and
Corollary~\ref{cor:perimeter} for $\alpha>0$. The two parts are
independent of each other, and each has its own notation. Section~\ref{sec:verif} describes the
implementation, the verification and the remaining caveats. Appendix~\ref{p-sec:large} contains a
strong-coupling bound for individual triangles that is not used in the proof.

\part{The case \texorpdfstring{$\alpha<0$}{alpha < 0}}
\section{Preliminaries}\label{sec:prelim}

In this part $\alpha<0$, and the notation is independent of Part~II. Throughout, $|T|=1$ by scaling,
and $\alpha=-\beta<0$. Each map
$\alpha\mapsto\int|\nabla u|^2+\alpha\int_{\partial T}u^2$ is affine. By~\eqref{eq:var} we therefore have
the following.

\begin{lemma}\label{lem:concave}
For fixed $T$, the function $\alpha\mapsto\lambda_1(T,\alpha)$ is concave and nondecreasing on $\R$, and
$\lambda_1(T,0)=0$. Consequently $\beta\mapsto|\lambda_1(T,-\beta)|/\beta$ is nondecreasing on
$(0,\infty)$. For each fixed $u\in H^1(T)\setminus\{0\}$, the Rayleigh quotient $R_\alpha[u]$ is affine
in $\alpha$, and $R_\alpha[u]\ge\lambda_1(T,\alpha)$, with equality only if $u$ is a ground state.
\end{lemma}

\begin{proof}
An infimum of nondecreasing affine functions is concave and nondecreasing. If $f$ is concave with
$f(0)=0$ and $\alpha<\alpha'<0$, then $f(\alpha')\ge(\alpha'/\alpha)f(\alpha)$, so
$f(\alpha')/\alpha'\le f(\alpha)/\alpha$. The last statement is the variational principle.
\end{proof}

\subsection{The equilateral triangle}

Let $\E$ be the equilateral triangle of area $1$. Its height is $H=3^{1/4}$, its side length
$\ell=2/H$, its inradius $r=H/3$ and its perimeter $P_*=6/H=\sqrt{12\sqrt3}$. Let $d_i(x)$ be the
distance from $x$ to the $i$th side. For $M>0$ put
\begin{equation}\label{eq:param}
  t=\tanh M,\quad L=-\artanh(t/2),\quad K=M-L,\quad k=\frac{2K}{H},\quad
  \beta(M)=\frac{tK}{H},\quad \lambda_0(M)=-\frac{4K^2}{H^2}.
\end{equation}

\begin{proposition}[{cf.~\cite{McCartin,KLV}}]\label{prop:equi}
The map $M\mapsto\beta(M)$ is a smooth increasing bijection $(0,\infty)\to(0,\infty)$. For
$\beta=\beta(M)$, the function
\[
  \psi_0=\sum_{i=0}^2\varphi_i,\qquad \varphi_i=\cosh(L+k\,d_i),
\]
is a positive ground state of the Robin Laplacian on $\E$ with $\alpha=-\beta$, and
$\lambda_1(\E,-\beta)=\lambda_0(M)=-4\beta^2/\tanh^2M$. If $n_i$ is the inward unit normal of side~$i$,
then $\psi_0$ is a combination of the six plane waves $e^{\pm k n_i\cdot x}$.
\end{proposition}

\begin{proof}
The functions $t$, $K$, and hence $\beta=tK/H$, are increasing in $M$, with limits $0$ and $\infty$. Since
$|\nabla d_i|=1$, each $\varphi_i$ satisfies $\Delta\varphi_i=k^2\varphi_i$, so $-\Delta\psi_0=\lambda_0\psi_0$.
On side~$0$ we have $d_0=0$, $\partial_\nu d_0=-1$, $\partial_\nu d_1=\partial_\nu d_2=\frac12$ (the
normals meet at $120^\circ$) and $d_1+d_2=H$. Write $y=kd_1-K$, so that $L+kd_1=M+y$ and $L+kd_2=M-y$.
On side~$0$,
\[
  \psi_0=\cosh L+2\cosh M\cosh y,\qquad
  \partial_\nu\psi_0=k\bigl(-\sinh L+\sinh M\cosh y\bigr).
\]
Hence $\partial_\nu\psi_0-\beta\psi_0=0$ there, for all $y$, if and only if $k\sinh L=-\beta\cosh L$ and
$k\sinh M=2\beta\cosh M$, that is, $\tanh L=-\beta/k$ and $\tanh M=2\beta/k$. Both hold
by~\eqref{eq:param}. By symmetry the same holds on the other sides. A positive eigenfunction is a ground
state, because eigenfunctions with different eigenvalues are $L^2$-orthogonal and the ground state does
not change sign. Finally $\lambda_0=-k^2=-4\beta^2/t^2$. With the centroid at the origin we have
$d_i=n_i\cdot x+r$, so $\varphi_i$ is a combination of $e^{\pm kn_i\cdot x}$.
\end{proof}

Let $N=\int_\E\psi_0^2$, $\mathcal G=\int_\E|\nabla\psi_0|^2$ and $B=\int_{\partial\E}\psi_0^2$, so that
$\lambda_0=(\mathcal G-\beta B)/N$.

\begin{lemma}\label{lem:HF}
Let $K_M=dK/dM=1+\frac{2(1-t^2)}{4-t^2}$. Then
\[
  \frac BN=\frac{d\lambda_1(\E,\alpha)}{d\alpha}=\frac{8KK_M}{H\,[(1-t^2)K+tK_M]},\qquad
  \frac{\mathcal G}N=\lambda_0+\beta\frac BN,
\]
\[
  g:=\frac{\mathcal G/N}{2\beta B/N}=\frac14-\frac{K}{2\sinh(2M)K_M}\le\frac{\tanh^2M}{4}.
\]
Moreover $K$ is increasing in $M$ and $K_M$ is decreasing.
\end{lemma}

\begin{proof}
For every $\alpha'$ we have $\lambda_1(\E,\alpha')\le R_{\alpha'}[\psi_0]=\lambda_0+(\alpha'-\alpha)B/N$.
So $B/N$ is a supergradient at $\alpha$ of the concave function $\lambda_1(\E,\cdot)$. By
Proposition~\ref{prop:equi} that function is a smooth function of $M$, and $d\beta/dM>0$, so it is
differentiable. Hence $B/N$ equals the derivative $(d\lambda_0/dM)/(d\alpha/dM)$, which gives the stated
formula. The identity for $\mathcal G/N$ is $\lambda_0=(\mathcal G-\beta B)/N$.

Substituting gives $g=\frac14-\frac{(1-t^2)K}{4tK_M}$, and $(1-t^2)/t=2/\sinh 2M$. For the bound we use
$K\ge M+\frac12\tanh M\ge\frac32\tanh M$ and $K_M\le\frac32$. These give
$K/(\sinh 2M\,K_M)\ge1/(2\cosh^2M)$, hence $g\le\frac14-\frac1{4\cosh^2M}=\frac14\tanh^2M$. Finally $K_M$
decreases because $s\mapsto(1-s)/(4-s)$ decreases on $[0,1]$.
\end{proof}

\subsection{Affine transplantation}

Let $\E$ have its centroid at the origin and vertices $V_0=(-\ell/2,-r)$, $V_1=(\ell/2,-r)$, $V_2=(0,2r)$.
Side~$s$ joins $V_s$ to $V_{s+1}$ and has unit direction $e_s$. For $a>0$ and $b\in\R$ let
\[
  D_{a,b}=\begin{pmatrix}a&b\\0&1/a\end{pmatrix},\qquad T_{a,b}=D_{a,b}\E .
\]

\begin{lemma}\label{lem:shape}
Every triangle is similar to exactly one $T_{a,b}$ with $(a,b)$ in
\[
  \mathcal D=\Bigl\{a\ge1,\ -\tfrac{a}{\sqrt3}\Bigl(\sqrt{4-3a^{-4}}-1\Bigr)\le b\le0\Bigr\}.
\]
The equilateral triangle corresponds to $(1,0)$. If $T_{a,b}$ has perimeter $P$, then $P/P_*\ge 2a/3$.
\end{lemma}

\begin{proof}
Normalise a longest side to $[0,1]$ and the apex to $(x,y)$ with $0\le x\le\frac12$ and
$(1-x)^2+y^2\le1$, then scale by $\sigma=\sqrt{2/y}$ to area $1$. The linear map that takes the edge
vectors $V_1-V_0$, $V_2-V_0$ of $\E$ to $(\sigma,0)$ and $(\sigma x,\sigma y)$ is $D_{a,b}$ with
$a=\sigma/\ell$ and $b=\sigma(x-\frac12)/H$. A direct computation turns the constraints on $(x,y)$ into
$(a,b)\in\mathcal D$. Finally $P=\sigma(1+|(x,y)|+|(1-x,y)|)\ge2\sigma=2a\ell$ and $P_*=3\ell$.
\end{proof}

Let $D$ be a real $2\times2$ matrix with $\det D=1$, and let $T=D\E+x_0$. For $u\in H^1(\E)$ let
$\tilde u(x)=u(D^{-1}(x-x_0))$. Then
\begin{equation}\label{eq:transplant}
  \int_{T}|\nabla\tilde u|^2=\int_\E\nabla u^{\top}\mathcal A\,\nabla u,\qquad
  \int_T\tilde u^2=\int_\E u^2,\qquad
  \int_{\partial T}\tilde u^2=\sum_{s=0}^2\rho_s\int_{\text{side }s}u^2,
\end{equation}
where $\mathcal A=D^{-1}D^{-\top}$ and $\rho_s=|De_s|$ is the ratio of the length of side~$s$ of $T$ to
$\ell$. For $D=D_{a,b}$ we have
$\mathcal A=\left(\begin{smallmatrix}a^{-2}+b^2&-ab\\-ab&a^2\end{smallmatrix}\right)$.

\begin{lemma}[transplanted ground state]\label{lem:transplant}
Let $T$ be a triangle of area $1$ with sides $a_0,a_1,a_2$. Let $p=(a_0+a_1+a_2)/P_*\ge1$ and
$\tau=(a_0^2+a_1^2+a_2^2)/(2\sqrt3)\ge2$. Then
\[
  \lambda_1(T,-\beta)\le\frac{\mathcal G}N\cdot\frac\tau2-\beta\frac BN\,p
  =\lambda_0+\frac{\mathcal G}N\Bigl(\frac\tau2-1\Bigr)-\beta\frac BN(p-1).
\]
\end{lemma}

\begin{proof}
Take $u=\psi_0$ in~\eqref{eq:transplant}. The matrix $\int_\E\nabla\psi_0\nabla\psi_0^{\top}$ commutes with the
rotation by $120^\circ$, so it equals $\frac{\mathcal G}2I$. Hence the gradient term is
$\frac{\mathcal G}2\tr\mathcal A$. For a $2\times2$ matrix of determinant one,
$\tr(D^{-1}D^{-\top})=\tr(D^\top D)$. Since $\sum_se_se_s^\top=\frac32I$, we get
$\sum_s a_s^2=\ell^2\sum_s|De_s|^2=\frac32\ell^2\tr(D^\top D)$. As $\frac32\ell^2=2\sqrt3$, this gives
$\tr\mathcal A=\tau$. By symmetry each side of $\E$ carries $B/3$, so the boundary term is
$\frac B3\sum_s a_s/\ell=Bp$. Finally, $\tau\ge2$ because $\tr(D^\top D)\ge2\det D$, and $p\ge1$ is the
isoperimetric inequality for triangles.
\end{proof}

\section{Small coupling: proof of Theorem~\ref{thm:A}}\label{sec:A}

Set $\beta_A=6/5$, $G_A=\clGA$ and $q_A=\clqA$.

\begin{lemma}\label{lem:gq}
For $0<\beta\le\beta_A$: (i) $g(\beta)<G_A$; (ii)
$q(\beta):=|\lambda_0(\beta)|/(\beta P_*)\le q(\beta_A)<q_A$.
\end{lemma}

\begin{proof}
(ii) The first inequality is Lemma~\ref{lem:concave}. For the second we enclose $M(\beta_A)$ by
bisection in ball arithmetic. This gives $\clqBetaALo\le q(\beta_A)\le\clqBetaAHi$.

(i) We cover $[0,M(\beta_A)]$ by $\clthmAstepOne$ intervals $[m_1,m_2]$. On the first interval we use
$g\le\tanh^2(m_2)/4$. On the others we use $g\le\frac14-K(m_1)/(2\sinh(2m_2)K_M(m_1))$, which is valid by
the monotonicity in Lemma~\ref{lem:HF}. Every bound is below $G_A$. The largest value of $g$ on this
range is $g(\beta_A)=\clgBetaA\ldots$.
\end{proof}

\begin{lemma}\label{lem:geomA}
Let $T$ be a non-equilateral triangle of area $1$ with $p(T)\le q_A$. Then
$\Phi(T):=(p-1)-G_A(\tau-2)>0$.
\end{lemma}

\begin{proof}
We use the apex coordinates $(x,y)$ of the proof of Lemma~\ref{lem:shape}; in them $p$ and $\tau$ are
explicit algebraic functions. If $y<19/100$, then $p\ge2\sqrt{2/y}/P_*>q_A$. The equilateral apex
$z_0=(\frac12,\frac{\sqrt3}2)$ minimises both $p$ and $\tau$ over all apex positions. Hence
$\Phi(z_0)=0$ and $\nabla\Phi(z_0)=0$. On the box $|x-\frac12|,|y-\frac{\sqrt3}2|\le\frac1{40}$ we evaluate
the Hessian of $\Phi$ in interval arithmetic, by second-order automatic differentiation in ball
arithmetic on $\clthmAsub$ sub-boxes. It is positive definite on every sub-box. Hence
$\Phi(z)=\int_0^1(1-s)(z-z_0)^\top\nabla^2\Phi(z_0+s(z-z_0))(z-z_0)\dd s>0$ for $z\ne z_0$ in the box.
An adaptive subdivision of $[0,\frac12]\times[\frac{19}{100},\frac{87}{100}]$ ends in $\clthmAboxes$ boxes.
For each box, one of the following holds:
\begin{itemize}
\item it lies inside the local box;
\item interval arithmetic or a mean-value form shows that it lies outside the fundamental domain;
\item it shows that $p>q_A$ on the box;
\item it shows that $\Phi>0$ on the box.
\end{itemize}
\end{proof}

\begin{proof}[Proof of Theorem~\ref{thm:A}]
Let $0<\beta\le6/5$ and let $T$, of area $1$, be non-equilateral. First suppose $p(T)\ge q(\beta)$. The
constant trial function gives $\lambda_1(T,-\beta)<R[1]=-\beta P_*p\le\lambda_0(\beta)$. The first
inequality is strict because constants are not eigenfunctions when $\alpha\neq0$.

Otherwise $p<q(\beta)<q_A$ by Lemma~\ref{lem:gq}(ii), so $\Phi(T)>0$ by Lemma~\ref{lem:geomA}. Since
$\tau>2$, Lemma~\ref{lem:gq}(i) gives $p-1>G_A(\tau-2)>g(\tau-2)$. Lemma~\ref{lem:transplant} then gives
$\lambda_1(T,-\beta)-\lambda_0\le\beta\frac BN\bigl(g(\tau-2)-(p-1)\bigr)<0$.
\end{proof}

\begin{remark}\label{rem:A}
The transplanted ground state alone works near $\E$ exactly when $g(\beta)$ is below
$\lim_{T\to\E}(p-1)/(\tau-2)$. Write the side ratios as $\rho_i=1+\delta_i+\bar\delta$ with
$\sum\delta_i=0$. The area constraint (Lemma~\ref{lem:rho}(c) below) gives
$p-1\sim\frac12\sum\delta_i^2$ and $\tau-2\sim\frac83\sum\delta_i^2$, so the limit is $3/16=0.1875$ in
every direction. In floating point, $g(\beta)=3/16$ at $\beta\approx\clgLocalBeta$, and $g(\beta)<3/16$
below that value, so the local limit of this argument is $\beta\approx\clgLocalBeta$. This matches the numerical threshold of~\cite{KLV}.
Replacing the constant by the disc bound $\lambda_1(T,\alpha)\le\lambda_1(D_{r_T},\alpha)$, which holds
for tangential polygons by the method of parallel coordinates ($r_T$ the inradius), would raise $6/5$ to
about $1.29$. This is a floating-point estimate that we have not certified, and it is not needed.
\end{remark}

\section{Large coupling: proof of Theorem~\ref{thm:C}}\label{sec:C}

\subsection{The corner bound}

\begin{lemma}\label{lem:corner}
Let $T$ be a triangle of any area, let $\theta$ be one of its angles, and let $\beta>0$. Then
\[
  \lambda_1(T,-\beta)<-\frac{\beta^2}{\sin^2(\theta/2)} .
\]
\end{lemma}

\begin{proof}
Let $V$ be the vertex with angle $\theta$, and let $e$ be the unit vector along the bisector at $V$,
pointing into $T$. For $\kappa>0$ put $w=-\kappa e$ and $u(x)=e^{w\cdot(x-V)}$. Then
$|\nabla u|^2=\kappa^2u^2$ and $\operatorname{div}(w\,u^2)=2\kappa^2u^2$. By the divergence theorem,
\[
  2\kappa^2\int_Tu^2=\sum_{s}(w\cdot\nu_s)\int_{\text{side }s}u^2 ,
\]
where $\nu_s$ is the outer unit normal of side $s$. On the two sides through $V$ we have
$e\cdot\nu_s=-\sin(\theta/2)$, so $w\cdot\nu_s=\kappa\sin(\theta/2)$. On the opposite side
$e\cdot\nu_s>0$, because the bisector leaves $T$ through that side. Hence $w\cdot\nu_s<0<\kappa\sin(\theta/2)$
there. Since $u>0$,
\[
  2\kappa^2\int_Tu^2<\kappa\sin(\theta/2)\int_{\partial T}u^2 .
\]
So the Rayleigh quotient of $u$ is
$\kappa^2-\beta\int_{\partial T}u^2/\int_Tu^2<\kappa^2-2\beta\kappa/\sin(\theta/2)$. Now take
$\kappa=\beta/\sin(\theta/2)$.
\end{proof}

The trial function is the ground state $e^{-\kappa e\cdot(x-V)}$ of the infinite
sector~\cite[Example~2.5, Lemma~2.6]{LP}, see also~\cite{KP}. It was used for triangles
in~\cite[Theorem~5.3]{KLV}, where the bound $\lambda_1(T,\alpha)\le-\frac{\alpha^2}{\sin^2(\theta/2)}
\bigl(1-2e^{2\alpha L'\cot(\theta/2)}\bigr)$ is obtained ($L'$ the shorter adjacent side), with an
exponentially small loss; the divergence identity removes that loss. Lemma~\ref{lem:corner} also appears,
with the same proof, as (7)--(8) in Ma's preprint~\cite{Ma}. It is the corner analogue of the bound
$\lambda_1(\Omega,-\beta)\le-\beta^2$ for smooth domains, which~\cite{BP} attribute
to~\cite[Theorem~2.1]{GS}. The asymptotics $\lambda_1(T,-\beta)\sim-\beta^2/\sin^2(\theta_{\min}/2)$ as
$\beta\to\infty$ go back to Lacey, Ockendon and Sabina~\cite{LOS} and Levitin and Parnovski~\cite{LP};
exponentially precise asymptotics are in~\cite{Khalile,HP}. We could not consult~\cite{LOS} or the
survey~\cite{BFK}.

\begin{corollary}\label{cor:far}
Let $|T|=1$, $\beta=\beta(M)$ and $t=\tanh M$, and let $\theta_{\min}$ be the smallest angle of $T$. If
$\sin(\theta_{\min}/2)\le t/2$, then $\lambda_1(T,-\beta)<\lambda_0(M)$.
\end{corollary}

\begin{proof}
By Lemma~\ref{lem:corner} and Proposition~\ref{prop:equi},
$\lambda_1(T,-\beta)<-\beta^2/\sin^2(\theta_{\min}/2)\le-4\beta^2/t^2=\lambda_0(M)$.
\end{proof}

It remains to treat the \emph{near region} $\theta_{\min}>\theta(t)$, where
\begin{equation}\label{eq:near}
  \theta(t):=2\arcsin(t/2)\in(0,\pi/3),\qquad \pi/3-\theta(t)\sim\tfrac{2}{\sqrt3}(1-t)\quad(t\to1).
\end{equation}

\subsection{Three trial functions}

Let $T$ be a triangle of area $1$ with sides $a_0,a_1,a_2$, numbered counterclockwise like those of $\E$,
and let $\theta_i$ be the angle opposite side~$i$. Put
\[
  \rho_i=\frac{a_i}{\ell},\qquad S=\rho_0^2+\rho_1^2+\rho_2^2,\qquad
  W=\sum_i(\rho_i-\bar\rho)^2,\qquad \bar\rho=\tfrac13\sum_i\rho_i ,
\]
where $\ell=2/H$ is the side length of $\E$. Let $x\mapsto Dx+x_0$ be the affine map of $\E$ onto $T$ that
takes side~$i$ to side~$i$. Then $\det D=1$ and $\rho_i=|De_i|$, and~\eqref{eq:transplant} holds with
$\mathcal A=D^{-1}D^{-\top}$. Let $\mu\ge1$ be the larger eigenvalue of $\mathcal A$. Our trial functions
are the transplants $\tilde\varphi_i$ of the three summands $\varphi_i=\cosh(L+kd_i)$ of $\psi_0$. On $T$
they are the functions $\cosh(L+2K\,b_i)$, where $b_0,b_1,b_2$ are the barycentric coordinates of $T$
and $b_i$ vanishes on side~$i$. Each $\varphi_i$ is the sum of the plane wave that decays away from the
vertex opposite side~$i$ and its reflection in side~$i$. In this section $\{i,j,m\}=\{0,1,2\}$ always
denotes the three distinct indices.

\begin{lemma}\label{lem:rho}
Let $n_i$ be the inward unit normal of side~$i$ of $\E$.
\begin{enumerate}
\item[(a)] $n_i^\top\mathcal An_i=\rho_i^2$ and $n_i^\top\mathcal An_j=\frac12(\rho_m^2-\rho_i^2-\rho_j^2)$.
\item[(b)] $\tr\mathcal A=\frac23S$, so $\mu=\frac12\bigl(\tau+\sqrt{\tau^2-4}\bigr)$ with $\tau=\frac23S$.
\item[(c)] $S^2-2\sum_i\rho_i^4=3$. Consequently $S\ge3$ and
$S-3=\dfrac{6\mathcal V}{S+3}\le\mathcal V\le4\rho_{\max}^2W$, where $\mathcal V=\sum_i(\rho_i^2-S/3)^2$.
\item[(d)] $\rho_i^2=\frac{\sqrt3}{2}\,\dfrac{\sin\theta_i}{\sin\theta_j\sin\theta_m}
=\frac{\sqrt3}{2}(\cot\theta_j+\cot\theta_m)$, and $S=\sqrt3\sum_i\cot\theta_i$.
\end{enumerate}
\end{lemma}

\begin{proof}
(a) For a $2\times2$ matrix of determinant one, $D^{-\top}=JDJ^{-1}$, where $J$ is the rotation by
$90^\circ$. Since $J^{-1}n_i=\pm e_i$, we get $|D^{-\top}n_i|=|De_i|=\rho_i$. The second formula follows
from $n_0+n_1+n_2=0$.

(b) Sum (a) over $i$ and use $\sum_in_in_i^\top=\frac32I$. The eigenvalues of $\mathcal A$ have product
$1$ and sum $\tau$.

(c) Heron's formula $16|T|^2=(\sum a_i^2)^2-2\sum a_i^4$ and $\ell^4=16/3$ give the identity. Since
$\sum\rho_i^4=\mathcal V+S^2/3$, it reads $S^2=9+6\mathcal V$. Finally
$\mathcal V=\frac13\sum_{i<j}(\rho_i^2-\rho_j^2)^2\le\frac{4\rho_{\max}^2}3\sum_{i<j}(\rho_i-\rho_j)^2
=4\rho_{\max}^2W$.

(d) By the law of sines, $a_i=2R\sin\theta_i$ with $R$ the circumradius, and $1=|T|=2R^2\prod_j\sin\theta_j$.
So $a_i^2=2\sin\theta_i/(\sin\theta_j\sin\theta_m)$. Divide by $\ell^2=4/\sqrt3$ and use
$\sin\theta_i=\sin(\theta_j+\theta_m)$.
\end{proof}

By the symmetry of $\E$, all integrals of products of the $\varphi_i$, and of their gradients, over $\E$
and over its sides reduce to eight numbers:
\begin{align*}
 N_d&=\int_\E\varphi_0^2, & N_o&=\int_\E\varphi_0\varphi_1, \\
 \Gamma_d&=k^2\!\int_\E\sinh^2(L+kd_0), & \Gamma_o&=k^2\!\int_\E\sinh(L+kd_0)\sinh(L+kd_1),\\
 \Sigma_1&=\int_{\text{side }0}\varphi_0^2, & \Sigma_2&=\int_{\text{side }1}\varphi_0^2,\\
 \Sigma_3&=\int_{\text{side }0}\varphi_0\varphi_1, & \Sigma_4&=\int_{\text{side }2}\varphi_0\varphi_1 .
\end{align*}
Thus $\int_\E|\nabla\varphi_i|^2=\Gamma_d$ and $\int_\E\nabla\varphi_i\,\nabla\varphi_j^\top=\Gamma_o\,n_in_j^\top$
for $i\ne j$. We also set
\begin{gather*}
 G_0=\Gamma_d-\Gamma_o,\qquad \Gp=\Gamma_d+\tfrac12\Gamma_o,\qquad
 \Bp=\Sigma_1+2\Sigma_2-2\Sigma_3-\Sigma_4,\qquad \Np=N_d-N_o,\\
 b_{\rm own}=\Sigma_1+2\Sigma_3,\qquad b_{\rm adj}=\Sigma_2+\Sigma_3+\Sigma_4,\qquad
 \chi=b_{\rm adj}-b_{\rm own} .
\end{gather*}
Here $b_{\rm own}=\int_{\text{side }i}\varphi_i\psi_0$ and $b_{\rm adj}=\int_{\text{side }j}\varphi_i\psi_0$ for
$j\ne i$. Hence $B=3(b_{\rm own}+2b_{\rm adj})$. Also $\mathcal G=\sum_{i,j}n_i\cdot n_j\,\Gamma_{ij}=3G_0$,
where $\Gamma_{ii}=\Gamma_d$ and $\Gamma_{ij}=\Gamma_o$; in particular $G_0>0$. The subscript $\perp$ refers to
the complement of $\psi_0$ in the span of the $\varphi_i$:

\begin{lemma}\label{lem:vperp}
Let $v\in\R^3$ with $v_0+v_1+v_2=0$, and let $u_v=\sum_iv_i\varphi_i$. Then
\[
  \int_\E|\nabla u_v|^2=\Gp|v|^2,\qquad \int_{\partial\E}u_v^2=\Bp|v|^2,\qquad
  \int_\E u_v^2=\Np|v|^2 .
\]
In particular $\Gp,\Bp\ge0$ and $\Np>0$. Moreover $\gamma:=\Gp-\beta \Bp-\lambda_0\Np>0$.
\end{lemma}

\begin{proof}
Since $n_i\cdot n_j=-\frac12$ for $i\ne j$ and $\sum_{i\ne j}v_iv_j=-|v|^2$, the first integral is
$\Gamma_d|v|^2-\frac12\Gamma_o\sum_{i\neq j}v_iv_j=\Gp|v|^2$. The other two are computed in the same way;
for the boundary, $\sum_s\int_{\text{side }s}\varphi_i^2=\Sigma_1+2\Sigma_2$ and
$\sum_s\int_{\text{side }s}\varphi_i\varphi_j=2\Sigma_3+\Sigma_4$. The $\varphi_i$ are linearly independent.
Indeed, on side~$0$ they restrict to $\cosh L$, $\cosh(L+y)$ and $\cosh(U-y)$ for a coordinate $y$, where
$U=L+2K$ and $L+U=2M\ne0$. Hence $\Np>0$. Finally, $\gamma|v|^2$ is the value at $u_v$ of the quadratic form
of $-\Delta-\lambda_0$ on $\E$. This form is nonnegative and vanishes only on multiples of the ground state
$\psi_0=u_{\one}$, $\one=(1,1,1)$; and $u_v$ is not such a multiple when $v\ne0$.
\end{proof}

For $c\in\R^3$ let $\tilde u_c=\sum_ic_i\tilde\varphi_i$, and let
\[
  \mathcal P_{ij}=\int_T\nabla\tilde\varphi_i\cdot\nabla\tilde\varphi_j
   -\beta\int_{\partial T}\tilde\varphi_i\tilde\varphi_j-\lambda_0\int_T\tilde\varphi_i\tilde\varphi_j .
\]
Then $\lambda_1(T,-\beta)<\lambda_0$ as soon as $c^\top\mathcal Pc<0$ for some $c$.

\begin{lemma}\label{lem:P}
\begin{enumerate}
\item[(a)] $(\mathcal P\one)_i=G_0(\rho_i^2-1)+\beta\chi(\rho_i-1)-\beta b_{\rm adj}\bigl(\sum_j\rho_j-3\bigr)$.
\item[(b)] $\one^\top\mathcal P\one=G_0(S-3)-\beta\frac B3\bigl(\sum_j\rho_j-3\bigr)\le G_0(S-3)$.
\item[(c)] If $v_0+v_1+v_2=0$, then $v^\top\mathcal Pv\le\Lambda_T|v|^2$, where
$\Lambda_T:=\mu \Gp-\beta\rho_{\min}\Bp-\lambda_0\Np$.
\end{enumerate}
\end{lemma}

\begin{proof}
(a) By~\eqref{eq:transplant} and Lemma~\ref{lem:rho}(a),
\[
 \int_T\nabla\tilde\varphi_i\cdot\nabla\tilde\psi_0
 =\rho_i^2\Gamma_d+\Gamma_o\,n_i^\top\mathcal A(n_j+n_m)=\rho_i^2(\Gamma_d-\Gamma_o)=\rho_i^2G_0,
\]
and $\int_{\partial T}\tilde\varphi_i\tilde\psi_0=\rho_ib_{\rm own}+(\rho_j+\rho_m)b_{\rm adj}$. Hence
$(\mathcal P\one)_i=\rho_i^2G_0-\beta\rho_ib_{\rm own}-\beta(\rho_j+\rho_m)b_{\rm adj}-\lambda_0(N_d+2N_o)$.
For $T=\E$ all $\rho_i=1$, and the left side vanishes because $\psi_0$ is an eigenfunction with
eigenvalue $\lambda_0$. Subtracting the two identities gives (a).

(b) Sum (a) over $i$. The inequality holds because $B>0$ and $\sum_j\rho_j\ge3$, which is the
isoperimetric inequality for triangles.

(c) By~\eqref{eq:transplant},
$\int_T|\nabla\tilde u_v|^2=\int_\E\nabla u_v^\top\mathcal A\nabla u_v\le\mu\int_\E|\nabla u_v|^2$ and
$\int_{\partial T}\tilde u_v^2=\sum_s\rho_s\int_{\text{side }s}u_v^2\ge\rho_{\min}\int_{\partial\E}u_v^2$.
Now use Lemma~\ref{lem:vperp}.
\end{proof}

\begin{proposition}[criterion]\label{prop:crit}
Let $T$ be a non-equilateral triangle of area $1$. Suppose that $\chi\ge0$, and that $\Lambda>0$ satisfies
$v^\top\mathcal Pv\le\Lambda|v|^2$ for all $v\perp\one$. If
\begin{equation}\label{eq:crit}
  \Lambda<\frac{(2G_0\rho_{\min}+\beta\chi)^2}{4G_0\,\rho_{\max}^2},
\end{equation}
then $\lambda_1(T,-\beta)<\lambda_0$.
\end{proposition}

\begin{proof}
It suffices to find $c$ with $c^\top\mathcal Pc<0$. Let $q=\mathcal P\one$ and let $q^\perp=q-\bar q\,\one$
be its component orthogonal to $\one$. For $c=\one-q^\perp/\Lambda$,
\[
 c^\top\mathcal Pc=\one^\top\mathcal P\one-\frac{2|q^\perp|^2}{\Lambda}
 +\frac{(q^{\perp})^\top\mathcal Pq^\perp}{\Lambda^2}
 \le\one^\top\mathcal P\one-\frac{|q^\perp|^2}{\Lambda}.
\]
By Lemma~\ref{lem:P}(a), $q_i-q_j=(\rho_i-\rho_j)\bigl(G_0(\rho_i+\rho_j)+\beta\chi\bigr)$, and the second
factor is at least $2G_0\rho_{\min}+\beta\chi>0$. Hence
\[
 |q^\perp|^2=\tfrac13\sum_{i<j}(q_i-q_j)^2\ge(2G_0\rho_{\min}+\beta\chi)^2\,W .
\]
By Lemma~\ref{lem:P}(b) and Lemma~\ref{lem:rho}(c),
$\one^\top\mathcal P\one\le G_0(S-3)\le4G_0\rho_{\max}^2W$. Therefore
\[
 c^\top\mathcal Pc\le W\Bigl(4G_0\rho_{\max}^2-\frac{(2G_0\rho_{\min}+\beta\chi)^2}{\Lambda}\Bigr)<0,
\]
because $W>0$ for a non-equilateral triangle.
\end{proof}

The mechanism is as follows. For large $\beta$, the Rayleigh quotient of the transplanted ground state
$\tilde\psi_0$ alone exceeds $\lambda_0$ by a positive multiple of $W$. The deformation, however, couples
$\tilde\psi_0$ to the two-dimensional space $\{\tilde u_v:v\perp\one\}$ with strength
$|q^\perp|\approx3G_0\sqrt W$. That space lies only $\gamma\approx3(1-t)G_0$ above $\lambda_0$; this is the
tunnelling splitting. Second-order perturbation theory then gains $|q^\perp|^2/\Lambda$, which exceeds the
loss by a factor of order $(1-t)^{-1}$.

\subsection{The near region}

\begin{lemma}\label{lem:geomC}
Let $\pi/4\le\theta_*\le\pi/3$, and let $T$ be a triangle of area $1$ all of whose angles are at least
$\theta_*$. Then
\[
  \rho_{\rm lo}^2:=\sqrt3\tan\frac{\theta_*}2\le\rho_i^2\le\sqrt3\,\cot\theta_*=:\rho_{\rm hi}^2,
  \qquad
  \mu\le\mu_{\rm hi}:=\frac{\tau_{\rm hi}+\sqrt{\tau_{\rm hi}^2-4}}2,\quad
  \tau_{\rm hi}:=\frac{3\cot^2\theta_*+1}{\sqrt3\,\cot\theta_*}.
\]
Moreover $\rho_{\rm lo}\le1\le\mu_{\rm hi}$.
\end{lemma}

\begin{proof}
By Lemma~\ref{lem:rho}(d), $\rho_i^2=\frac{\sqrt3}2(\cot\theta_j+\cot\theta_m)\le\sqrt3\cot\theta_*$.
Also $\sin\theta_j\sin\theta_m=\frac12(\cos(\theta_j-\theta_m)+\cos\theta_i)\le\cos^2(\theta_i/2)$, so
$\rho_i^2\ge\sqrt3\tan(\theta_i/2)\ge\sqrt3\tan(\theta_*/2)$. For $\mu$, Lemma~\ref{lem:rho}(b),(d) give
$\tau=\frac2{\sqrt3}\sum_i\cot\theta_i$. The angles lie in $[\theta_*,\pi-2\theta_*]\subset(0,\pi/2]$, where
$\cot$ is convex. So $\sum_i\cot\theta_i$ is a convex function on the triangle
$\{\theta_i\ge\theta_*,\ \sum\theta_i=\pi\}$, and its maximum is at a vertex
$(\theta_*,\theta_*,\pi-2\theta_*)$, where it equals $2\cot\theta_*-\cot2\theta_*=(3\cot^2\theta_*+1)/(2\cot\theta_*)$.
Finally, $\mu$ is increasing in $\tau$. Since $\theta_*\le\pi/3$, we have
$\rho_{\rm lo}^2\le\sqrt3\tan(\pi/6)=1$, and $\mu_{\rm hi}\ge1$ because $\tau_{\rm hi}\ge2$.
\end{proof}

For $t\in(0,1)$ let $\rho_{\rm lo}(t)$, $\rho_{\rm hi}(t)$ and $\mu_{\rm hi}(t)$ be the quantities of
Lemma~\ref{lem:geomC} for $\theta_*=\theta(t)$. Since $\sin(\theta(t)/2)=t/2$,
\begin{equation}\label{eq:geomt}
  \tan\frac{\theta(t)}2=\frac{t}{\sqrt{4-t^2}},\qquad \cot\theta(t)=\frac{2-t^2}{t\sqrt{4-t^2}} ,
\end{equation}
and $\theta(t)\ge\pi/4$ is equivalent to $\cot\theta(t)\le1$. For $M>0$ and $t'\in(0,1)$ put
\begin{equation}\label{eq:Lhi}
  \Lambda_{\rm hi}(M;t')=\mu_{\rm hi}(t')\Gp-\beta\rho_{\rm lo}(t')\Bp-\lambda_0\Np,\qquad
  \mathcal R(M;t')=\frac{G_0}{\rho_{\rm hi}(t')^2}\Bigl(\rho_{\rm lo}(t')+\frac{\beta\chi}{2G_0}\Bigr)^2 ,
\end{equation}
where $\beta,\lambda_0,G_0,\Gp,\Bp,\Np,\chi$ are evaluated at $M$. The geometric quantities depend on a
lower bound $t'$ for $t=\tanh M$. This is the form in which the inequality is certified.

\begin{lemma}[certified]\label{lem:certC}
Let $J_a=[M_a,M_a+10^{-3}]$ with $M_a=\frac{13}{10}+\frac a{1000}$, for $a=0,\dots,\clthmCcells-1$, so that
the $J_a$ cover $[13/10,8]$; let $t_a=\tanh M_a$, and let $J_{\rm tail}=[8,\infty)$ with
$t_{\rm tail}=\tanh 8$. For each of these intervals $J$, with its value $t_J$:
\begin{enumerate}
\item[(i)] $\cot\theta(t_J)\le1$;
\item[(ii)] for every $M\in J$: $G_0>0$, $\chi\ge0$ and $\Lambda_{\rm hi}(M;t_J)<\mathcal R(M;t_J)$.
\end{enumerate}
Moreover $\clbetaClo\le\beta(13/10)\le\clbetaChi$.
\end{lemma}

The geometric bounds are monotone in $t'$. As $t'$ increases, $\theta(t')$ increases, so $\rho_{\rm lo}$
increases while $\rho_{\rm hi}$ and $\mu_{\rm hi}$ decrease. Moreover $\Lambda_{\rm hi}$ is increasing in
$\mu_{\rm hi}$ and decreasing in $\rho_{\rm lo}$, because $\Gp,\beta \Bp\ge0$. And $\mathcal R$ is decreasing
in $\rho_{\rm hi}$ and increasing in $\rho_{\rm lo}$, because $\chi\ge0$. Hence (ii) implies the pointwise
inequality $\Lambda_{\rm hi}(M;\tanh M)<\mathcal R(M;\tanh M)$ for all $M\ge13/10$. The proof of
Theorem~\ref{thm:C} below, however, uses (i) and (ii) directly.

\begin{proof}
All quantities have closed forms. Put $U=2M-L=L+2K$. With $y=kd_i\in[0,2K]$, the area element of $\E$ is
$dy_0\,dy_1/(2K^2)$ on $\{y_0,y_1\ge0,\ y_0+y_1\le2K\}$, and the arclength on a side is $dy/(HK)$.
Elementary integration gives
\begin{align*}
 I_1&:=\int_0^{2K}(2K-y)\cosh^2(L+y)\dd y=K^2-\tfrac K2\sinh2L+\tfrac18(\cosh2U-\cosh2L),\\
 I_1^s&:=\int_0^{2K}(2K-y)\sinh^2(L+y)\dd y=-K^2-\tfrac K2\sinh2L+\tfrac18(\cosh2U-\cosh2L),\\
 I_2&:=\iint\cosh(L+y_0)\cosh(L+y_1)=K\sinh2M-\sinh L\,(\sinh U-\sinh L),\\
 I_2^s&:=\iint\sinh(L+y_0)\sinh(L+y_1)=K\sinh2M-\cosh L\,(\cosh U-\cosh L),\\
 J_2&:=\int_0^{2K}\cosh^2(L+y)\dd y=K+\tfrac14(\sinh2U-\sinh2L),\\
 J_3&:=\int_0^{2K}\cosh(L+y)\dd y=\sinh U-\sinh L,\\
 J_4&:=\int_0^{2K}\cosh(L+y)\cosh(U-y)\dd y=K\cosh2M+\tfrac12\sinh2K ,
\end{align*}
and then
\begin{gather*}
 N_d=\frac{I_1}{2K^2},\qquad N_o=\frac{I_2}{2K^2},\qquad \Gamma_d=\frac{2I_1^s}{H^2},\qquad
 \Gamma_o=\frac{2I_2^s}{H^2},\\
 \Sigma_1=\frac{2\cosh^2L}{H},\qquad \Sigma_2=\frac{J_2}{HK},\qquad
 \Sigma_3=\frac{\cosh L\,J_3}{HK},\qquad \Sigma_4=\frac{J_4}{HK}.
\end{gather*}
With $\beta=tK/H$ and $\lambda_0=-4K^2/H^2$, this gives
\begin{gather*}
 G_0=\tfrac2{H^2}(I_1^s-I_2^s),\qquad \Gp=\tfrac2{H^2}\bigl(I_1^s+\tfrac12I_2^s\bigr),\qquad
 -\lambda_0\Np=\tfrac2{H^2}(I_1-I_2),\\
 \beta \Bp=\tfrac2{H^2}\cdot\tfrac t2\bigl(2K\cosh^2L+2J_2-2\cosh L\,J_3-J_4\bigr),\\
 \beta\chi=\tfrac2{H^2}\cdot\tfrac t2\bigl(J_2+J_4-\cosh L\,J_3-2K\cosh^2L\bigr).
\end{gather*}
The inequality $\Lambda_{\rm hi}<\mathcal R$ and the sign conditions are homogeneous of degree one in
the five numbers $G_0$, $\Gp$, $-\lambda_0\Np$, $\beta \Bp$ and $\beta\chi$. We may therefore divide them
by $\frac2{H^2}X$, where $X=e^{2U}/16$. Let $s=e^{-2M}$, $\omega=e^{2L}=(2-t)/(2+t)$ and $u=Ks$. Then
$t=(1-s)/(1+s)$, $e^{2U}=1/(s^2\omega)$ and $e^{2K}=1/(s\omega)$. The hyperbolic functions of $L$ are
rational in $t$; for example $\cosh^2L=4/(4-t^2)$ and $\sinh2L=-4t/(4-t^2)$. After division by $X$, each of
$I_1,\dots,J_4$ and $K\cosh^2L$ becomes a polynomial in $s$, $u$, $\omega$ and these rational functions.
For example,
$I_2^s/X=8\omega u(1-s^2)-4s\bigl(\omega(1+s^2)+1+\omega^2s^2\bigr)+64\omega s^2/(4-t^2)$. These expressions
remain bounded as $s\to0$.

We evaluate them in Arb ball arithmetic. On each $J_a$ we enclose $s,t,K,u$ over the whole interval and
evaluate the geometric bounds~\eqref{eq:geomt} at $t'=t_a$. On the tail we use the single box
$s\in[0,e^{-16}]$, $t\in[\tanh8,1]$, $u\in[0,K(8)e^{-16}]$, with the geometric bounds at $t'=\tanh8$. This
box contains $(s,t,u)$ for every $M\ge8$, because $u=Ke^{-2M}$ is decreasing for $M\ge3/4$ (as
$dK/dM\le\frac32\le2K$). On every interval the script certifies (i), $G_0>0$, $\chi\ge0$ and
$\Lambda_{\rm hi}<\mathcal R$. The ratio $\Lambda_{\rm hi}/\mathcal R$ is at most $\clthmCworst$ on
$[13/10,8]$ and at most $\clthmCtail$ on the tail. The value $\beta(13/10)$ is enclosed directly.
\end{proof}

\begin{proof}[Proof of Theorem~\ref{thm:C}]
Let $\beta\ge\beta_C$, so that $\beta=\beta(M)$ with $M\ge13/10$, and put $t=\tanh M$. Let $T$ be a
non-equilateral triangle of area $1$. If $\sin(\theta_{\min}/2)\le t/2$, Corollary~\ref{cor:far} applies.

Otherwise every angle of $T$ exceeds $\theta(t)$. Let $J$ be an interval of Lemma~\ref{lem:certC}
containing $M$. Then $t_J\le t$, so every angle exceeds $\theta(t_J)$, and $\theta(t_J)\ge\pi/4$ by
Lemma~\ref{lem:certC}(i). Lemma~\ref{lem:geomC} with $\theta_*=\theta(t_J)$ gives
$\rho_{\rm lo}(t_J)\le\rho_{\min}$, $\rho_{\max}\le\rho_{\rm hi}(t_J)$ and $\mu\le\mu_{\rm hi}(t_J)$. Since
$\Gp,\Bp\ge0$, Lemma~\ref{lem:P}(c) gives
$v^\top\mathcal Pv\le\Lambda_T|v|^2\le\Lambda_{\rm hi}(M;t_J)|v|^2$ for $v\perp\one$. Moreover
\[
\Lambda_{\rm hi}(M;t_J)-\gamma=(\mu_{\rm hi}(t_J)-1)\Gp+(1-\rho_{\rm lo}(t_J))\beta \Bp\ge0
\]
by Lemma~\ref{lem:geomC}, so $\Lambda_{\rm hi}(M;t_J)\ge\gamma>0$ by Lemma~\ref{lem:vperp}. By
Lemma~\ref{lem:certC}(ii) we have $\chi\ge0$ and
\[
 \Lambda_{\rm hi}(M;t_J)<\mathcal R(M;t_J)=\frac{(2G_0\rho_{\rm lo}(t_J)+\beta\chi)^2}{4G_0\rho_{\rm hi}(t_J)^2}
 \le\frac{(2G_0\rho_{\min}+\beta\chi)^2}{4G_0\rho_{\max}^2}.
\]
So Proposition~\ref{prop:crit} applies with $\Lambda=\Lambda_{\rm hi}(M;t_J)$.
\end{proof}

\begin{remark}[size of the margin]\label{rem:asymp}
As $M\to\infty$:
\begin{itemize}
\item $\Gp/G_0\to1$, $\beta \Bp/G_0\to2$ and $\beta\chi/G_0\to1$;
\item $\gamma/G_0\sim3(1-t)$, $\mu_{\rm hi}(t)-1\sim\frac83(1-t)$ and $1-\rho_{\rm lo}(t)\sim\frac23(1-t)$;
\item $1-t\sim6e^{-2\cdot3^{1/4}\beta}$.
\end{itemize}
Since $\Lambda_{\rm hi}=\gamma+(\mu_{\rm hi}-1)\Gp+(1-\rho_{\rm lo})\beta \Bp$, it follows that
$\Lambda_{\rm hi}\sim7(1-t)\,G_0\approx42\,e^{-2\cdot3^{1/4}\beta}G_0$, while $\mathcal R$ tends to
$\frac94G_0$. So the margin of the inequality grows exponentially with $\beta$; see Table~\ref{tab:C}. In
floating point, $\Lambda_{\rm hi}(M;\tanh M)<\mathcal R(M;\tanh M)$ holds for $M>\clthrM$ (that is, for
$\beta>\clthrBeta$) and fails below. So the method alone reaches below $6/5$, but not down to small
$\beta$.
\end{remark}

\begin{table}[ht]\centering
\begin{tabular}{rrrrrr}\toprule
$\beta$ & $M$ & $\theta(\tanh M)$ (deg) & $\gamma/G_0$ & $\Lambda_{\rm hi}/G_0$ & $\mathcal R/G_0$\\\midrule
\clTCrowa
\clTCrowb
\clTCrowc
\clTCrowd
\clTCrowe
\clTCrowf
\bottomrule
\end{tabular}
\caption{Lemma~\ref{lem:certC} in floating point, with $t'=\tanh M$; an illustration only (the
certified statement is the lemma). The near region is $\theta_{\min}>\theta(\tanh M)$.}\label{tab:C}
\end{table}

\section{Proof of Theorem~\ref{thm:main}(2) and of Corollary~\ref{cor:perimeter} for
\texorpdfstring{$\alpha<0$}{alpha < 0}}\label{sec:main}

\begin{proof}[Proof of Theorem~\ref{thm:main}\textup{(2)}]
By scaling we may take $|T|=1$, so that $\E_T=\E$ and $\beta=-\alpha$. If $\beta\le6/5$, use
Theorem~\ref{thm:A}. If $\beta\ge6/5$, then $\beta\ge\beta_C$ since $\beta_C\le\clbetaChi<6/5$, and we use
Theorem~\ref{thm:C}. The two ranges overlap on $[\beta_C,6/5]$.
\end{proof}

\begin{proof}[Proof of Corollary~\ref{cor:perimeter} for $\alpha<0$]
By the isoperimetric inequality for triangles, $|T|\le|\E^T|$, with equality only for equilateral $T$.
Theorem~\ref{thm:main}(2) gives $\lambda_1(T,\alpha)\le\lambda_1(\E_T,\alpha)$, with strict inequality unless
$T$ is equilateral. Moreover $\lambda_1(\E_T,\alpha)\le\lambda_1(\E^T,\alpha)$, because for $\alpha<0$ the
lowest eigenvalue of an equilateral triangle is increasing in its area~\cite[Prop.~2.3]{KLV}. This also
follows from Proposition~\ref{prop:equi} and scaling: for the equilateral triangle of area $A$,
$\lambda_1=-4\alpha^2/\tanh^2M$, where $M$ increases with $\beta=-\alpha\sqrt A$.
\end{proof}

\section{Independent cross-checks for Part I}\label{sec:cross}

None of the checks in this section is used in the proof.

\subsection{A second certificate of Lemma~\ref{lem:certC}}
An independent AI referee wrote a second certificate of Lemma~\ref{lem:certC}, sharing no code with the
first. It uses the interval arithmetic of mpmath at $\clindepDps$ digits instead of Arb. On
$[13/10,\clindepMbig]$ it evaluates the un-normalised closed forms directly from $M$ and from the
definitions of the $\Sigma_k$, $\Gamma$ and $N$, on $\clindepCells$ adaptively chosen cells. On the tail
$M\ge\clindepMbig$ it uses a normalisation derived separately with a computer algebra system; the
resulting expressions have integer coefficients. It computes the geometric bounds from half- and
double-angle formulas rather than from~\eqref{eq:geomt}. Its first cell starts at an exact dyadic
number below $13/10$. It confirms (i) and (ii) on every cell, with
worst ratio $\Lambda_{\rm hi}/\mathcal R\le\clindepWorst$ on the body and $\le\clindepTail$ on the tail,
and it confirms $\beta(13/10)<6/5$. At $M=2,5,10,30$ its two forms agree to more than $100$ digits. The same referee found no fatal or serious issue in the arguments of
Sections~\ref{sec:prelim}--\ref{sec:main}.

\subsection{A six-wave certificate for \texorpdfstring{$6/5\le\beta\le3$}{6/5 <= beta <= 3}}\label{sec:six}
Before Theorem~\ref{thm:C} was found, the range $6/5\le\beta\le3$ was proved by a much larger computation,
which now gives an independent check of Theorem~\ref{thm:C} there. Let $w_{2s}=kn_s$ and $w_{2s+1}=-kn_s$
($s=0,1,2$) be the six wave vectors of Proposition~\ref{prop:equi}. Transplanting
$u_c=\sum_jc_je^{w_j\cdot x}$ to $T_{a,b}$ gives the Rayleigh quotient $c^\top A(a,b;\alpha)c/c^\top Nc$,
where
\[
  A_{jk}=(w_j^\top\mathcal A\,w_k)N_{jk}+\alpha\sum_s\rho_sS^s_{jk},\qquad
  N_{jk}=\int_\E e^{(w_j+w_k)\cdot x},\quad S^s_{jk}=\int_{\text{side }s}e^{(w_j+w_k)\cdot x}.
\]
These are divided differences of $\exp$, which we evaluate in ball arithmetic: by difference quotients
when the nodes are separated, and by power series with explicit tail bounds when they are close. The shape
domain $\mathcal D\setminus\{(1,0)\}$ of Lemma~\ref{lem:shape} is covered by four regions; write
$\delta=(a-1,b)$.
\begin{enumerate}
\item[(a)] \emph{Local boxes $|\delta|_\infty\le r$.} Use the trial vector $c_0(M)+(a-1)c_1+bc_2$, where
$c_0$ is the coefficient vector of $\psi_0$ (an exact eigenvector at $\delta=0$). Then
$F=c^\top(A-\lambda_0N)c$ has a critical point with value $0$ at $\delta=0$, and an interval Hessian in
$\delta$, with a mean-value form in $M$ obtained from dual numbers, is negative definite.
\item[(b)] \emph{A near annulus $r<|\delta|_\infty\le R$.} The same trial vector, with $F<0$ checked on
cells.
\item[(c)] \emph{Far cells in $(a,b,\beta)$.} Fixed trial data on each cell, with a coefficient vector
affine in $(a,b)$. The Rayleigh quotient is affine in $\alpha$ and $\lambda_1(\E,\cdot)$ is concave
(Lemma~\ref{lem:concave}), so it suffices to check the two ends of each $\beta$-interval.
\item[(d)] \emph{Thin triangles $a>6$.} The constant trial function and $P/P_*\ge2a/3$.
\end{enumerate}
The certificate consists of $\clcellsLocal$ local $M$-intervals (with $r$ from $\clrLocMin$ to
$\clrLocMax$), $\clcellsNear$ near $M$-intervals (with $R$ from $\clRMin$ to $\clRMax$), and
$\clcellsFar$ far cells in $\clcellsFarFiles$ $\beta$-ranges that tile $[6/5,3]$. The local and near
intervals cover $[M(6/5),M(3)]$. The cell count grows quickly with $\beta$, because the near-equilateral
regime shrinks like $e^{-2\cdot3^{1/4}\beta}$. Removing that obstruction is what Theorem~\ref{thm:C}
achieves.

\subsection{Floating-point checks}
\begin{itemize}
\item \emph{Closed forms.} The closed forms of the eight integrals agree with Gauss quadrature of
$\cosh(L+2Kb_i)$ on actual scalene triangles to $\clquadErr$. The normalised polynomial expressions
agree with the closed forms to rounding error. The referee's independent quadrature of the pencil on $\E$
reproduces the closed forms to $\cldirectErr$.
\item \emph{The chain of inequalities.} Every inequality used in Lemma~\ref{lem:P},
Proposition~\ref{prop:crit} and Lemma~\ref{lem:geomC} holds, and $c^\top\mathcal Pc<0$, on
$\clchainSamples$ random triangles of the near region with $M\in[1.3,5]$. In this test all matrices are
computed by quadrature on the triangle itself. On $\clrunChainTested$ further near-region triangles
(referee's code, $M\in[1.3,6.5]$), the only failures occur at $W<\clrunChainW$, where $c^\top\mathcal Pc$ is
below the rounding level.
\item \emph{The corner bound.} The divergence identity of Lemma~\ref{lem:corner} holds to $\cldivErr$ on
random triangles, and the corner bound holds on $\clcornerSamples$ (triangle, vertex, $\beta$) samples.
\item \emph{Finite elements.} On $\clnearFemSamples$ triangles of the near region with
$\beta\in\{1.16,1.5,2,3\}$, a $P_2$ finite element value of $\lambda_1(T,-\beta)$ (an upper bound up to
rounding) lies below the three-function Ritz value, to within the finite element error on $\E$, and the
Ritz value lies below $\lambda_0$. At $\beta=2$, a $P_1$ finite element computation with Richardson
extrapolation gives $\lambda_1(\E,-2)\approx\cllamFEMtwo$, compared with $\lambda_0=\cllamZeroTwo$.
\item \emph{Consistency of the two parts.} The eight integrals reproduce $g$, $B/N$ and $\lambda_0$ of
Lemma~\ref{lem:HF}. Sampling $\clgeomAsamples$ random triangles with $p\le q_A$ gives
$\min(p-1)/(\tau-2)=\clgeomAmin\ldots>G_A$, in line with Lemma~\ref{lem:geomA}.
\item \emph{Negative tests.} The certificate of Lemma~\ref{lem:certC} fails, as it should, when started at
$M=21/20$, where the inequality is false, or at $M=1$, where $\cot\theta(t)>1$. The certificate of
Lemma~\ref{lem:geomA} fails with $G_A=0.146$, which exceeds the minimum of $(p-1)/(\tau-2)$.
\end{itemize}
\part{The case \texorpdfstring{$\alpha>0$}{alpha > 0}}

\section{Overview of Part II}\label{p-sec:overview}

In this part $\alpha>0$, and the notation is independent of Part~I.

\paragraph{Notation.} $T$ is a triangle with area $A$, perimeter $P$, side lengths
$a_1,a_2,a_3$, inradius $r=2A/P$, incentre $I$, altitudes $h_j=2A/a_j$ and
$R_I=\max_k|v_k-I|$ over the vertices $v_k$. $E$ is the equilateral triangle of area $A$; its
perimeter, inradius and altitude are $P_E=\sqrt{12\sqrt3\,A}$, $r_E=2A/P_E$ and $H=3r_E$.
We write $\lambda_0(\alpha)=\lambda_1(E,\alpha)$ and denote by
$\Lambda_1(T)<\Lambda_2(T)$ the first two Dirichlet eigenvalues of $T$. Norms and inner products
without subscript are those of $L^2$ of the domain in question.

\paragraph{Three explicit results.} The following statements are proved without computation, apart
from an interval check of an elementary inequality in two variables in Theorem~\ref{p-thm:uniform}.
They are used in the proof of Theorem~\ref{thm:main}(1).

\begin{theorem}[thin triangles]\label{p-thm:thin}
If the longest side of $T$ satisfies $a_{\max}\ge P_E$, that is $a_{\max}^2\ge12\sqrt3\,A$, then
$\lambda_1(T,\alpha)>\lambda_0(\alpha)$ for every $\alpha>0$.
\end{theorem}

\begin{theorem}[weak coupling]\label{p-thm:small}
For every triangle $T$ and every $\alpha>0$,
\begin{equation}\label{p-eq:small}
  \lambda_1(T,\alpha)\ \ge\ \alpha\frac PA-\alpha^2\frac{R_I^2}{r^2}.
\end{equation}
Consequently $\lambda_1(T,\alpha)>\lambda_0(\alpha)$ whenever $T$ is not equilateral and
$0<\alpha\le (P-P_E)\,r^2/(A\,R_I^2)$.
\end{theorem}

\begin{theorem}[uniform weak coupling]\label{p-thm:uniform}
If $0<\alpha\sqrt A\le1/100$, then $\lambda_1(T,\alpha)>\lambda_0(\alpha)$ for every
non-equilateral triangle $T$ of area $A$.
\end{theorem}

\paragraph{Structure of the proof of Theorem~\ref{thm:main}(1).} By scaling we may take $A=1$.
The shapes of triangles of area $1$ are parametrised by the hyperbolic plane; we use its Klein
model, a point $w$ of the unit disc, with $w=0$ for $E$ and $|w|=\tanh2t$, where $e^{t}$ is the
stretch factor of the symmetric linear map taking $E$ to $T$ (Section~\ref{p-sec:shapes}).
The parameter set $\{(w,\alpha)\}$ is covered by four regions.
\begin{itemize}
\item[(A)] \emph{Thin triangles} ($a_{\max}\ge P_E$, all $\alpha$): Theorem~\ref{p-thm:thin},
by slicing perpendicular to the longest side.
\item[(S)] \emph{Weak coupling} ($\alpha\le1/100$, all shapes): Theorem~\ref{p-thm:uniform}.
\item[(C)] \emph{Near $E$} ($|w|\le0.1194$, which is slightly less than $\tanh(2\cdot\frac3{50})$;
$\alpha\ge1/100$): Proposition~\ref{p-prop:local}. A
Schur-complement criterion (Lemma~\ref{p-lem:schur}) reduces the inequality to
$D[\psi_0]>\cR/(1-\eta')$, where $D[\psi_0]>0$ is the explicit second-order gain of the frozen
ground state $\psi_0$ of $E$ and $\cR$ is a second-order perturbation sum. The difficulty is an
\emph{upper} bound for $\cR$, uniform in $\alpha\in[1/100,\infty)$; in floating point the gain exceeds
it only by a factor $\clpSpare$ in the worst case (Remark~\ref{p-rem:local}). We obtain it from certified
\emph{lower} bounds for the lowest eigenvalue of two perturbed problems at distance $s=1/10$
from $E$ in opposite directions (Lemma~\ref{p-lem:ring}); by symmetry one number $R_1(\alpha)$
suffices, and the resulting bound is within a few per cent of the true value.
\item[(B)] \emph{Middle region} ($|w|\ge0.1194$, not thin, $\alpha\ge1/100$):
Proposition~\ref{p-prop:middle}. The lowest eigenvalue is a concave function of the coefficients
$(G,m)$ of the quadratic form pulled back to $E$. Positive combinations of the matrices $G$ are
exactly the points of Euclidean triangles in the Klein model, and for such a cell
(Lemma~\ref{p-lem:cell}) $\lambda_1(T_w,\alpha)\ge c\sum_i\theta_i\lambda_1(T_{w_i},\alpha)$ with
the explicit constant $c=(\theta^{\mathsf T}\Gamma\theta)^{-1/2}$,
$\Gamma_{ik}=\cosh2d(w_i,w_k)$. Together with concavity in $\alpha$ this turns certified lower
bounds at the vertices of a prism into a lower bound on the whole prism, with an error of
second order in its size.
\end{itemize}
The certified lower bounds in (B) and (C) are Crouzeix--Raviart bounds
(Lemma~\ref{p-lem:CR}); the discrete eigenvalue is bounded from below by a Collatz--Wielandt
quotient evaluated in exact integer arithmetic (Lemma~\ref{p-lem:CW}), and the interpolation
constant is proved by hand (Lemma~\ref{p-lem:interp}). No verified matrix factorisation is needed,
apart from a small congruence test for the second eigenvalue of $E$ (Lemma~\ref{p-lem:inertia}).
\section{The equilateral triangle}\label{p-sec:equi}

Let $d_1,d_2,d_3$ be the distances to the sides of $E$; then $d_1+d_2+d_3=H$ in $E$.

\begin{proposition}\label{p-prop:equi}
For $\alpha>0$ let $\kappa=\kappa(\alpha)$ be the unique root in $(0,2\pi/H)$ of
\begin{equation}\label{p-eq:kappa}
  \frac{\kappa H}{2}=\arctan\frac{\alpha}{\kappa}+\arctan\frac{2\alpha}{\kappa},
\end{equation}
and let $\delta=\arctan(\alpha/\kappa)$. Then $\lambda_0(\alpha)=\kappa^2$, with eigenfunction
$\psi_0=\sum_{i=1}^3\cos(\kappa d_i-\delta)$. Moreover $\lambda_0(\alpha)<\alpha P_E/A$ and
$\lambda_0(\alpha)<\Lambda_1(E)=4\pi^2/H^2$.
\end{proposition}

\begin{proof}
The left side of \eqref{p-eq:kappa} increases from $0$ to $\pi$ on $(0,2\pi/H]$ and the right side
decreases from $\pi$ to a value below $\pi$, so the root exists and is unique; it depends
continuously on $\alpha$. Since $|\nabla d_i|=1$, $-\Delta\psi_0=\kappa^2\psi_0$. The inward unit
normals satisfy $\nabla d_i\cdot\nabla d_j=-\tfrac12$ for $i\ne j$. On the side $d_1=0$, with
outer normal $\nu=-\nabla d_1$, write $d_2=H/2+\tau$, $d_3=H/2-\tau$. Then
\[
  \partial_\nu\psi_0+\alpha\psi_0
  =\bigl(\alpha\cos\delta-\kappa\sin\delta\bigr)
  +\cos(\kappa\tau)\Bigl(2\alpha\cos\bigl(\tfrac{\kappa H}{2}-\delta\bigr)
     -\kappa\sin\bigl(\tfrac{\kappa H}{2}-\delta\bigr)\Bigr),
\]
and both brackets vanish because $\tan\delta=\alpha/\kappa$ and
$\tan(\kappa H/2-\delta)=2\alpha/\kappa$ by \eqref{p-eq:kappa}. By symmetry the boundary condition
holds on all sides, so $\kappa^2$ is a Robin eigenvalue of $E$ with eigenfunction $\psi_0$;
here $\psi_0\not\equiv0$ because, in the coordinates $(d_1,d_2)$,
$\partial^2\psi_0/\partial d_1\partial d_2=-\kappa^2\cos(\kappa d_3-\delta)\not\equiv0$. As
$\alpha\to0$ we have $\kappa\to0$ (the right side of \eqref{p-eq:kappa} is at most $3\alpha/\kappa$)
and $\alpha/\kappa\to0$ (otherwise the right side of \eqref{p-eq:kappa} would stay bounded away
from $0$ along a sequence while the left side tends to $0$), hence $\delta\to0$ and
$\psi_0\to3$ uniformly; so $\psi_0>0$ and $\kappa(\alpha)^2=\lambda_0(\alpha)$ for small $\alpha$.
The set of $\alpha>0$ with $\kappa(\alpha)^2=\lambda_0(\alpha)$ is closed by continuity. It is
open: the first two eigenvalues $\lambda_0(\alpha)<\lambda_2(\alpha)$ are continuous in $\alpha$
by the min--max principle, so for $\alpha$ near a point $\alpha^*$ of the set the only eigenvalue
below $\frac12(\lambda_0(\alpha^*)+\lambda_2(\alpha^*))$ is $\lambda_0(\alpha)$, and the
eigenvalue $\kappa(\alpha)^2$, being continuous, lies below that number. Hence the set is all of
$(0,\infty)$.
(The formula is a special case of McCartin's \cite{McCartin}.) The bound
$\lambda_0<\alpha P_E/A$ follows from the constant trial function, which is not an
eigenfunction, and $\lambda_0<\Lambda_1(E)$ from $\kappa<2\pi/H$.
\end{proof}

\begin{lemma}\label{p-lem:lam0}
The function $\lambda_0$ is positive, strictly increasing, concave and continuously
differentiable on $(0,\infty)$, with
\begin{equation}\label{p-eq:lam0prime}
  \lambda_0'(\alpha)=\frac{\int_{\partial E}\psi_0^2}{\int_E\psi_0^2}
  =2\kappa\,\frac{\dfrac{\kappa}{\kappa^2+\alpha^2}+\dfrac{2\kappa}{\kappa^2+4\alpha^2}}
  {\dfrac H2+\dfrac{\alpha}{\kappa^2+\alpha^2}+\dfrac{2\alpha}{\kappa^2+4\alpha^2}},
  \qquad \kappa=\kappa(\alpha),
\end{equation}
and $\int_E|\nabla\psi_0|^2=(\lambda_0-\alpha\lambda_0')\int_E\psi_0^2$. Moreover
$\alpha\lambda_0'(\alpha)\le2\bigl(\lambda_0(\alpha)-\lambda_0(\alpha/2)\bigr)$, and for every
$\alpha^*>0$ the tangent $\alpha\mapsto\lambda_0(\alpha^*)+\lambda_0'(\alpha^*)(\alpha-\alpha^*)$
lies above $\lambda_0$ on $(0,\infty)$. If $\beta_1<\dots<\beta_M$ are positive numbers and $U$
is the minimum of the tangents $t_1,\dots,t_M$ at these points, then $U\ge\lambda_0$, and on
each interval $[\beta_m,\beta_{m+1}]$ we have $U=\min(t_m,t_{m+1})$, a piecewise linear function
whose only possible kink is the point of $[\beta_m,\beta_{m+1}]$ where $t_m=t_{m+1}$.
\end{lemma}

\begin{proof}
$\lambda_0$ is an infimum of affine functions of $\alpha$ by \eqref{eq:var}, hence concave. The
implicit function theorem applied to \eqref{p-eq:kappa} (whose left side minus right side has a
positive derivative in $\kappa$) gives $\kappa\in C^1$ and the second expression in
\eqref{p-eq:lam0prime} for $(\kappa^2)'$; it is positive, so $\lambda_0$ is strictly increasing. For the first expression let $\alpha,\beta>0$ and insert
the normalised ground state for $\alpha$ into the quotient for $\beta$, and conversely:
$(\beta-\alpha)\int_{\partial E}\psi_{0,\beta}^2\le\lambda_0(\beta)-\lambda_0(\alpha)\le
(\beta-\alpha)\int_{\partial E}\psi_{0,\alpha}^2$ for $\|\psi_{0,\cdot}\|=1$; since
$\psi_{0,\beta}\to\psi_{0,\alpha}$ uniformly as $\beta\to\alpha$ by the explicit formula, the
derivative is $\int_{\partial E}\psi_{0,\alpha}^2$. The formula for $\int|\nabla\psi_0|^2$ is
$\lambda_0=\int|\nabla\psi_0|^2+\alpha\int_{\partial E}\psi_0^2$. By concavity the slope of the
chord over $[\alpha/2,\alpha]$ is at least $\lambda_0'(\alpha)$, which is the next claim, and the
claim on tangents is concavity again. For the last claim let $j<m$: then
$t_j(\beta_m)\ge\lambda_0(\beta_m)=t_m(\beta_m)$ and the slope of $t_j$ is at least that of
$t_m$, so $t_j\ge t_m$ on $[\beta_m,\infty)$; similarly $t_j\ge t_{m+1}$ on $(0,\beta_{m+1}]$
for $j>m+1$. Finally $t_m\le t_{m+1}$ at $\beta_m$ and $t_m\ge t_{m+1}$ at $\beta_{m+1}$.
\end{proof}


\section{Thin triangles: proof of Theorem~\ref{p-thm:thin}}\label{p-sec:thin}

For $\ell>0$ let $\mu(\ell,\alpha)=k^2$, where $k\in(0,\pi/\ell)$ is the root of
$k\tan(k\ell/2)=\alpha$; this is the lowest eigenvalue of $-w''=\mu w$ on $(0,\ell)$ with Robin
parameter $\alpha$ at both ends, and it is decreasing in $\ell$.

\begin{lemma}\label{p-lem:slice}
For every triangle $T$ and $\alpha>0$, $\lambda_1(T,\alpha)\ge\mu(h_{\min},\alpha)$, where
$h_{\min}=2A/a_{\max}$ is the altitude on the longest side.
\end{lemma}

\begin{proof}
Place the longest side on $\{y=0\}$, so that $T=\{0<x<b,\ 0<y<\ell(x)\}$ with $\ell$ piecewise
linear and $0<\ell\le h_{\min}$ (the angles adjacent to the longest side are acute, so the
foot of the altitude lies on it). On the two other sides $\dd s\ge\dd x$. Hence for
$u\in C^1(\overline T)$, and then for $u\in H^1(T)$ by density,
\begin{align*}
  \int_T|\nabla u|^2+\alpha\int_{\partial T}u^2
  &\ge\int_0^b\Bigl(\int_0^{\ell(x)}|\partial_yu|^2\dd y
      +\alpha\,u(x,0)^2+\alpha\,u(x,\ell(x))^2\Bigr)\dd x\\
  &\ge\int_0^b\mu(\ell(x),\alpha)\int_0^{\ell(x)}u^2\dd y\dd x ,
\end{align*}
and $\mu(\ell(x),\alpha)\ge\mu(h_{\min},\alpha)$.
\end{proof}

\begin{proof}[Proof of Theorem~\ref{p-thm:thin}]
The hypothesis is $h_{\min}\le r_E$, so by Lemma~\ref{p-lem:slice} it suffices to show
$\mu(r_E,\alpha)>\lambda_0(\alpha)$. Let $k$ be the root of $k\tan(kr_E/2)=\alpha$, that is
$kr_E/2=\arctan(\alpha/k)$. The function
$g(\kappa)=\kappa H/2-\arctan(\alpha/\kappa)-\arctan(2\alpha/\kappa)$ is increasing and vanishes
at $\kappa(\alpha)$. Since $H=3r_E$,
\[
  g(k)=3\arctan\frac{\alpha}{k}-\arctan\frac{\alpha}{k}-\arctan\frac{2\alpha}{k}
  =2\arctan\frac{\alpha}{k}-\arctan\frac{2\alpha}{k}>0,
\]
because $\arctan$ is strictly concave on $[0,\infty)$ and vanishes at $0$. Hence
$k>\kappa(\alpha)$, that is $\mu(r_E,\alpha)>\lambda_0(\alpha)$ by
Proposition~\ref{p-prop:equi}.
\end{proof}

\section{Weak coupling: proof of Theorem~\ref{p-thm:small}}\label{p-sec:small}

\begin{lemma}[dual principle]\label{p-lem:dual}
Let $\Omega$ be a bounded Lipschitz domain and $V\in C^1(\overline\Omega;\R^2)$ with
$V\cdot\nu\le\alpha$ on $\partial\Omega$. Then
$\lambda_1(\Omega,\alpha)\ge\inf_\Omega(\operatorname{div}V-|V|^2)$.
\end{lemma}

\begin{proof}
For $u\in H^1(\Omega)$, $0\le|\nabla u+uV|^2$ gives
$|\nabla u|^2\ge-V\cdot\nabla(u^2)-|V|^2u^2$. Integrating by parts,
$\int_\Omega|\nabla u|^2\ge\int_\Omega(\operatorname{div}V-|V|^2)u^2
-\int_{\partial\Omega}(V\cdot\nu)u^2$, and the last term is at least
$-\alpha\int_{\partial\Omega}u^2$.
\end{proof}

\begin{proof}[Proof of Theorem~\ref{p-thm:small}]
Every side of $T$ is tangent to the incircle, so $(x-I)\cdot\nu=r$ on $\partial T$. Take
$V(x)=\alpha(x-I)/r$. Then $V\cdot\nu=\alpha$ on $\partial T$,
$\operatorname{div}V=2\alpha/r=\alpha P/A$ and $|V|^2\le\alpha^2R_I^2/r^2$ on $T$, since $|x-I|$
is maximal at a vertex. Lemma~\ref{p-lem:dual} gives \eqref{p-eq:small}. If
$\alpha R_I^2/r^2\le(P-P_E)/A$ then \eqref{p-eq:small} gives
$\lambda_1(T,\alpha)\ge\alpha P_E/A>\lambda_0(\alpha)$ by Proposition~\ref{p-prop:equi}; here
$P>P_E$ by the isoperimetric inequality for triangles.
\end{proof}

\section{Shapes, concavity and cells}\label{p-sec:shapes}

From now on $A=1$ (this is no restriction, since $\lambda_1(sT,\alpha)=s^{-2}\lambda_1(T,s\alpha)$).
Fix the equilateral triangle $E$ of area $1$ with centroid $0$, unit side directions
$\tau_j=(\cos\vartheta_j,\sin\vartheta_j)$, $\vartheta_j\in\{0,2\pi/3,4\pi/3\}$, and sides
$\Gamma_1,\Gamma_2,\Gamma_3$ of length $s_E=2/H$, $H=3^{1/4}$. For a symmetric positive definite
$G\in\R^{2\times2}$ and $m\in[0,\infty)^3$, $m\neq0$, put
\[
  q_{G,m}[v]=\int_E\nabla v\cdot G\nabla v+\sum_{j=1}^3m_j\int_{\Gamma_j}v^2,\qquad
  \Lambda(G,m)=\inf_{0\ne v\in H^1(E)}\frac{q_{G,m}[v]}{\int_Ev^2}.
\]
The infimum is attained. If $B$ is an invertible matrix and $u=v\circ B^{-1}$ on $B(E)$, then
\begin{equation}\label{p-eq:pullback}
  q_{G,m}[v]=\frac1{|\det B|}\Bigl(\int_{B(E)}\nabla u\cdot BGB^{\mathsf T}\nabla u
  +\sum_j\frac{m_j|\det B|}{|B\tau_j|}\int_{B\Gamma_j}u^2\Bigr),\qquad
  \int_Ev^2=\frac1{|\det B|}\int_{B(E)}u^2 .
\end{equation}
With $B=G^{-1/2}$ this shows that $\Lambda(G,m)$ is the lowest eigenvalue of the Laplacian on
the triangle $G^{-1/2}(E)$ with the Robin coefficient $\beta_j=m_j(\det G)^{-1/2}/|G^{-1/2}\tau_j|$
on its $j$-th side. In particular, let $\cP$ be the set of symmetric positive definite matrices
of determinant one and, for $G\in\cP$,
\[
  T_G=G^{-1/2}(E),\qquad \ell_j(G)=|G^{-1/2}\tau_j|=(\tau_j^{\mathsf T}G^{-1}\tau_j)^{1/2};
  \qquad\text{then}\quad \lambda_1(T_G,\alpha)=\Lambda\bigl(G,\alpha\,\ell(G)\bigr),
\]
$|T_G|=1$, and the sides of $T_G$ have lengths $s_E\ell_j(G)$.

\begin{lemma}[concavity]\label{p-lem:concave}
$\Lambda$ is concave, positively homogeneous of degree one, and nondecreasing with respect to the
Loewner order in $G$ and the componentwise order in $m$. Consequently, if
$\theta_i\ge0$, $\sum_i\theta_i=1$, $c>0$ and
$G\succeq c\sum_i\theta_iG_i$, $m\ge c\sum_i\theta_im_i$, then
$\Lambda(G,m)\ge c\sum_i\theta_i\Lambda(G_i,m_i)$.
\end{lemma}

\begin{proof}
For fixed $v$ the quotient is linear in $(G,m)$ and nondecreasing; an infimum of such
functions is concave, homogeneous and nondecreasing.
\end{proof}

In particular $\alpha\mapsto\lambda_1(T,\alpha)$ is concave and nondecreasing for every triangle
$T$.

\paragraph{Shape coordinates.} For $w=(w_1,w_2)$ in the open unit disc $\mathbb D$ let
\[
  S(w)=\begin{pmatrix}w_1&w_2\\ w_2&-w_1\end{pmatrix},\qquad
  G(w)=\frac{I-S(w)}{\sqrt{1-|w|^2}},\qquad T_w=T_{G(w)},\qquad
  e_j=(\cos2\vartheta_j,\sin2\vartheta_j).
\]

\begin{lemma}\label{p-lem:shape}
\begin{enumerate}
\item[\textup{(a)}] $w\mapsto G(w)$ is a bijection from $\mathbb D$ onto $\cP$, and
$G(w)^{-1}=(I+S(w))/\sqrt{1-|w|^2}$. Hence
$\ell_j(G(w))^2=(1+w\cdot e_j)/\sqrt{1-|w|^2}$, and the angle $\omega_j$ of $T_w$ opposite to its
$j$-th side satisfies $\cot\omega_j=(1-2\,w\cdot e_j)/\sqrt{3(1-|w|^2)}$.
\item[\textup{(b)}] Every triangle of area $1$ is congruent to $T_w$ for some $w$ in the closed
sector $\Sigma=\{w\in\mathbb D:\ 0\le\arg w\le\pi/3\}\cup\{0\}$; $T_w$ is equilateral only for $w=0$.
\item[\textup{(c)}] If $(1+w_1)^2\ge81(1-|w|^2)$, then the longest side of $T_w$ is at least $P_E$.
\item[\textup{(d)}] If $L=\exp(tS_\phi)$ with $t\ge0$ and
$S_\phi=S((\cos2\phi,\sin2\phi))$, then $L(E)=T_w$ with $w=\tanh(2t)(\cos2\phi,\sin2\phi)$;
moreover $G(w)=\cosh(2t)I-\sinh(2t)S_\phi$ and
$\ell_j(G(w))^2=\cosh2t+\sinh2t\,\cos(2\phi-2\vartheta_j)$.
\item[\textup{(e)}] For $w,w'\in\mathbb D$,
$\Gamma(w,w'):=\tfrac12\tr\bigl(G(w)G(w')^{-1}\bigr)=(1-w\cdot w')/\sqrt{(1-|w|^2)(1-|w'|^2)}\ge1$.
\end{enumerate}
\end{lemma}

\begin{proof}
(a) $S(w)^2=|w|^2I$, so $(I-S(w))(I+S(w))=(1-|w|^2)I$ and $\det(I-S(w))=1-|w|^2$; thus
$G(w)\in\cP$ with the stated inverse. Conversely a matrix in $\cP$ is $aI+bS((1,0))+cS((0,1))$
with $a>0$, $a^2-b^2-c^2=1$, which is $G(w)$ exactly for $w=-(b,c)/a$. Next
$\tau_j^{\mathsf T}S(w)\tau_j=w_1\cos2\vartheta_j+w_2\sin2\vartheta_j=w\cdot e_j$. For the angle
use $\cot\omega_j=(\sum_ka_k^2-2a_j^2)/(4|T_w|)$ with $a_k^2=s_E^2\ell_k^2$, $s_E^2=4/\sqrt3$ and
$\sum_ke_k=0$.
(b) Write the triangle as $M(E)+b$; composing with a reflection of $E$ we may assume
$\det M=1$. Then $M=Q(M^{\mathsf T}M)^{1/2}$ with $Q$ a rotation, so the triangle is congruent
to $T_G$ with $G=(M^{\mathsf T}M)^{-1}\in\cP$. If $g$ is one of the six orthogonal symmetries of
$E$, then $T_{g^{\mathsf T}Gg}=g^{\mathsf T}G^{-1/2}g(E)=g^{\mathsf T}(T_G)$ is congruent to
$T_G$, and $g^{\mathsf T}S(w)g=S(w')$ where $w'$ is obtained from $w$ by a rotation through a
multiple of $2\pi/3$, possibly composed with $w_2\mapsto-w_2$ (rotating $E$ by $\vartheta$
rotates $w$ by $-2\vartheta$). These maps form the dihedral group of order six acting on
$\mathbb D$, for which $\Sigma$ is a fundamental domain. If $T_w$ is equilateral then
$\ell_1=\ell_2=\ell_3$, so $w\cdot e_j$ does not depend on $j$, hence $w=0$.
(c) The hypothesis is $\ell_1^2\ge9$, and the sides of $T_w$ are $s_E\ell_j$, $P_E=3s_E$.
(d) $\exp(tS_\phi)=\cosh t\,I+\sinh t\,S_\phi$ is symmetric with determinant one, so
$L(E)=T_G$ with $G=L^{-2}=\exp(-2tS_\phi)=\cosh2t\,I-\sinh2t\,S_\phi=G(w)$ for the stated $w$,
because $\cosh2t=(1-\tanh^22t)^{-1/2}$. The formula for $\ell_j$ follows from (a).
(e) $\tr(S(w)S(w'))=2w\cdot w'$ and $\tr S(w)=0$ give the formula. The matrix
$G(w)G(w')^{-1}$ is similar to a positive definite matrix of determinant one, so its
eigenvalues are $\mu,1/\mu$ with $\mu>0$, and $\frac12(\mu+1/\mu)\ge1$.
\end{proof}

By (a), $\cP$ with the distance $d(G,G')=\frac12\operatorname{arcosh}\Gamma$ is the hyperbolic
plane and $w$ is its Klein (projective) model; we do not use this, but it explains the next
lemma: Euclidean segments in $\mathbb D$ are the hyperbolic geodesics, and the constant $c$ below
is $1-O(\diam^2)$ for a cell of hyperbolic diameter $\diam$.

\begin{lemma}[cells]\label{p-lem:cell}
Let $w_1,\dots,w_N\in\mathbb D$, $G_i=G(w_i)$ and $\Gamma_{ik}=\Gamma(w_i,w_k)$. Let $w$ be a
point of the convex hull of $w_1,\dots,w_N$. Then there are $\theta_i\ge0$ with
$\sum_i\theta_i=1$ such that, with
$c=(\theta^{\mathsf T}\Gamma\theta)^{-1/2}=\bigl(\sum_{i,k}\theta_i\theta_k\Gamma_{ik}\bigr)^{-1/2}\in(0,1]$,
\begin{equation}\label{p-eq:cell}
  G(w)=c\sum_i\theta_iG_i\qquad\text{and}\qquad
  \lambda_1(T_w,\alpha)\ \ge\ c\sum_i\theta_i\,\lambda_1(T_{w_i},\alpha)\quad\text{for every }\alpha>0 .
\end{equation}
\end{lemma}

\begin{proof}
Write $w=\sum_is_iw_i$ with $s_i\ge0$, $\sum s_i=1$, put $a_i=(1-|w_i|^2)^{-1/2}$ and
$\theta_i=(s_i/a_i)/\sum_k(s_k/a_k)$. Then
$\bar G:=\sum_i\theta_iG_i=\sum_i\theta_ia_i(I-S(w_i))=\bigl(\sum_i\theta_ia_i\bigr)(I-S(w))$,
because $S$ is linear and $\theta_ia_i$ is proportional to $s_i$. Thus $\bar G$ is a positive
multiple of $G(w)$, and since $\det G(w)=1$, $G(w)=c\,\bar G$ with $c=(\det\bar G)^{-1/2}$. For
$2\times2$ matrices $\det X=\frac12\tr(X\adj X)$ and $\adj$ is linear, so
$\det\bar G=\frac12\sum_{i,k}\theta_i\theta_k\tr(G_i\adj G_k)=\theta^{\mathsf T}\Gamma\theta$,
as $\adj G_k=G_k^{-1}$. Since $\Gamma_{ik}\ge1$ (Lemma~\ref{p-lem:shape}(e)),
$\theta^{\mathsf T}\Gamma\theta\ge(\sum\theta_i)^2=1$ and $c\le1$.

Next, for $G\in\cP$ we have $G^{-1}=\adj G=J^{\mathsf T}GJ$ with $J$ the rotation through
$\pi/2$; hence $\ell_j(G)^2=(J\tau_j)^{\mathsf T}G(J\tau_j)$ is the restriction to $\cP$ of a
\emph{linear} function of $G$. Therefore
$\sum_i\theta_i\ell_j(G_i)^2=(J\tau_j)^{\mathsf T}\bar G(J\tau_j)=\ell_j(G(w))^2/c$, and by
the Cauchy--Schwarz inequality
$\sum_i\theta_i\ell_j(G_i)\le\bigl(\sum_i\theta_i\ell_j(G_i)^2\bigr)^{1/2}=c^{-1/2}\ell_j(G(w))$.
So $\ell_j(G(w))\ge c^{1/2}\sum_i\theta_i\ell_j(G_i)\ge c\sum_i\theta_i\ell_j(G_i)$. Now
Lemma~\ref{p-lem:concave}, applied to the points $(G_i,\alpha\ell(G_i))$, gives \eqref{p-eq:cell}.
\end{proof}

\begin{corollary}[prisms]\label{p-cor:prism}
Let $w_1,w_2,w_3\in\mathbb D$, $\Gamma_{ik}=\Gamma(w_i,w_k)$, $0<\alpha_0<\alpha_1$, and let
$b_i^0,b_i^1>0$ satisfy $b_i^0\le\lambda_1(T_{w_i},\alpha_0)$,
$b_i^1\le\lambda_1(T_{w_i},\alpha_1)$ $(i=1,2,3)$. Let $U:[\alpha_0,\alpha_1]\to\R$ satisfy
$U\ge\lambda_0$.
\begin{enumerate}
\item[\textup{(a)}] If for all $1\le i\le k\le3$ and all $\alpha=(1-u)\alpha_0+u\alpha_1$,
$u\in[0,1]$,
\begin{equation}\label{p-eq:prism}
  (1-u)\sqrt{b_i^0b_k^0}+u\sqrt{b_i^1b_k^1}\ >\ \sqrt{\Gamma_{ik}}\;U(\alpha),
\end{equation}
then $\lambda_1(T_w,\alpha)>\lambda_0(\alpha)$ for every $w$ in the closed triangle with
vertices $w_1,w_2,w_3$ and every $\alpha\in[\alpha_0,\alpha_1]$.
\item[\textup{(b)}] If $\sqrt{b_i^0b_k^0}>\sqrt{\Gamma_{ik}}\,\Lambda_1(E)$ for all
$1\le i\le k\le3$, then $\lambda_1(T_w,\alpha)>\lambda_0(\alpha)$ for every $w$ in that triangle
and every $\alpha\ge\alpha_0$.
\end{enumerate}
\end{corollary}

\begin{proof}
(a) Fix $\alpha$ and $u$. By concavity in $\alpha$,
$\lambda_1(T_{w_i},\alpha)\ge b_i:=(1-u)b_i^0+ub_i^1$. The geometric mean of two positive
affine functions of $u$ is concave, so $\sqrt{b_ib_k}$ is at least the left side of
\eqref{p-eq:prism}; hence $b_ib_k>\Gamma_{ik}\lambda_0(\alpha)^2$ for all $i,k$ (for $i>k$ by
symmetry). For $w$ in the triangle take $\theta$ and $c$ from Lemma~\ref{p-lem:cell}. Then
$(\sum_i\theta_ib_i)^2=\sum_{i,k}\theta_i\theta_kb_ib_k>\lambda_0(\alpha)^2\,
\theta^{\mathsf T}\Gamma\theta$, that is $c\sum_i\theta_ib_i>\lambda_0(\alpha)$, and
\eqref{p-eq:cell} gives the claim. (b) For $\alpha\ge\alpha_0$,
$\lambda_1(T_{w_i},\alpha)\ge\lambda_1(T_{w_i},\alpha_0)\ge b_i^0$, and the same computation
gives $\lambda_1(T_w,\alpha)>\Lambda_1(E)>\lambda_0(\alpha)$ by Proposition~\ref{p-prop:equi}.
\end{proof}


\section{Certified lower bounds}\label{p-sec:cert}

\begin{lemma}[interpolation constant]\label{p-lem:interp}
Let $K$ be a triangle with longest side $h$ and let $f\in H^1(K)$ have mean value zero on each
of the three sides of $K$. Then
\[
  \|f\|_{L^2(K)}\le\kappa_0\,h\,\|\nabla f\|_{L^2(K)},\qquad
  \kappa_0=\Bigl(\frac1{48}+\frac1{j_{1,1}^2}\Bigr)^{1/2}<0.2983,
\]
where $j_{1,1}=3.8317\ldots$ is the first positive zero of the Bessel function $J_1$. The same
holds with $j_{1,1}$ replaced by $\pi$, that is with the constant
$(1/48+1/\pi^2)^{1/2}<0.3496$.
\end{lemma}

\begin{proof}
(This is the argument of \cite{CGal}; see also \cite{CG}.) Let $\bar f$ be the mean of $f$ over
$K$. Then $\|f\|^2=\|f-\bar f\|^2+|K|\bar f^2$, and $\|f-\bar f\|\le(h/j_{1,1})\|\nabla f\|$
by the optimal Poincar\'e inequality for triangles of Laugesen and Siudeja \cite{LaSiPoincare}
(with $\pi$ in place of $j_{1,1}$ this is the Payne--Weinberger inequality for convex domains
\cite{PW,Bebendorf}). Let $P_i$ be the vertex opposite to the side $e_i$. The field $x-P_i$ is
tangent to the two sides through $P_i$, and $(x-P_i)\cdot\nu=2|K|/|e_i|$ on $e_i$. The
divergence theorem for $f\,(x-P_i)$ gives
$2\int_Kf+\int_K(x-P_i)\cdot\nabla f=(2|K|/|e_i|)\int_{e_i}f=0$. Averaging over $i=1,2,3$,
$2\int_Kf=-\int_K(x-c)\cdot\nabla f$ with $c$ the centroid, so
$|K|\,|\bar f|\le\frac12\|x-c\|_{L^2(K)}\|\nabla f\|$. The polar moment of inertia of a
triangle with sides $a,b,c'$ about its centroid is $\|x-c\|^2_{L^2(K)}=|K|(a^2+b^2+c'^2)/36\le
|K|h^2/12$. Hence $|K|\bar f^2\le(h^2/48)\|\nabla f\|^2$.
\end{proof}

Let $\mathcal T_h$ be a triangulation of a triangle $\Omega$ (conforming: two elements meet in a
common side, a common vertex, or not at all), $V_h$ the Crouzeix--Raviart space \cite{CR} on
$\mathcal T_h$, i.e.\ the piecewise affine functions that are continuous at the midpoints of the
interior edges, with no boundary condition, and $\Pi_h:H^1(\Omega)\to V_h$ the interpolation
operator defined by $\int_e\Pi_hu=\int_eu$ for every edge $e$. Let
$\beta:\partial\Omega\to[0,\infty)$ be constant on each side of $\Omega$ and not identically
zero, and
\[
  a_h(u,v)=\sum_{K\in\mathcal T_h}\int_K\nabla u\cdot\nabla v
  +\sum_{e\subset\partial\Omega}\beta_e\,|e|\,\bar u_e\,\bar v_e ,
\]
where $e$ runs over the boundary edges and $\bar u_e$ is the mean of $u$ on $e$. Let
$\lambda_{1,h}\le\lambda_{2,h}\le\dots$ be the eigenvalues of $a_h$ on $V_h$ with respect to the
$L^2(\Omega)$ inner product, and $\lambda_1\le\lambda_2\le\dots$ the eigenvalues of the form
$q_\beta[u]=\int_\Omega|\nabla u|^2+\int_{\partial\Omega}\beta u^2$ on $H^1(\Omega)$.

\begin{lemma}[Crouzeix--Raviart lower bound]\label{p-lem:CR}
Let $C_h=\kappa_0h$, where $h$ is the largest diameter of an element of $\mathcal T_h$ and
$\kappa_0$ is a constant for which Lemma~\ref{p-lem:interp} holds. Then for $k=1$ and $k=2$,
and every $\sigma$ with $0<\sigma\le\lambda_{k,h}$,
\[
  \lambda_k\ \ge\ \frac{\lambda_{k,h}}{1+C_h^2\lambda_{k,h}}\ \ge\ \frac{\sigma}{1+C_h^2\sigma}.
\]
\end{lemma}

\begin{proof}
(For $\beta=0$ on a part of the boundary and Dirichlet conditions on the rest this is due to
\cite{CG,Liu}; the Robin term requires only the observation that it is diagonal in the edge
means.) First, $a_h$ is positive definite on $V_h$: if $a_h(w,w)=0$ then $w$ is constant on each
element, continuous at the midpoints, hence constant, and a boundary edge with $\beta_e>0$
forces $w=0$. Extend $a_h$ to $H^1(\Omega)+V_h$ by the same formula. For $v_h\in V_h$,
$\partial_\nu v_h$ is constant on each edge of each element, so
$\sum_K\int_K\nabla(u-\Pi_hu)\cdot\nabla v_h=0$; the boundary term of $a_h(u-\Pi_hu,v_h)$
vanishes because the edge means of $u-\Pi_hu$ vanish. Thus
\begin{equation}\label{p-eq:pyth}
  a_h(u,u)=a_h(\Pi_hu,\Pi_hu)+a_h(u-\Pi_hu,u-\Pi_hu)\qquad(u\in H^1(\Omega)).
\end{equation}
By Lemma~\ref{p-lem:interp} on each element,
$\|u-\Pi_hu\|\le C_h\,a_h(u-\Pi_hu,u-\Pi_hu)^{1/2}$. Moreover $|e|\,\bar u_e^2\le\int_eu^2$, so
$a_h(u,u)\le q_\beta[u]$.

Let $U\subset H^1(\Omega)$ be spanned by eigenfunctions for $\lambda_1,\dots,\lambda_k$, so that
$q_\beta[u]\le\lambda_k\|u\|^2$ on $U$. If $\Pi_h$ is not injective on $U$, take $u\in U$ with
$\|u\|=1$ and $\Pi_hu=0$; then $1=\|u-\Pi_hu\|\le C_ha_h(u,u)^{1/2}\le C_h\lambda_k^{1/2}$, so
$\lambda_k\ge C_h^{-2}$, which is larger than the asserted bound. Otherwise $\Pi_hU$ has
dimension $k$ and by the min--max principle in $V_h$ there is $u\in U\setminus\{0\}$ with
$a_h(\Pi_hu,\Pi_hu)\ge\lambda_{k,h}\|\Pi_hu\|^2$. For this $u$, by \eqref{p-eq:pyth} and the
Cauchy--Schwarz inequality,
\begin{align*}
  \|u\|\le\|\Pi_hu\|+\|u-\Pi_hu\|
  &\le\lambda_{k,h}^{-1/2}a_h(\Pi_hu,\Pi_hu)^{1/2}+C_h\,a_h(u-\Pi_hu,u-\Pi_hu)^{1/2}\\
  &\le\bigl(\lambda_{k,h}^{-1}+C_h^2\bigr)^{1/2}a_h(u,u)^{1/2}
  \le\bigl(\lambda_{k,h}^{-1}+C_h^2\bigr)^{1/2}\lambda_k^{1/2}\|u\| .
\end{align*}
The second inequality of the lemma holds because $x\mapsto x/(1+C_h^2x)$ is increasing.
\end{proof}

\paragraph{The matrices.} In the Crouzeix--Raviart basis (one function per edge, equal to $1$
at its midpoint; on an element it is $1-2\mu$, with $\mu$ the barycentric coordinate of the
opposite vertex) the mass matrix $M_h$ is diagonal, with entries $\frac13$ of the total area of
the (one or two) elements containing the edge; the boundary term of $a_h$ is diagonal with
entries $\beta_e|e|$; and the element stiffness matrix has off-diagonal entries
$-2\cot\omega$, where $\omega$ is the angle of the element opposite to its third edge, and zero
row sums. Hence, on a triangulation without obtuse angles, the matrix $K_h$ of $a_h$ has
nonpositive off-diagonal entries. We use two triangulations of a triangle $\Omega$, depending on
an integer $n$:
\begin{itemize}
\item[(U)] if $\Omega$ is acute: the subdivision into $n^2$ congruent triangles similar to
$\Omega$;
\item[(A)] in general: $\Omega$ is cut along the altitude on its longest side into two right
triangles, and each of them is subdivided into $n^2$ congruent right triangles similar to it
(the two subdivisions induce the same $n$ equal segments on the altitude, so the union is
conforming).
\end{itemize}
In both cases all data of $K_h$ and $M_h$ are determined by $n$ and by the squared side
lengths, the area and the Robin coefficients of $\Omega$: the cotangent of the angle of
$\Omega$ opposite to the side $a_j$ is $(\sum_ka_k^2-2a_j^2)/(4|\Omega|)$; in (A) the two right
triangles have the acute angles $\omega_B,\pi/2-\omega_B$ and $\omega_C,\pi/2-\omega_C$, where
$\omega_B,\omega_C$ are the angles of $\Omega$ adjacent to its longest side, their legs on that
side have lengths $h_\Omega\cot\omega_B$ and $h_\Omega\cot\omega_C$ with $h_\Omega$ the altitude,
and the coupling across a right angle is exactly zero. The largest element diameter is the
longest side of $\Omega$ divided by $n$ in (U), and the longer of the two other sides divided by
$n$ in (A).

\begin{lemma}[Collatz--Wielandt certificate]\label{p-lem:CW}
Let $K$ be a real symmetric matrix with nonpositive off-diagonal entries and $M$ a diagonal
matrix with positive diagonal. If $x$ has positive entries and $Kx\ge\sigma Mx$ entrywise, then
every eigenvalue of the pencil $(K,M)$ is at least $\sigma$.
\end{lemma}

\begin{proof}
Let $B=K-\sigma M$ and $X=\operatorname{diag}(x)$. The symmetric matrix $XBX$ has nonpositive
off-diagonal entries and row sums $x_i(Bx)_i\ge0$, so it is diagonally dominant with nonnegative
diagonal, hence positive semidefinite by Gershgorin's theorem. So $B$ is positive semidefinite.
\end{proof}

\begin{lemma}[second eigenvalue]\label{p-lem:inertia}
Let $K$ be real symmetric, $M$ symmetric positive definite, $x$ a vector, $g\ge0$ and
$\sigma\in\R$. If $K-\sigma M+g\,(Mx)(Mx)^{\mathsf T}$ is positive definite, then the second
eigenvalue of the pencil $(K,M)$ is at least $\sigma$. A symmetric matrix $B$ is positive
definite if there is a matrix $S$ such that $S^{\mathsf T}BS$ is strictly diagonally dominant
with positive diagonal.
\end{lemma}

\begin{proof}
On the hyperplane $\{y:(Mx)^{\mathsf T}y=0\}$ we have
$y^{\mathsf T}Ky-\sigma y^{\mathsf T}My=y^{\mathsf T}By>0$ for $y\ne0$, and the second eigenvalue
is the maximum over hyperplanes of the minimum of the Rayleigh quotient. If
$C=S^{\mathsf T}BS$ is strictly diagonally dominant with positive diagonal, it is positive
definite by Gershgorin's theorem; then $S$ is injective, hence invertible, and
$B=S^{-\mathsf T}CS^{-1}$ is positive definite.
\end{proof}

\paragraph{How the bounds are certified.} For given $(a_1^2,a_2^2,a_3^2,|\Omega|,\beta,n)$, a
floating-point approximation $x>0$ of the discrete ground state is computed (sparse
shift-and-invert Lanczos, followed by damped Jacobi sweeps, which make the small components
near the corners accurate); it is not trusted. The vector $x$ is converted \emph{exactly} into a
vector of positive integers (a common power of two times $x$). The real coefficients
($2\cot\omega$ for the angle classes, $\beta_e|e|$, and $\frac13$ of the element areas; at
most a dozen numbers) are enclosed in Arb balls \cite{Arb} and replaced by rational lower or
upper bounds with denominator $2^{60}$, the direction of rounding being chosen for each term
according to the sign of its exact integer cofactor. The inequality $Kx\ge\sigma Mx$ of
Lemma~\ref{p-lem:CW} then becomes a sign test on exact integers, and the nonnegativity of the
cotangents is checked on their lower bounds. This gives $\lambda_{1,h}\ge\sigma$ for a dyadic
rational $\sigma$, and Lemma~\ref{p-lem:CR} gives a lower bound for $\lambda_1(\Omega,\beta)$,
evaluated in ball arithmetic with an upper bound for $h$. The hypothesis $\sigma>0$ of
Lemma~\ref{p-lem:CR} is checked explicitly wherever the lemma is applied (for $\sigma$ below
$-1/C_h^2$ the right-hand side would be positive and large), and a re-check tests the stored
rational $\sigma$ itself, without rounding. For the second eigenvalue of $E$ the
matrices of Lemma~\ref{p-lem:inertia} are formed in ball arithmetic, with $x$ and $S$
floating-point matrices that are entered exactly ($S$ approximates an inverse Cholesky factor).


\section{The Schur-complement criterion and uniform weak coupling}\label{p-sec:uniform}

The following criterion is the tool for all shapes close to $E$: for $\alpha\le1/100$ in this
section and for $\alpha\ge1/100$ in Section~\ref{p-sec:local}.

\begin{lemma}[Schur-complement criterion near $E$]\label{p-lem:schur}
Fix $\alpha>0$, let $q_0=q_{I,\alpha(1,1,1)}$, let $\psi_0$ be its $L^2$-normalised ground state
and $\lambda_0<\lambda_2$ its first two eigenvalues (so $\lambda_0=\lambda_1(E,\alpha)$ and
$\lambda_2$ is the second eigenvalue). Let $(G,m)$ satisfy
$q_{G,m}[w]\ge(1-\eta)q_0[w]$ for all $w\in H^1(E)$, with $0\le\eta<1-\lambda_0/\lambda_2$, and put
$D=q_{G,m}-q_0$ (a symmetric bilinear form), $\eta'=\eta\lambda_2/(\lambda_2-\lambda_0)<1$ and
\[
  \mathcal R=\sup\Bigl\{\frac{D(\psi_0,w)^2}{q_0[w]-\lambda_0\|w\|^2}:\ 0\ne w\in H^1(E),\
  \langle w,\psi_0\rangle=0\Bigr\}.
\]
If $D[\psi_0]\ge\mathcal R/(1-\eta')$, then $\Lambda(G,m)\ge\lambda_0$; if the inequality is
strict, then $\Lambda(G,m)>\lambda_0$.
\end{lemma}

\begin{proof}
Let $v=c\psi_0+w$ with $w\perp\psi_0$ in $L^2(E)$; then $q_0(\psi_0,w)=0$ and
$q_0[w]\ge\lambda_2\|w\|^2$. Hence
\[
  q_{G,m}[v]-\lambda_0\|v\|^2=c^2D[\psi_0]+2c\,D(\psi_0,w)+q_{G,m}[w]-\lambda_0\|w\|^2 ,
\]
and $q_{G,m}[w]-\lambda_0\|w\|^2\ge(1-\eta)q_0[w]-\lambda_0\|w\|^2\ge
(1-\eta')(q_0[w]-\lambda_0\|w\|^2)$, the last step being equivalent to
$(\eta'-\eta)q_0[w]\ge\eta'\lambda_0\|w\|^2$, which holds since
$(\eta'-\eta)\lambda_2=\eta'\lambda_0$. Writing $N=q_0[w]-\lambda_0\|w\|^2$ and using
$|D(\psi_0,w)|\le(\mathcal RN)^{1/2}$,
\[
  q_{G,m}[v]-\lambda_0\|v\|^2\ge c^2D[\psi_0]-2|c|(\mathcal RN)^{1/2}+(1-\eta')N\ge
  c^2\Bigl(D[\psi_0]-\frac{\mathcal R}{1-\eta'}\Bigr)\ge0 .
\]
If $D[\psi_0]>\mathcal R/(1-\eta')$, the right side is positive for $c\ne0$, and for $c=0$,
$w\ne0$ the left side is at least $(1-\eta')N>0$; since the infimum defining $\Lambda(G,m)$ is
attained, $\Lambda(G,m)>\lambda_0$.
\end{proof}
\begin{remark}\label{p-rem:lam2}
Lemma~\ref{p-lem:schur} and its proof remain valid if $\lambda_2$ is replaced throughout (in the
condition on $\eta$ and in the definition of $\eta'$) by a number $\underline\lambda_2$ with
$\lambda_0<\underline\lambda_2\le\lambda_2$, because only the inequality
$q_0[w]\ge\underline\lambda_2\|w\|^2$ for $w\perp\psi_0$ is used.
\end{remark}

\subsection*{Proof of Theorem~\ref{p-thm:uniform}}
By the scaling $\lambda_1(sT,\alpha)=s^{-2}\lambda_1(T,s\alpha)$ we may take $A=1$; then
$H=3^{1/4}$, $P_E=6/H$, $r_E=H/3$ and the side of $E$ is $2/H$. Put $\alpha_1=1/100$. Every
triangle of area $1$ is congruent to $L(E)$ with $L=\exp(tS_\phi)$, $t\ge0$,
$S_\phi=\bigl(\begin{smallmatrix}\cos2\phi&\sin2\phi\\ \sin2\phi&-\cos2\phi\end{smallmatrix}\bigr)$
(write the triangle as $M(E)+b$ with $\det M=1$, after composing with a reflection of $E$ if
necessary, and take the polar decomposition of $M$). If $\tau_j=(\cos\vartheta_j,\sin\vartheta_j)$
are the side directions of $E$, $\vartheta_j\in\{0,2\pi/3,4\pi/3\}$, then
\begin{equation}\label{p-eq:ell}
  G=\exp(-2tS_\phi),\qquad \ell_j^2=|L\tau_j|^2=\cosh2t+\sinh2t\,\cos(2\phi-2\vartheta_j),
\end{equation}
the sides of $T$ are $a_j=2\ell_j/H$, and $t=0$ exactly when $T$ is equilateral. We write
$\bar\ell=\frac13\sum_j\ell_j=P/P_E$, $\hat p=\bar\ell-1\ge0$,
$|z_0|=\bigl(\tfrac23\sum_j(\ell_j-\bar\ell)^2\bigr)^{1/2}$ and $\zeta=2\hat p+|z_0|$.

\begin{lemma}[the equilateral triangle at weak coupling]\label{p-lem:Eweak}
Let $A=1$, $0<\alpha\le\alpha_1$, and let $\psi_0$ be the $L^2(E)$-normalised ground state. Then
\begin{enumerate}
\item[\textup{(a)}] $\alpha P_E-\alpha^2/(1-4\alpha/P_E)\le\lambda_0(\alpha)\le
\alpha P_E-\alpha^2+\tfrac{11}{5}\alpha^3/P_E$;
\item[\textup{(b)}] $\int_{\partial E}\psi_0^2\ge P_E-g_2\alpha$ and
$\|\nabla\psi_0\|_{L^2(E)}\le g_1\alpha$, where, with $c=1.18$,
\[
  g_2=\frac1{c-1}\Bigl(\frac{c^2}{1-4c\alpha_1/P_E}-1+\frac{11\alpha_1}{5P_E}\Bigr)<\clpGtwo,
  \qquad g_1=\Bigl(g_2-1+\frac{11\alpha_1}{5P_E}\Bigr)^{1/2}<\clpGone;
\]
\item[\textup{(c)}] the second Robin eigenvalue satisfies
$\lambda_2(E,\alpha)\ge\mu_2:=4\pi^2/(3\sqrt3)$.
\end{enumerate}
\end{lemma}

\begin{proof}
(a) With $y=\alpha/\kappa$, \eqref{p-eq:kappa} reads
$\lambda_0=\kappa^2=\frac{2\alpha}{H}\,(\arctan y+\arctan2y)/y$. For $y\ge0$,
$y-y^3/3\le\arctan y\le y-y^3/3+y^5/5$ (the differences are $\int_0^ys^4/(1+s^2)\dd s$ and
$\int_0^ys^6/(1+s^2)\dd s$). Since $6/H=P_E$,
\[
  \alpha P_E(1-y^2)\le\lambda_0\le\alpha P_E\bigl(1-y^2+\tfrac{11}{5}y^4\bigr),\qquad
  y^2=\alpha^2/\lambda_0 .
\]
Theorem~\ref{p-thm:small} for $E$ (where $R_I=2r$) gives $\lambda_0\ge\alpha P_E-4\alpha^2$, and
$\lambda_0\le\alpha P_E$; hence $\alpha/P_E\le y^2\le\alpha/(P_E-4\alpha)$. The lower bound
follows. For the upper bound note that $u\mapsto1-u+\frac{11}{5}u^2$ is decreasing on
$[0,5/22]$, which contains the range of $y^2$, and evaluate at $u=\alpha/P_E$.

(b) $\lambda_0$ is concave and, being a simple eigenvalue of an analytic family of forms,
differentiable with $\lambda_0'(\alpha)=\int_{\partial E}\psi_0^2$. By concavity and (a) (the
lower bound in (a) holds for every $\alpha$ with $4\alpha<P_E$),
\[
  \lambda_0'(\alpha)\ge\frac{\lambda_0(c\alpha)-\lambda_0(\alpha)}{(c-1)\alpha}
  \ge P_E-\frac{\alpha}{c-1}\Bigl(\frac{c^2}{1-4c\alpha/P_E}-1+\frac{11\alpha}{5P_E}\Bigr)
  \ge P_E-g_2\alpha .
\]
Then $\|\nabla\psi_0\|^2=\lambda_0-\alpha\lambda_0'(\alpha)\le
\alpha^2\bigl(g_2-1+\tfrac{11}{5}\alpha/P_E\bigr)\le g_1^2\alpha^2$.

(c) By the min--max principle $\lambda_2(E,\alpha)\ge\lambda_2(E,0)$, the first nonzero Neumann
eigenvalue of $E$, which equals $16\pi^2/(9s^2)$ for side $s$ \cite{McCartin2}; here $s^2=4/\sqrt3$.
\end{proof}

\begin{lemma}\label{p-lem:Z}
Let $\nu_j$ be the outer normal on $\Gamma_j$, $\beta=\hat p/r_E$,
$z_0=\frac23\sum_k(\ell_k-\bar\ell)\nu_k$ and $Z(x)=\beta x+z_0$. Then $Z\cdot\nu_j=\ell_j-1$ on
$\Gamma_j$, $\sup_E|Z|\le\zeta$, and for $f\in W^{1,1}(E)$,
$\sum_j(\ell_j-1)\int_{\Gamma_j}f=\int_E(2\beta f+Z\cdot\nabla f)$.
\end{lemma}

\begin{proof}
On $\Gamma_j$, $x\cdot\nu_j=r_E$ and, since $\nu_k\cdot\nu_j=-\frac12$ for $k\ne j$ and
$\sum_k(\ell_k-\bar\ell)=0$, $z_0\cdot\nu_j=\ell_j-\bar\ell$; so
$Z\cdot\nu_j=\hat p+\ell_j-\bar\ell=\ell_j-1$. The identity is the divergence theorem for $fZ$.
The convex function $|Z|$ attains its maximum over $E$ at a vertex, where $|x|=2r_E$; and
$|z_0|^2=\frac49\cdot\frac32\sum_k(\ell_k-\bar\ell)^2$ because
$|\sum_kx_k\nu_k|^2=\frac32\sum_kx_k^2$ when $\sum_kx_k=0$.
\end{proof}

\begin{proposition}\label{p-prop:L}
Let $A=1$, $0<\alpha\le\alpha_1$, $T=\exp(tS_\phi)(E)$ with $t>0$, and
$\varepsilon_1=\alpha_1P_E/\mu_2$, $\varepsilon_2=\alpha_1P_E/(\mu_2-\alpha_1P_E)$. If
\begin{equation}\label{p-eq:L}\tag{L}
  \hat p\,(P_E-g_2\alpha_1)\bigl(e^{-2t}(1+\varepsilon_2)-\varepsilon_2\bigr)
  >\alpha_1\Bigl(\frac{g_1(e^{2t}-1)+\zeta}{(1-\varepsilon_1)^{1/2}}
  +\frac{g_1\alpha_1\zeta}{(\mu_2-\alpha_1P_E)^{1/2}}\Bigr)^2,
\end{equation}
then $\lambda_1(T,\alpha)>\lambda_0(\alpha)$.
\end{proposition}

\begin{proof}
We verify the hypotheses of Lemma~\ref{p-lem:schur} for $(G,m)=(G,\alpha\ell)$. The eigenvalues of
$G$ are $e^{\pm2t}$ and $\ell_j\ge e^{-t}$ by \eqref{p-eq:ell}, so
$q_{G,\alpha\ell}\ge e^{-2t}q_0$ and we take $\eta=1-e^{-2t}$. By Proposition~\ref{p-prop:equi} and
Lemma~\ref{p-lem:Eweak}(c), $\lambda_0/\lambda_2\le\varepsilon_1$ and
$\lambda_0/(\lambda_2-\lambda_0)\le\varepsilon_2$, hence
\[
  1-\eta'=e^{-2t}-(1-e^{-2t})\frac{\lambda_0}{\lambda_2-\lambda_0}
  \ge e^{-2t}(1+\varepsilon_2)-\varepsilon_2 ,
\]
which is positive by \eqref{p-eq:L}; in particular $\eta<1-\lambda_0/\lambda_2$. Next, by the
$D_3$-invariance of $\psi_0$ (so that $\int_{\Gamma_j}\psi_0^2$ does not depend on $j$ and
$\int_E\partial_a\psi_0\partial_b\psi_0=\frac12\delta_{ab}\|\nabla\psi_0\|^2$), $\tr G=2\cosh2t$
and Lemma~\ref{p-lem:Eweak}(b),
\[
  D[\psi_0]=(\cosh2t-1)\|\nabla\psi_0\|^2+\alpha\hat p\int_{\partial E}\psi_0^2
  \ge\alpha\hat p\,(P_E-g_2\alpha).
\]
For $w\perp\psi_0$, Lemma~\ref{p-lem:Z} with $f=\psi_0w$ and $\int_E\psi_0w=0$ gives
\[
  D(\psi_0,w)=\int_E\bigl((G-I)\nabla\psi_0+\alpha\psi_0Z\bigr)\cdot\nabla w
  +\alpha\int_E(Z\cdot\nabla\psi_0)\,w ,
\]
so $|D(\psi_0,w)|\le X\|\nabla w\|+\alpha Y\|w\|$ with
$X\le(e^{2t}-1)\|\nabla\psi_0\|+\alpha\zeta\le\alpha\bigl(g_1(e^{2t}-1)+\zeta\bigr)$ and
$Y\le\zeta\|\nabla\psi_0\|\le g_1\alpha\zeta$. With $N=q_0[w]-\lambda_0\|w\|^2$ we have
$N\ge(1-\lambda_0/\lambda_2)\|\nabla w\|^2$ and $N\ge(\lambda_2-\lambda_0)\|w\|^2$, therefore
\[
  \mathcal R\le\alpha^2\Bigl(\frac{g_1(e^{2t}-1)+\zeta}{(1-\varepsilon_1)^{1/2}}
  +\frac{g_1\alpha\zeta}{(\mu_2-\alpha_1P_E)^{1/2}}\Bigr)^2\le\alpha\,\alpha_1B^2 ,
\]
where $B$ denotes the bracket on the right of \eqref{p-eq:L} (we used $\alpha\le\alpha_1$ twice).
Since also $D[\psi_0]\ge\alpha\hat p\,(P_E-g_2\alpha_1)$, inequality \eqref{p-eq:L} multiplied by
$\alpha$ gives $D[\psi_0](1-\eta')>\mathcal R$. Now apply Lemma~\ref{p-lem:schur}.
\end{proof}

\begin{proof}[Proof of Theorem~\ref{p-thm:uniform}]
Let $A=1$, $0<\alpha\le\alpha_1$ and $T=\exp(tS_\phi)(E)$, $t>0$. It suffices that one of the
following holds:
\begin{itemize}
\item[(W)] $\max_j\ell_j\ge3$; then $a_{\max}\ge P_E$ and Theorem~\ref{p-thm:thin} applies;
\item[(S)] $\alpha_1R_I^2\le(P-P_E)r^2$, where $r=2/P$ and
$R_I^2=\max_k\bigl(r^2+(P/2-a_k)^2\bigr)$; then Theorem~\ref{p-thm:small} applies;
\item[(L)] inequality \eqref{p-eq:L}; then Proposition~\ref{p-prop:L} applies.
\end{itemize}
All three conditions are explicit in $(t,\phi)$ through \eqref{p-eq:ell}. Since the largest of the
three numbers $\cos(2\phi-2\vartheta_j)$ is at least $\frac12$ and
$\cosh(2t)+\frac12\sinh(2t)\ge9$ for $t\ge1.2413$, (W) holds for $t\ge1.2413$. For $0<t\le t_0=1/100$
we use the elementary bounds $\cosh2t-1\le2C_ct^2$, $\sinh2t\le2C_st$, $e^{2t}-1\le2C_et$,
$e^t-1\le C_1t$, where $C_c,C_s,C_e,C_1$ are the values of the corresponding quotients at $t_0$;
with $u_j=\ell_j^2-1$, Taylor's formula $\sqrt{1+u}\ge1+u/2-u^2/8-K|u|^3$,
$K=\frac1{16}(1-u_0)^{-5/2}$, $u_0=e^{2t_0}-1$, and $\frac13\sum u_j=\cosh2t-1\ge2t^2$,
$\frac13\sum u_j^2=(\cosh2t-1)^2+\frac12\sinh^22t$, $|u_j|\le e^{2t}-1$,
$\hat p\le\frac16\sum u_j=\frac12(\cosh2t-1)$ (from $\sqrt{1+u}\le1+u/2$) and
$|z_0|^2\le\frac23\sum(\ell_j-1)^2\le2(e^t-1)^2$ (the mean minimises the sum of squares, and
$e^{-t}\le\ell_j\le e^t$), they give
\[
  \hat p\ge c_pt^2,\qquad \zeta\le c_\zeta t,\qquad
  c_p=1-\tfrac18\bigl(4t_0^2C_c^2+2C_s^2\bigr)-8Kt_0C_e^3,\quad
  c_\zeta=2C_ct_0+\sqrt2\,C_1 ,
\]
so that \eqref{p-eq:L} follows, after division by $t^2$, from
$c_p(P_E-g_2\alpha_1)(e^{-2t_0}(1+\varepsilon_2)-\varepsilon_2)>
\alpha_1\bigl((2g_1C_e+c_\zeta)(1-\varepsilon_1)^{-1/2}
+g_1\alpha_1c_\zeta(\mu_2-\alpha_1P_E)^{-1/2}\bigr)^2$. In ball arithmetic the left side exceeds $\clpWeakLhs$ and the right side is below $\clpWeakRhs$.
On the remaining rectangle $[t_0,1.2413]\times[0,\pi]$ the script \texttt{weak\_cert.py} verifies
in ball arithmetic (Arb, through python-flint, $80$ bits), by adaptive bisection starting from a
$64\times96$ grid, that each box of the final subdivision satisfies (W), (S) or (L) in its
entirety. The run uses $\clpWeakBoxes$ boxes, of which $\clpWeakLeaves$ are leaves ($\clpWeakW$ of type
(W), $\clpWeakS$ of type (S), $\clpWeakL$ of type (L)), with bisection depth at most $\clpWeakDepth$ and
no failure. The same script evaluates the inequality $\cosh(2t)+\frac12\sinh(2t)\ge9$ at $t=1.2413$ and
the two displayed numbers.
\end{proof}

\begin{remark}
The value $1/100$ is limited by condition (S) for tall isosceles triangles close to the thin regime
($t\approx1.24$), where the threshold of Theorem~\ref{p-thm:small} is about $\clpSmallThr$; near $E$
inequality \eqref{p-eq:L} holds with a factor $\clpWeakFactor$ to spare (the quotient of the two
displayed numbers).
\end{remark}

\section{A neighbourhood of the equilateral triangle}\label{p-sec:local}

Throughout this section $A=1$, $\alpha>0$ is fixed, $q_0=q_{I,\alpha(1,1,1)}$, $\psi_0$ is its
$L^2$-normalised ground state, $\lambda_0=\lambda_0(\alpha)$, $\lambda_2=\lambda_2(E,\alpha)$
is the second eigenvalue, $\underline\lambda_2$ is a number with
$\lambda_0<\underline\lambda_2\le\lambda_2$, and
\[
  N[w]=q_0[w]-\lambda_0\|w\|^2,\qquad
  \gamma=\frac{\underline\lambda_2}{\underline\lambda_2-\lambda_0}.
\]
For $w\perp\psi_0$ we have $q_0[w]\ge\lambda_2\|w\|^2$, hence
\begin{equation}\label{p-eq:Ngap}
  N[w]\ \ge\ \gamma^{-1}q_0[w]\qquad\text{and}\qquad N[w]\ \ge\ (\underline\lambda_2-\lambda_0)\|w\|^2
  \qquad(w\perp\psi_0).
\end{equation}
A \emph{direction} is a pair $\theta=(\Theta,\mu)$ of a symmetric $2\times2$ matrix and a
vector $\mu\in\R^3$. For a direction let
\[
  D_\theta(u,v)=\int_E\nabla u\cdot\Theta\nabla v+\alpha\sum_j\mu_j\int_{\Gamma_j}uv,\qquad
  R(\theta)=\sup\Bigl\{\frac{D_\theta(\psi_0,w)^2}{N[w]}:\ 0\ne w\in H^1(E),\ w\perp\psi_0\Bigr\}.
\]
Thus $q_{G,m}-q_0=D_\theta$ with $\theta=(G-I,(m_j/\alpha-1)_j)$, and the quantity $\cR$ of
Lemma~\ref{p-lem:schur} is $R(\theta)$. By \eqref{p-eq:Ngap}, $N^{1/2}$ is a Hilbert norm on
$\{w\perp\psi_0\}$ equivalent to the $H^1$ norm, and $R(\theta)^{1/2}$ is the norm of the
functional $D_\theta(\psi_0,\cdot)$ in the dual space. Since $D_\theta$ is linear in $\theta$,
\begin{equation}\label{p-eq:seminorm}
  \theta\mapsto R(\theta)\ \text{is a nonnegative quadratic form, and}\ R^{1/2}\ \text{is a seminorm.}
\end{equation}

\subsection{Symmetry}

Let $\mathcal G$ be the group of the six orthogonal maps $g$ with $g(E)=E$. For $v\in H^1(E)$
put $v^g=v\circ g$, and for a direction let $g\cdot\theta=(g\Theta g^{\mathsf T},\mu')$, where
$\mu'_k=\mu_j$ if $g\Gamma_j=\Gamma_k$. The substitution $x\mapsto gx$ gives
\begin{equation}\label{p-eq:equivar}
  D_\theta(u^g,v^g)=D_{g\cdot\theta}(u,v),\qquad q_0[v^g]=q_0[v],\qquad \|v^g\|=\|v\| .
\end{equation}
Since $\lambda_0$ is simple and $\psi_0>0$, $\psi_0^g=\psi_0$. Hence $N[w^g]=N[w]$,
$w\mapsto w^g$ preserves $\{w\perp\psi_0\}$, and
\begin{equation}\label{p-eq:Rinv}
  R(g\cdot\theta)=R(\theta)\qquad(g\in\mathcal G).
\end{equation}
Let $\mathsf Pw=\frac16\sum_{g\in\mathcal G}w^g$.

\begin{lemma}\label{p-lem:sym}
\begin{enumerate}
\item[\textup{(a)}] If $\sum_{g\in\mathcal G}g\cdot\theta=0$, then $D_\theta[\psi_0]=0$ and
$D_\theta(\psi_0,\mathsf Pw)=0$ for all $w\in H^1(E)$.
\item[\textup{(b)}] If $g\cdot\theta=\theta$ for all $g\in\mathcal G$, then
$D_\theta(\psi_0,w)=D_\theta(\psi_0,\mathsf Pw)$ for all $w\in H^1(E)$.
\item[\textup{(c)}] If $\theta_A$ is as in \textup{(b)} and $\theta_E$ as in \textup{(a)}, then
$R(\theta_A+\theta_E)\le R(\theta_A)+R(\theta_E)$.
\item[\textup{(d)}] For $\phi\in\R$ let $c_j(\phi)=\cos(2\phi-2\vartheta_j)$ and
$\Theta_\phi=\bigl(-2S_\phi,(c_j(\phi))_{j}\bigr)$. Then $\sum_gg\cdot\Theta_\phi=0$, and
$R(\Theta_\phi)$ does not depend on $\phi$. We denote its value by $R_1=R_1(\alpha)$.
\end{enumerate}
\end{lemma}

\begin{proof}
By \eqref{p-eq:equivar} and $\psi_0^g=\psi_0$,
$D_\theta(\psi_0,w^g)=D_\theta(\psi_0^g,w^g)=D_{g\cdot\theta}(\psi_0,w)$; averaging over $g$
gives $D_\theta(\psi_0,\mathsf Pw)=\frac16D_{\sum_gg\cdot\theta}(\psi_0,w)$, which is $0$ in case
(a) and $D_\theta(\psi_0,w)$ in case (b). With $w=\psi_0=\mathsf P\psi_0$, (a) follows.

(c) Let $w\perp\psi_0$. Since $(w^g)^h=w^{gh}$ and $N$ is invariant, the bilinear form of $N$
satisfies $N(w^g,w^h)=N(w^{gh^{-1}},w)$, so $N(\mathsf Pw,\mathsf Pw)=N(\mathsf Pw,w)$, that is
$N[w]=N[\mathsf Pw]+N[w-\mathsf Pw]$; also $\mathsf Pw\perp\psi_0$. By (a) and (b),
\begin{align*}
  |D_{\theta_A+\theta_E}(\psi_0,w)|&=|D_{\theta_A}(\psi_0,\mathsf Pw)+D_{\theta_E}(\psi_0,w-\mathsf Pw)|\\
  &\le R(\theta_A)^{1/2}N[\mathsf Pw]^{1/2}+R(\theta_E)^{1/2}N[w-\mathsf Pw]^{1/2}
  \le\bigl(R(\theta_A)+R(\theta_E)\bigr)^{1/2}N[w]^{1/2}.
\end{align*}

(d) Let $g$ be the rotation through $2\pi/3$. Then $gS_\phi g^{\mathsf T}=S_{\phi+2\pi/3}$
(the matrix $S_\phi$ has the eigenvector $(\cos\phi,\sin\phi)$ for the eigenvalue $1$), and
$g\Gamma_j$ is the side with direction angle $\vartheta_j+2\pi/3$, on which the new coefficient
is $c_j(\phi)=\cos(2(\phi+2\pi/3)-2(\vartheta_j+2\pi/3))$. Hence
$g\cdot\Theta_\phi=\Theta_{\phi+2\pi/3}$. Similarly the reflection $x_1\mapsto-x_1$ maps
$\Theta_\phi$ to $\Theta_{-\phi}$. So the orbit of $\Theta_\phi$ consists of
$\Theta_{\pm\phi+2\pi k/3}$, $k=0,1,2$, and since $\Theta_\phi$ is a linear combination of
$\cos2\phi$ and $\sin2\phi$ with fixed coefficients $\Theta_0,\Theta_{\pi/4}$, the sum over the
orbit vanishes. By \eqref{p-eq:seminorm},
$R(\Theta_\phi)=R(\cos2\phi\,\Theta_0+\sin2\phi\,\Theta_{\pi/4})=a_0+a_1\cos4\phi+a_2\sin4\phi$
with constants $a_i$, and by \eqref{p-eq:Rinv} this function has period $2\pi/3$. A function
$a_1\cos4\phi+a_2\sin4\phi$ that is not identically zero has the periods $k\pi/2$,
$k\in\mathbb Z$, only; hence $a_1=a_2=0$.
\end{proof}

\subsection{An upper bound for \texorpdfstring{$R_1$}{R1} from lower bounds on a ring}

\begin{lemma}[ring]\label{p-lem:ring}
Let $0<s<\frac12$ and, for the two signs,
\[
  \Lambda_\pm=\Lambda\bigl(I\mp2sS_0,\ \alpha(1\pm s\cos2\vartheta_j)_{j=1,2,3}\bigr),
\]
the lowest eigenvalues of the forms $q_0\pm sD_{\Theta_0}$. Then $\Lambda_\pm\le\lambda_0$, and
if $\delta_\pm>0$ satisfy $\delta_\pm\ge\lambda_0-\Lambda_\pm$, then
\begin{equation}\label{p-eq:ring}
  R_1\ \le\ \frac{2+(\delta_++\delta_-)/(\underline\lambda_2-\lambda_0)}
  {s^2\,(1/\delta_++1/\delta_-)} .
\end{equation}
As a function of the four independent variables $\delta_+,\delta_-,\lambda_0,\underline\lambda_2$
(with $\underline\lambda_2>\lambda_0$), the right side is increasing in $\delta_+$, $\delta_-$ and
$\lambda_0$, and decreasing in $\underline\lambda_2$.
\end{lemma}

\begin{proof}
$D_{\Theta_0}[\psi_0]=0$ by Lemma~\ref{p-lem:sym}(a),(d), so the Rayleigh quotient of $\psi_0$ for
$q_0\pm sD_{\Theta_0}$ is $\lambda_0$. Let $w\perp\psi_0$, $F=D_{\Theta_0}(\psi_0,w)$ and
$c\in\R$. Since $q_0(\psi_0,w)=\lambda_0\langle\psi_0,w\rangle=0$,
\begin{align*}
  0&\le(q_0\pm sD_{\Theta_0})[c\psi_0+w]-\Lambda_\pm\|c\psi_0+w\|^2\\
  &=c^2(\lambda_0-\Lambda_\pm)\pm2csF+N[w]\pm sD_{\Theta_0}[w]+(\lambda_0-\Lambda_\pm)\|w\|^2\\
  &\le c^2\delta_\pm\pm2csF+N[w]\pm sD_{\Theta_0}[w]+\delta_\pm\|w\|^2 .
\end{align*}
Choosing $c=\mp sF/\delta_\pm$ gives
$s^2F^2/\delta_\pm\le N[w]\pm sD_{\Theta_0}[w]+\delta_\pm\|w\|^2$. Adding the two
inequalities, the terms $\pm sD_{\Theta_0}[w]$ cancel, and with \eqref{p-eq:Ngap},
\[
  s^2F^2\Bigl(\frac1{\delta_+}+\frac1{\delta_-}\Bigr)\le2N[w]+(\delta_++\delta_-)\|w\|^2
  \le N[w]\Bigl(2+\frac{\delta_++\delta_-}{\underline\lambda_2-\lambda_0}\Bigr).
\]
The monotonicity statements are evident.
\end{proof}

The point of using both signs is that the first-order error terms $\pm sD_{\Theta_0}[w]$
cancel exactly; the bound \eqref{p-eq:ring} exceeds $R_1$ only by terms of relative order $s^2$
and by the defect of the bounds $\delta_\pm$ divided by $s^2$. In a floating-point comparison
with $R_1$ computed by $P_2$ finite elements, with $s=1/10$ and the certified values of $\delta_\pm$
and $\underline\lambda_2$ used below, the bound exceeds $R_1$ by $\clpNodeLo$ to $\clpNodeHi$ per cent
at the grid points $\alpha_k$; the interpolation between grid points in Lemma~\ref{p-lem:Ccheck} raises
this to $\clpIvNear$ per cent near $\alpha=2$ (and to at most $\clpIvMax$ per cent, for small $\alpha$,
where the margin is large).

\subsection{Elementary shape bounds}

For $T_w$ with $|w|=\tanh2t$ and $w=|w|(\cos2\phi,\sin2\phi)$ write $\ell_j=\ell_j(G(w))$,
$\bar\ell=\frac13\sum_j\ell_j$, $c_j=c_j(\phi)$ and
\[
  \rho_j=\ell_j-\bar\ell-\tfrac12\sinh(2t)\,c_j .
\]
By Lemma~\ref{p-lem:shape}(d), $u_j:=\ell_j^2-1=(\cosh2t-1)+\sinh2t\,c_j$. From
$\ell-1=(\ell^2-1)/(\ell+1)$ one gets the identity
$\ell_j-1=\tfrac12u_j-u_j^2/(2(1+\ell_j)^2)$, and since $\sum_jc_j=0$,
\begin{equation}\label{p-eq:lbar}
  \bar\ell-1=\tfrac12(\cosh2t-1)-t^2\,\bar v,\qquad \rho_j=-t^2\,(v_j-\bar v),\qquad
  v_j=\frac{(u_j/t)^2}{2(1+\ell_j)^2},\quad\bar v=\tfrac13\sum_jv_j .
\end{equation}

\begin{lemma}\label{p-lem:shapeconst}
Let $t_*=3/50$. For $0<t\le t_*$ and all $\phi$,
\[
  \bar\ell-1\ \ge\ c_\ell\,t^2,\qquad \max_j|\rho_j|\ \le\ \hat\rho\,t^2,\qquad
  c_\ell=\clpCell,\quad\hat\rho=\clpRho .
\]
\end{lemma}

\begin{proof}
Replacing $\phi$ by $\phi+\pi/3$ permutes the numbers $c_j$, so it suffices to consider
$\phi\in[0,\pi/3]$. The functions $(\cosh2t-1)/t^2$ and $\sinh(2t)/t$ are power series in $t^2$
with positive coefficients, hence increasing in $t>0$ with limit $2$ at $0$, and
$u_j/t=t\,(\cosh2t-1)/t^2+c_j\sinh(2t)/t$. The script \texttt{local\_check.py} evaluates
\eqref{p-eq:lbar} in ball arithmetic on a subdivision of $(0,t_*]\times[0,\pi/3]$ into
$40\times60$ boxes, using these monotonicity properties on each $t$-interval, and obtains the
bounds $c_\ell\ge\clpCellComp$ and $\hat\rho\le\clpRhoComp$, which imply the stated ones. (The
independent verifier of Section~\ref{sec:verif}, with a finer subdivision, obtains $\clpCellIndep$ and
$\clpRhoIndep$.)
\end{proof}

\subsection{The local proposition}

\begin{proposition}\label{p-prop:local}
Let $A=1$, $\alpha\ge1/100$ and $0<|w|\le\tanh(2t_*)$, $t_*=3/50$. Then
$\lambda_1(T_w,\alpha)>\lambda_0(\alpha)$.
\end{proposition}

\begin{proof}
In this proof $w$ is the shape parameter and test functions in $H^1(E)$ are called $v$.
Let $|w|=\tanh2t$, $0<t\le t_*$, and $w=|w|(\cos2\phi,\sin2\phi)$. We apply
Lemma~\ref{p-lem:schur} (with $\underline\lambda_2$ in place of $\lambda_2$, see
Remark~\ref{p-rem:lam2}) to $(G,m)=(G(w),\alpha\ell)$, $\ell=(\ell_j(G(w)))_j$, for which
$\Lambda(G,m)=\lambda_1(T_w,\alpha)$. The eigenvalues of $G(w)$ are $e^{\pm2t}$ and
$\ell_j\ge e^{-t}$, so $q_{G,m}\ge(1-\eta)q_0$ with $\eta=1-e^{-2t}$, and $\eta'=\eta\gamma$.
By Lemma~\ref{p-lem:shape}(d), $D=q_{G,m}-q_0=D_\theta$ with $\theta=\theta_A+\theta_E$,
\[
  \theta_A=\bigl((\cosh2t-1)I,\ (\bar\ell-1)(1,1,1)\bigr),\qquad
  \theta_E=\bigl(-\sinh2t\,S_\phi,\ (\ell_j-\bar\ell)_j\bigr)
  =\tfrac12\sinh(2t)\,\Theta_\phi+(0,\rho).
\]
$\theta_A$ is invariant under $\mathcal G$; $\Theta_\phi$ has vanishing orbit sum by
Lemma~\ref{p-lem:sym}(d), and so has $(0,\rho)$ because $\sum_j\rho_j=0$.

\emph{The gain.} By Lemma~\ref{p-lem:sym}(a), $D[\psi_0]=D_{\theta_A}[\psi_0]$. Hence, by
Lemma~\ref{p-lem:lam0} (recall $\|\psi_0\|=1$), $\cosh2t-1\ge2t^2$ and Lemma~\ref{p-lem:shapeconst},
\begin{equation}\label{p-eq:gain}
  D[\psi_0]=(\cosh2t-1)(\lambda_0-\alpha\lambda_0')+\alpha(\bar\ell-1)\lambda_0'
  \ \ge\ t^2\bigl(2\lambda_0-(2-c_\ell)\,\alpha\lambda_0'\bigr).
\end{equation}

\emph{The invariant part.} $D_{\theta_A}$ is a nonnegative form, and
$D_{\theta_A}[v]\le(\cosh2t-1)\,q_0[v]$ because
$0\le\bar\ell-1\le(\frac13\sum_j\ell_j^2)^{1/2}-1=(\cosh2t)^{1/2}-1\le\cosh2t-1$. By the
Cauchy--Schwarz inequality for $D_{\theta_A}$ and \eqref{p-eq:Ngap},
$D_{\theta_A}(\psi_0,v)^2\le D_{\theta_A}[\psi_0]\,(\cosh2t-1)\gamma N[v]$ for $v\perp\psi_0$, that is
$R(\theta_A)\le(\cosh2t-1)\,\gamma\,D[\psi_0]$.

\emph{The remainder $(0,\rho)$.} For $v\perp\psi_0$, by \eqref{p-eq:Ngap} and
$\int_{\partial E}\psi_0^2=\lambda_0'$,
\[
  |D_{(0,\rho)}(\psi_0,v)|\le\max_j|\rho_j|\,\Bigl(\alpha\int_{\partial E}\psi_0^2\Bigr)^{1/2}
  \Bigl(\alpha\int_{\partial E}v^2\Bigr)^{1/2}\le\max_j|\rho_j|\,(\alpha\lambda_0'\,\gamma N[v])^{1/2}.
\]
Hence, by \eqref{p-eq:seminorm}, Lemma~\ref{p-lem:sym}(d) and Lemma~\ref{p-lem:shapeconst},
\[
  R(\theta_E)^{1/2}\le\tfrac12\sinh(2t)R_1^{1/2}+\hat\rho\,t^2(\alpha\lambda_0'\gamma)^{1/2}
  \le t\,\bigl(s_*R_1^{1/2}+\hat\rho\,t_*(\alpha\lambda_0'\gamma)^{1/2}\bigr),
  \qquad s_*=\frac{\sinh2t_*}{2t_*}.
\]

\emph{Conclusion.} By Lemma~\ref{p-lem:sym}(c), $\cR=R(\theta)\le R(\theta_A)+R(\theta_E)$. With
$\eta_*=1-e^{-2t_*}$ and $\varepsilon_*=\cosh2t_*-1$, the hypothesis
$D[\psi_0](1-\eta')>\cR$ of Lemma~\ref{p-lem:schur} therefore follows from
$D[\psi_0](1-(\eta_*+\varepsilon_*)\gamma)>R(\theta_E)$, and by \eqref{p-eq:gain} this follows
from the inequality
\begin{equation}\label{p-eq:C}\tag{C}
  \bigl(2\lambda_0-(2-c_\ell)\,\alpha\lambda_0'\bigr)\bigl(1-(\eta_*+\varepsilon_*)\gamma\bigr)
  \ >\ \Bigl(s_*\overline R_1^{1/2}+\hat\rho\,t_*(\alpha\lambda_0'\gamma)^{1/2}\Bigr)^2
\end{equation}
for some $\overline R_1\ge R_1(\alpha)$ and some admissible $\underline\lambda_2$. (The first
factor on the left of \eqref{p-eq:C} is at least $c_\ell\lambda_0>0$, because
$\alpha\lambda_0'\le\lambda_0$; so \eqref{p-eq:C} forces $(\eta_*+\varepsilon_*)\gamma<1$, in
particular $\eta\gamma<1$, which is the condition $\eta<1-\lambda_0/\underline\lambda_2$ of
Lemma~\ref{p-lem:schur}.) Inequality \eqref{p-eq:C}
involves $\alpha$ only. It is verified for all $\alpha\ge1/100$ in Lemma~\ref{p-lem:Ccheck}.
\end{proof}

\begin{lemma}\label{p-lem:Ccheck}
For every $\alpha\ge1/100$ there are $\underline\lambda_2\in(\lambda_0(\alpha),\lambda_2(E,\alpha)]$
and $\overline R_1\ge R_1(\alpha)$ such that \eqref{p-eq:C} holds.
\end{lemma}

\begin{proof}
Let $\alpha_k$, $k=0,1,\dots$, be the grid of rational numbers
$\alpha_k=\frac1{100}2^{\lfloor k/16\rfloor}r_{k\bmod16}$, where $r_0=1$ and $r_i$ is $2^{i/16}$
rounded to six decimals; so $\alpha_0=1/100$ and $\alpha_{k+1}/\alpha_k\approx1.0443$. Let
$K=336$, so that $\alpha_K=0.01\cdot2^{21}$, and $s=1/10$.

\emph{Certified data.} (i) For $k=0,\dots,K$ and both signs the script \texttt{local\_ring.py}
computes numbers $\mu_\pm^k\le\Lambda_\pm(\alpha_k)$ (notation of Lemma~\ref{p-lem:ring}) by the
method of Section~\ref{p-sec:cert}: by \eqref{p-eq:pullback}, $\Lambda_\pm$ is the lowest
eigenvalue of an acute isosceles triangle with side-dependent Robin coefficients, to which
Lemmas~\ref{p-lem:CR} and~\ref{p-lem:CW} are applied on the mesh (U) with $n=N_{\rm ring}$
($N_{\rm ring}=\clpRingN$ and $N_2=\clpLamN$ below for $\kappa_0=0.2983$; see Table~\ref{p-tab:runs}).
(ii) For $k\in\{0,4,8,\dots,K\}$ it computes numbers $\underline\lambda_2^k\le\lambda_2(E,\alpha_k)$
by Lemmas~\ref{p-lem:CR} ($k=2$) and~\ref{p-lem:inertia} on the mesh (U) of $E$ with $n=N_2$.
(iii) $\lambda_0(\alpha_k)$ and $\lambda_0'(\alpha_k)$ are enclosed using
Proposition~\ref{p-prop:equi} and \eqref{p-eq:lam0prime}.

\emph{An interval $I_k=[\alpha_k,\alpha_{k+1}]$, $0\le k<K$.} The functions
$\alpha\mapsto\Lambda_\pm(\alpha)$ are concave (Lemma~\ref{p-lem:concave}), so on $I_k$ they lie
above the chord through $(\alpha_k,\mu_\pm^k)$, $(\alpha_{k+1},\mu_\pm^{k+1})$. By
Lemma~\ref{p-lem:lam0}, $\lambda_0$ lies below the two tangents at $\alpha_k$ and $\alpha_{k+1}$.
The difference between the minimum of the two tangents and the chord is piecewise linear on
$I_k$ with one kink, so its maximum $\delta_\pm$ is attained at $\alpha_k$, $\alpha_{k+1}$ or at
the kink, and $\delta_\pm\ge\lambda_0(\alpha)-\Lambda_\pm(\alpha)$ on $I_k$; in particular
$\delta_\pm\ge\lambda_0(\alpha_k)-\Lambda_\pm(\alpha_k)\ge0$, and the script checks $\delta_\pm>0$. The second
eigenvalue $\lambda_2(E,\cdot)$ is nondecreasing, so
$\underline\lambda_2:=\underline\lambda_2^{4\lfloor k/4\rfloor}$ is admissible on $I_k$ provided
it exceeds $\lambda_0(\alpha_{k+1})$. On $I_k$ we further have
$\lambda_0(\alpha_k)\le\lambda_0\le\lambda_0(\alpha_{k+1})$ and
$\alpha\lambda_0'\le\alpha_{k+1}\lambda_0'(\alpha_k)$. By the monotonicity statement of
Lemma~\ref{p-lem:ring}, \eqref{p-eq:C} holds on $I_k$ if it holds with $\lambda_0$ replaced by
$\lambda_0(\alpha_k)$ in the first factor, $\alpha\lambda_0'$ by
$\alpha_{k+1}\lambda_0'(\alpha_k)$, $\gamma$ by
$\underline\lambda_2/(\underline\lambda_2-\lambda_0(\alpha_{k+1}))$ and $\overline R_1$ by the
right side of \eqref{p-eq:ring} evaluated with $\delta_\pm$, $\underline\lambda_2$ and
$\lambda_0(\alpha_{k+1})$.

\emph{The interval $[\alpha_K,\infty)$.} Here $\Lambda_\pm(\alpha)\ge\Lambda_\pm(\alpha_K)\ge\mu_\pm^K$
by monotonicity and $\lambda_0(\alpha_K)\le\lambda_0<\Lambda_1(E)=4\pi^2/\sqrt3$, so we take
$\delta_\pm=\Lambda_1(E)-\mu_\pm^K>0$, $\underline\lambda_2=\underline\lambda_2^K$ (which
exceeds $\Lambda_1(E)$), replace
$\lambda_0$ by $\Lambda_1(E)$ in $\gamma$ and in \eqref{p-eq:ring}, and use
$\alpha\lambda_0'\le2(\Lambda_1(E)-\lambda_0(\alpha_K/2))$ (Lemma~\ref{p-lem:lam0}).

The script \texttt{local\_check.py} evaluates these $K+1=\clpIneqC$ inequalities in ball arithmetic.
All hold; the smallest quotient of the left by the right side of \eqref{p-eq:C} is $\clpQmin$, attained
for $k=\clpQminK$ ($\alpha\approx\clpQminAlpha$), and the quotient exceeds $\clpQgeThirty$ for
$\alpha\ge30$ and $\clpQleSix$ for $\alpha\le0.6$.
\end{proof}

\begin{remark}\label{p-rem:local}
Floating-point values from $P_2$ finite elements on the grid $\alpha_0,\dots,\alpha_K$ (not used in
the proof): with $a(\alpha)=2\lambda_0-\frac54\alpha\lambda_0'$, the leading coefficient of
$D[\psi_0]/t^2$ as $t\to0$, the ratio $R_1/a$ is $\clpRaLow$ at $\alpha=0.1$, has its maximum
$\clpRaMax$ near $\alpha\approx\clpRaMaxAlpha$, and is $\clpRaLast$ at $\alpha_K\approx2\cdot10^4$; so the
criterion has a factor of about $\clpSpare$ to spare as $t\to0$, uniformly in $\alpha$, of which the
factor $1-(\eta_*+\varepsilon_*)\gamma$ and the over-estimation of $R_1$ consume a part at $t=t_*$. The
Hessian of $T\mapsto\lambda_1(T,\alpha)/\lambda_0(\alpha)$ at $E$ in the variable $t$ is
$2(a-R_1)/\lambda_0$, between $\clpHessMin$ (near $\alpha\approx\clpHessMinAlpha$) and $\clpHessMax$.
\end{remark}


\section{The middle region}\label{p-sec:middle}

Let $A=1$, $r_*=0.1194<\tanh(2t_*)=0.11942\ldots$, and
\[
  F(w)=(1+w_1)^2-81\,(1-|w|^2),\qquad
  \Omega_B=\bigl\{w\in\mathbb D:\ 0<\arg w<\pi/3,\ |w|>r_*,\ F(w)<0\bigr\}.
\]
By Lemma~\ref{p-lem:shape}(b),(c), Theorem~\ref{p-thm:thin} and Proposition~\ref{p-prop:local}, the
inequality $\lambda_1(T_w,\alpha)>\lambda_0(\alpha)$ for $\alpha\ge1/100$ remains to be proved
only for $w$ in the closure of $\Omega_B$ in $\mathbb D$.

\begin{lemma}[coverage]\label{p-lem:cover}
Let $\mathcal T$ be a finite set of closed, positively oriented, nondegenerate triangles in the
plane such that every side $[p,q]$ of a triangle of $\mathcal T$, traversed from $p$ to $q$ in
its boundary, either occurs exactly once more, as the side $[q,p]$ of another triangle, or does
not occur again in either direction (a \emph{boundary side}). Let $\Omega$ be a connected open
set that is disjoint from all boundary sides and contains a point that belongs to the interior
of some triangle of $\mathcal T$ and to no side of any triangle of $\mathcal T$. Then $\overline\Omega\subset\bigcup\mathcal T$.
\end{lemma}

\begin{proof}
For a point $x$ on no side, the number $n(x)$ of triangles of $\mathcal T$ containing $x$ is
the sum of the winding numbers of their boundaries around $x$, that is the winding number
around $x$ of the cycle formed by all oriented sides. The contributions of the sides that occur
in both directions cancel, so $n(x)$ is the winding number of the cycle of boundary sides, which
is locally constant on the complement of the boundary sides, in particular on $\Omega$; there
it equals its value at the given point, which is at least $1$. So every point of $\Omega$ on no
side lies in $\bigcup\mathcal T$; these points are dense in $\Omega$ and $\bigcup\mathcal T$ is
closed.
\end{proof}

\paragraph{The cells.} The file \texttt{coverB\_roots.json} contains a list $\mathcal T$ of
$\clpRoots$ triangles in $\mathbb D$ whose vertices have rational coordinates with denominator
$2^{20}$: the points $\tanh(2t_i)(\cos\psi,\sin\psi)$, rounded, for thirteen radii
$t_i\in[0.057,1.285]$ and angles $\psi\in\{0,\frac14\psi_{\max},\dots,\psi_{\max}\}$,
$\psi_{\max}=\pi/3+10^{-3}$ (with four additional points on the outermost circle), joined to a
triangulation. The grid of parameters $\alpha_k$ is the one of Lemma~\ref{p-lem:Ccheck},
continued to $k\le400$.

\begin{proposition}\label{p-prop:middle}
\begin{enumerate}
\item[\textup{(a)}] $\overline{\Omega_B}\cap\mathbb D\subset\bigcup\mathcal T$.
\item[\textup{(b)}] For every $\Delta\in\mathcal T$, either $F\ge0$ on $\Delta$, or there are
indices $0=k_0<k_1<\dots<k_p\le400$ and, for each vertex $w_i$ of $\Delta$ and each
$k\in\{k_0,\dots,k_p\}$, a certified number $b_i^k\le\lambda_1(T_{w_i},\alpha_k)$, such that
the hypothesis of Corollary~\ref{p-cor:prism}\textup{(a)} holds on each
$[\alpha_{k_{q-1}},\alpha_{k_q}]$, with $U$ the minimum of the tangents of $\lambda_0$ at the
grid points in that interval, and the hypothesis of Corollary~\ref{p-cor:prism}\textup{(b)}
holds at $\alpha_{k_p}$.
\end{enumerate}
Consequently $\lambda_1(T_w,\alpha)>\lambda_0(\alpha)$ for all $w\in\overline{\Omega_B}\cap\mathbb D$
and all $\alpha\ge1/100$.
\end{proposition}

\begin{proof}
(a) In exact rational arithmetic one checks that $\mathcal T$ satisfies the hypotheses of
Lemma~\ref{p-lem:cover} and that each of its $\clpBsides$ boundary sides lies in one of the closed sets
$\{|w|\le r_*\}$, $\{w_2\le0\}$, $\{w_2\ge\sqrt3\,w_1\}$, $\{F\ge0\}$, all of which are disjoint
from $\Omega_B$. (For the first three sets, which are convex, it suffices to test the end
points; $F$ is a convex quadratic polynomial, and its minimum on a segment is computed
exactly.) The set $\Omega_B$ is connected: for $0<\psi<\pi/3$ the function
$r\mapsto F(r\cos\psi,r\sin\psi)$ is increasing on $[0,1]$, negative at $r_*$ and positive at
$1$, so $\Omega_B=\{r_*<r<r_{\max}(\psi)\}$ in polar coordinates with a continuous
$r_{\max}>r_*$. Finally the point $(\frac12,\frac15)\in\Omega_B$ is interior to a triangle of
$\mathcal T$ and lies on no line that contains a side of a triangle of $\mathcal T$.

(b) is the result of the computation described in Section~\ref{p-sec:cert} and below, carried
out by \texttt{cover.py}: one triangle (adjacent to the point $w=(\tanh2.57,0)$) is recorded
with $F\ge0$; for each of the other $\clpCellsB$ (\clpThinBword\ of which also lie in $\{F\ge0\}$) the
certificate lists the indices $k_q$ and the bounds $b_i^k$. The vertex bounds are Crouzeix--Raviart bounds on the meshes (U) or (A) (the latter
whenever the smallest cotangent of $T_{w_i}$ is not certified to exceed $0.02$). The function
$U$ is piecewise linear, $U\ge\lambda_0$, and by Lemma~\ref{p-lem:lam0} its kinks are the
intersection points of consecutive tangents, one in each grid interval (the script encloses
each of them in its grid interval). Both sides of \eqref{p-eq:prism} are then piecewise linear in
$\alpha$; so \eqref{p-eq:prism} is verified, in ball arithmetic, at the end points and at these
kinks, for the six pairs $i\le k$. The steps $k_{q-1}\to k_q$ were chosen greedily as large as
possible, so some of these inequalities hold with a small relative margin (the smallest is
$\clpMargin$); they are decided by $100$-bit ball arithmetic applied to the stored exact
bounds.

The consequence follows from Corollary~\ref{p-cor:prism}, from Lemma~\ref{p-lem:shape}(c) and
Theorem~\ref{p-thm:thin} for the triangle with $F\ge0$, and from (a).
\end{proof}

In total the certificate for (b) consists of $\clpPrisms$ prisms and uses $\clpNodes$ vertex
bounds, with mesh parameters $n\le\clpNmax$ (at most $\clpNdof$ unknowns); the largest grid index
needed is $k_p=\clpKtop$ ($\alpha\approx\clpAtop$) for the cells closest to $E$. No cell had to be
subdivided.


\section{Proof of Theorem~\ref{thm:main}(1); the computations}\label{p-sec:proof}

\begin{proof}[Proof of Theorem~\ref{thm:main}\textup{(1)}]
By scaling, $A=1$. Let $T$ be a non-equilateral triangle of area $1$ and $\alpha>0$. If
$\alpha\le1/100$, Theorem~\ref{p-thm:uniform} applies. Let $\alpha\ge1/100$. By
Lemma~\ref{p-lem:shape}(b), $T$ is congruent to $T_w$ with $w\in\Sigma\setminus\{0\}$. If
$|w|\le r_*<\tanh(2t_*)$, Proposition~\ref{p-prop:local} applies. If $F(w)\ge0$, then
$a_{\max}\ge P_E$ by Lemma~\ref{p-lem:shape}(c) and Theorem~\ref{p-thm:thin} applies. Otherwise
$w\in\overline{\Omega_B}\cap\mathbb D$ and Proposition~\ref{p-prop:middle} applies.
\end{proof}

\begin{proof}[Proof of Corollary~\ref{cor:perimeter} for $\alpha>0$]
Let $E_s$ be the equilateral triangle of area $s^2|E|$. By the remark after
Lemma~\ref{p-lem:concave}, $\alpha\mapsto\lambda_1(E,\alpha)$ is concave, and it vanishes at $\alpha=0$; so for
$s>1$, $\lambda_1(E_s,\alpha)=s^{-2}\lambda_1(E,s\alpha)\le s^{-1}\lambda_1(E,\alpha)<\lambda_1(E,\alpha)$:
the eigenvalue of the equilateral triangle decreases strictly with its area. By the isoperimetric
inequality for triangles, $|T|\le|\E^T|$. Hence
$\lambda_1(T,\alpha)\ge\lambda_1(\E_T,\alpha)\ge\lambda_1(\E^T,\alpha)$, and the first inequality is strict
unless $T$ is equilateral, by Theorem~\ref{thm:main}(1).
\end{proof}

\paragraph{What is computed.} The proof uses the following machine computations, in
\texttt{code/positive/} of the repository; Section~\ref{sec:verif} describes how they are re-run and
checked.
\begin{enumerate}
\item \texttt{weak\_cert.py}: the interval check of the two-variable inequality in the proof of
Theorem~\ref{p-thm:uniform} ($\clpWeakBoxes$ boxes).
\item \texttt{local\_ring.py}: the $\clpRingRec=2\cdot\clpIneqC$ ring bounds $\mu_\pm^k$ and the
$\clpLamRec$ bounds $\underline\lambda_2^k$ of Lemma~\ref{p-lem:Ccheck}; \texttt{local\_check.py}: the
constants of Lemma~\ref{p-lem:shapeconst} and the $\clpIneqC$ inequalities \eqref{p-eq:C}.
\item \texttt{cover.py}: the vertex bounds and the prism inequalities of
Proposition~\ref{p-prop:middle}(b). The exact rational checks of
Proposition~\ref{p-prop:middle}(a) are done by \texttt{coverage\_check.py}.
\end{enumerate}
The certified statements are of three kinds: (1) sign tests on exact integers
(Lemma~\ref{p-lem:CW}); (2) inequalities between elementary functions of rational arguments,
evaluated in Arb ball arithmetic \cite{Arb} (through python-flint, $100$ bits), including
Gershgorin tests on small dense ball matrices (Lemma~\ref{p-lem:inertia}); (3) exact rational
arithmetic (the coverage). Floating-point arithmetic is used only to \emph{propose} objects whose
properties are then certified: approximate eigenvectors, an approximate inverse Cholesky
factor, the trial values of $\sigma$, the positions of the grid points and the choice of the mesh
parameters and of the side on which the altitude is dropped. Every certified
lower bound is stored after a further relative reduction of the discrete eigenvalue bound
$\sigma$ by $10^{-7}$ (by $10^{-6}$ for the second eigenvalue), which leaves room for an
independent re-verification with a different eigenvector. In all, the default run solves
$\clpEigCount$ sparse eigenvalue problems, in less than half an hour on one core.

\paragraph{Three interpolation constants.} The only external analytic input of the certified
lower bounds is the interpolation constant of Lemma~\ref{p-lem:interp}. The whole computation was
carried out three times, with three constants $\kappa_0$ in Lemma~\ref{p-lem:CR}:
$0.2983$ (Lemma~\ref{p-lem:interp} with the Laugesen--Siudeja inequality; this is the default run,
described in the text), $0.3496$ (Lemma~\ref{p-lem:interp} with the Payne--Weinberger inequality)
and $0.1893$ (the value obtained by Liu~\cite{Liu}, which is close to optimal; it is not proved by
hand but by a verified computation). Each run by itself proves Theorem~\ref{thm:main}(1); they differ
only in the mesh parameters. Table~\ref{p-tab:runs} summarises them.

\begin{table}[ht]\centering
\begin{tabular}{lccc}\toprule
constant $\kappa_0$ & $0.2983$ & $0.3496$ & $0.1893$\\\midrule
ring bounds: $n$ (unknowns) & \clpRingN\ (\clpRingNN) & \clpRingNpw\ (\clpRingNNpw) & \clpRingNliu\ (\clpRingNNliu)\\
$\underline\lambda_2$ bounds: $n$ (unknowns) & \clpLamN\ (\clpLamNN) & \clpLamNpw\ (\clpLamNNpw) & \clpLamNliu\ (\clpLamNNliu)\\
min.\ quotient in \eqref{p-eq:C} & \clpQmin & \clpQminpw & \clpQminliu\\
cells by prisms / thin & \clpCellsB\ / \clpCellsW & \clpCellsB\ / \clpCellsW & \clpCellsB\ / \clpCellsW\\
prisms & \clpPrisms & \clpPrismspw & \clpPrismsliu\\
vertex bounds used & \clpNodes & \clpNodespw & \clpNodesliu\\
largest $n$ in the cover (unknowns) & \clpNmax\ (\clpNdof) & \clpNmaxpw\ (\clpNdofpw) & \clpNmaxliu\ (\clpNdofliu)\\
largest $k_p$ ($\alpha_{k_p}$) & \clpKtop\ (\clpAtop) & \clpKtoppw\ (\clpAtoppw) & \clpKtopliu\ (\clpAtopliu)\\
time: ring, $\underline\lambda_2$, cover (min) & \clpTimels & \clpTimepw & \clpTimeliu\\\bottomrule
\end{tabular}
\caption{The three runs. ``min.\ quotient'' is the smallest quotient of the two sides of
\eqref{p-eq:C}; times are for one core of a loaded server.}\label{p-tab:runs}
\end{table}


\section{Implementation, verification and caveats}\label{sec:verif}

\paragraph{Implementation.} All certified steps use Arb~\cite{Arb} through python-flint~0.9.0: at
$128$ bits in Part~I, and at $100$ bits ($80$ bits in \texttt{weak\_cert.py}) in Part~II; Part~II also
uses exact integer and rational arithmetic. In Part~I, the certificate of Theorem~\ref{thm:A} consists of
the $\clthmAstepOne$ intervals of Lemma~\ref{lem:gq}, $\clthmAsub$ Hessian sub-boxes and
$\clthmAboxes$ boxes in Lemma~\ref{lem:geomA}, and that of Theorem~\ref{thm:C} of $\clthmCcells$
intervals and one tail box; each takes well under a second. Box endpoints are exact binary or rational
numbers, and every certified inequality is strict. The certificates of Part~II (the default run and the
two other runs of Table~\ref{p-tab:runs}, about $1.3$\,MB of JSON each) are stored in the repository and
are re-checked as described below; the data format is documented there.

Every number in this paper that a computation checks is kept in a single claims file. The script
\texttt{scripts/verify.py} reruns the computations, compares each computed enclosure or count with these
claims in exact rational arithmetic, and the macros of this paper are generated from the same file.
By default it reruns the certificates of Part~I, the second certificate of Lemma~\ref{lem:certC}, the
floating-point checks of \S\ref{sec:cross}, and for Part~II: \texttt{weak\_cert.py},
\texttt{local\_check.py} and \texttt{coverage\_check.py} from scratch, all prism and closing inequalities
from the stored bounds, the negative controls, and the independent verifier described below (with a
from-scratch re-certification of a sample of the stored bounds for the default run, and a replay of the
other two runs). Optional modes replay the six-wave cells of \S\ref{sec:six} and re-certify every stored
bound of all three runs of Part~II from scratch with both programs. Notation in the code of Part~I: $G_0$
is \texttt{D3}, $\Gp$, $\Bp$, $\Np$ are \texttt{G1}, \texttt{B1}, \texttt{N1}, the products $\beta\Bp$,
$\beta\chi$ and $|\lambda_0|\Np$ are \texttt{bB1}, \texttt{bkq} and \texttt{lN1}, and the $\Sigma_k$ are
\texttt{s1}, \dots, \texttt{s4}; in the second certificate, $\chi$ is \texttt{kappa} and $\beta\chi$ is
\texttt{bk}.

\paragraph{Checks of Part~II.} (i) \emph{Re-certification.} The script \texttt{recheck.py}
re-evaluates all prism inequalities and closing tests from the stored bounds and re-certifies every
stored bound from scratch: a freshly computed eigenvector must pass the exact test of
Lemma~\ref{p-lem:CW} for the \emph{stored} $\sigma$. (ii) \emph{Independent verifier.} A second program
was written by a separate AI instance from the text of Part~II and a description of the certificate
format only, without access to our code. It builds each triangle from coordinates and each mesh as an
explicit list of elements, assembles the Crouzeix--Raviart matrices element by element, computes its own
eigenvectors, and verifies $Kx\ge\sigma Mx$ for the stored $\sigma$ in outward-rounded floating-point
interval arithmetic; it treats the second eigenvalue by an interval $LDL^{\mathsf T}$ factorisation and
Sylvester's law of inertia, and re-derives the constants of Lemma~\ref{p-lem:shapeconst}, the
inequalities~\eqref{p-eq:C}, the prism inequalities and the coverage with a different interval library
(mpmath~\cite{mpmath}) and exact rationals. For each of the three runs it confirms every record: all
$\clpRingRec$ ring bounds, all $\clpLamRec$ bounds for $\lambda_2$, all stored vertex bounds, all
$\clpIneqC$ inequalities~\eqref{p-eq:C} (smallest quotient $\clpQmin$ at $k=\clpQminK$ for the default
run, as in our computation), all prisms and the coverage. (iii) \emph{Negative controls.} The two programs
reject perturbed statements; between them: a $\sigma$ larger by a factor $1+10^{-6}$ (resp.\
$1+10^{-4}$), mesh (U) on an obtuse triangle, $\lambda_0$ replaced by $1.01\lambda_0$ (the independent verifier rejects
$\clpIndepRej$ cells already for $1.0001\lambda_0$), the false bound
$\lambda_1(T_w,\alpha)\ge\lambda_0(1+0.6t^2)$ near $E$, $t_*$ replaced by $3/25$, ring bounds lowered by
$0.2$ per cent near $\alpha=2$, a cover with a triangle removed, a nonpositive $\sigma$ or $\sigma_2$
($-10^6$, $0$, and just below the pole $-1/C_h^2$ of $\sigma/(1+C_h^2\sigma)$), a vertex bound
$\le0$ in a cell, and, on the one-element mesh (U) of the triangle with sides $5,5,6$ and Robin
coefficients $6,6,5$ (discrete eigenvalue exactly $15/2$), the value $\sigma=15/2+2^{-61}$. (iv) \emph{Sanity checks in
floating point}, not used in the proof: the certified vertex bounds lie below conforming $P_2$ finite
element values, by relative amounts between $\clpSandLo$ and $\clpSandHi$ on random samples, and the
quantities of Remark~\ref{p-rem:local} are computed with $P_2$ elements. (v) An AI referee re-derived the
new arguments of Part~II line by line and reported no error; the gaps it found (the position of the kinks
of $U$, a point on no side in Lemma~\ref{p-lem:cover}, the use of the rounded constants
in~\eqref{p-eq:C}) are repaired. A second AI referee pass, made when the two parts were merged, again
checked the cell identity (Lemma~\ref{p-lem:cell}), the ring bound (Lemma~\ref{p-lem:ring}), the
uniformity in $\alpha$ (Lemma~\ref{p-lem:Ccheck}), the closing test of Corollary~\ref{p-cor:prism}(b),
Lemmas~\ref{p-lem:CR} and~\ref{p-lem:CW} and the coverage lemma, and the corresponding code, and found no
error. (vi) An external AI review (ChatGPT Pro) of the code found four defects, now corrected, none of
which affects a stored certificate: both checkers of the bounds for $\lambda_2$ (and our re-checks of
the other bounds) did not enforce $\sigma>0$ in Lemma~\ref{p-lem:CR}, so that a forged record with
a very negative $\sigma$ would have been accepted (all stored values of $\sigma$ and $\sigma_2$ are
positive); our re-check rounded a given $\sigma$ down to a multiple of $2^{-60}$ before the test of
Lemma~\ref{p-lem:CW} but then used the unrounded value (all stored values are such multiples); the
first cell of the second certificate of Lemma~\ref{lem:certC} started about $4\cdot10^{-17}$ above
$13/10$; and the reported lower end of the enclosure of $q(6/5)$ was evaluated at the upper end of the
enclosure of $M(6/5)$ (the upper bound $q(6/5)<\clqA$, which is what the proof uses, was correct). A
subsequent audit of both programs added explicit checks of further hypotheses: positive vertex bounds
in the prism tests (with a negative control), the ring parameter $s=1/10$ in our re-check, and $g\ge0$
in Lemma~\ref{p-lem:inertia}.

\paragraph{Check of prior work.} The repository also contains, clearly separated from our proof, the
AI referee's re-implementation of the middle-range criterion of Ma~\cite[Proposition~2]{Ma} in Arb, which
confirms it on $1\le\beta\le10$ with $\clmaBoxes$ boxes; it is run by \texttt{scripts/verify.py}.

\paragraph{Caveats.} The proofs are computer-assisted and, at the time of writing, have been checked
only by AI systems and by the programs described above; no human expert has reviewed them (see the
declaration below). In Part~I the main residual risks are the hand derivations, especially the closed
forms of the eight integrals of Lemma~\ref{lem:certC} and their transcription into code; they are
cross-checked against quadrature on actual triangles by two independently written scripts. In Part~II,
the local region has a comfortable margin (the smallest quotient in~\eqref{p-eq:C} is $\clpQmin$). In the
middle region no cell had to be subdivided, but the true gap $\lambda_1(T_w,\alpha)/\lambda_0(\alpha)-1$
is small next to the local disc (between $\clpGapLo$ and $\clpGapHi$ per cent at the inner vertices
$t=0.057$, in floating point), and individual prism inequalities hold with relative margins down to
$\clpMargin$; that region therefore rests entirely on the exactness of the Collatz--Wielandt test, of
the constants in Lemma~\ref{p-lem:CR} and of the ball arithmetic. The most delicate analytic point is the
neighbourhood of $E$: the criterion of Lemma~\ref{p-lem:schur} has only a factor of about $\clpSpare$ to
spare in the worst case, so the upper bound for $R_1$ has to be accurate to a few per cent;
Lemma~\ref{p-lem:ring} achieves this with two eigenvalue bounds per value of $\alpha$. We did not find a
computer-free bound for $R_1(\alpha)$. (By the ground state transformation
$N[\psi_0\chi]=\int_E\psi_0^2|\nabla\chi|^2$, one has $R_1=\min\int_E\psi_0^2|Y|^2$ over the vector fields
$Y$ with $\int_E\psi_0^2\,Y\cdot\nabla\chi=D_{\Theta_0}(\psi_0,\psi_0\chi)$ for all $\chi$ with
$\int_E\psi_0^2\chi=0$, so any explicit admissible field would give such a bound.) The only external
analytic inputs of Part~II are standard facts about the Robin form and the Poincar\'e inequality of
Laugesen and Siudeja~\cite{LaSiPoincare} (or of Payne and Weinberger~\cite{PW} in the second run); the
constant of the third run is Liu's~\cite{Liu}, which comes from a verified computation rather than a
hand proof, whereas the default run uses the constant $0.2983$ proved in Lemma~\ref{p-lem:interp}. The
remaining dependence is on the correctness of Arb/python-flint, of mpmath's interval arithmetic and of
Python's integers.

\paragraph{Availability.} The code, the certificates, the logs and the claims file are in the
accompanying repository\footnote{\url{https://github.com/jasondavies/robin-triangles}.}. The single
command \texttt{scripts/verify.py} reruns the checks described above and writes a PASS/FAIL summary.

\appendix
\section{A strong-coupling bound for \texorpdfstring{$\alpha>0$}{alpha > 0}}\label{p-sec:large}

This appendix uses the notation of Part~II and is independent of the proof of
Theorem~\ref{thm:main}. It gives an explicit lower bound for $\lambda_1(T,\alpha)$ of an individual
triangle as $\alpha\to\infty$.

\begin{theorem}[strong coupling]\label{p-thm:large}
Let $\psi_T$ be the torsion function of $T$ ($-\Delta\psi_T=1$ in $T$, $\psi_T=0$ on $\partial T$)
and $C_T=\sup|\partial_\nu\psi_T|\le h_{\max}/2$, the supremum being over the open sides. For every $\alpha>0$ the number
$\lambda=\lambda_1(T,\alpha)$ satisfies
\begin{equation}\label{p-eq:large}
  \alpha\ \le\ \frac{2\lambda^2}{r\,\Lambda_1(\Lambda_1-\lambda)}
  +C_T\,\frac{\lambda\,\Lambda_2}{\Lambda_2-\lambda},
  \qquad \Lambda_k=\Lambda_k(T).
\end{equation}
In particular, with $\alpha_*=C_T\Lambda_1\Lambda_2/(\Lambda_2-\Lambda_1)$,
\begin{equation}\label{p-eq:large2}
  \lambda_1(T,\alpha)\ \ge\ \Lambda_1(T)\Bigl(1-\frac{2}{r(\alpha-\alpha_*)}\Bigr)
  \qquad(\alpha>\alpha_*).
\end{equation}
\end{theorem}

\begin{lemma}\label{p-lem:harm}
If $h\in H^1(T)$ is harmonic in $T$, then $\int_Th^2\le C_T\int_{\partial T}h^2$, and
$C_T\le h_{\max}/2$.
\end{lemma}

\begin{proof}
Since $\psi_T\ge0$ vanishes on $\partial T$ and $\Delta(h^2)=2|\nabla h|^2$,
\[
  \int_Th^2=-\int_Th^2\Delta\psi_T
  =-\int_{\partial T}h^2\,\partial_\nu\psi_T-2\int_T\psi_T|\nabla h|^2
  \le C_T\int_{\partial T}h^2 .
\]
Here $C_T$ is understood as the supremum of $|\partial_\nu\psi_T|$ over the open sides, which
is finite by the bound proved below, and $\psi_T\in H^2(T)$ because $T$ is convex. The
computation is justified when $h$ is smooth up to the boundary. For general $h$ apply it on
the homothetic triangles $T_\varepsilon=I+(1-\varepsilon)(T-I)$, on whose closure $h$ is smooth,
with their own torsion functions, for which $C_{T_\varepsilon}=(1-\varepsilon)C_T$ by scaling;
then let $\varepsilon\to0$, using that $x\mapsto h(I+(1-\varepsilon)(x-I))$ converges to $h$ in
$H^1(T)$ and therefore its trace converges in $L^2(\partial T)$. For the bound on $C_T$ fix a
side, with altitude $h_j$ on it, and let $S$ be
the strip of width $h_j$ bounded by the line through that side and the parallel line through
the opposite vertex. Its torsion function $\psi_S(y)=y(h_j-y)/2$ satisfies
$\psi_S-\psi_T\ge0$ in $T$ by the maximum principle, and both vanish on the side, so
$|\partial_\nu\psi_T|\le|\partial_\nu\psi_S|=h_j/2$ there.
\end{proof}

\begin{proof}[Proof of Theorem~\ref{p-thm:large}]
Let $u$ be the Robin ground state, $\lambda=\lambda_1(T,\alpha)$, $g=u|_{\partial T}$. Restricting
\eqref{eq:var} to $H^1_0(T)$ gives $\lambda\le\Lambda_1$, and equality would make $\phi_1$ a Robin
eigenfunction, which it is not ($\partial_\nu\phi_1\not\equiv0=-\alpha\phi_1$ on $\partial T$); so
$\lambda<\Lambda_1$, and therefore $g\not\equiv0$. Write $u=h+v$ with
$h$ harmonic, $h=g$ on $\partial T$, and $v\in H^1_0(T)$. Then $-\Delta v-\lambda v=\lambda h$, so
$v=\lambda(-\Delta_D-\lambda)^{-1}h$, where $-\Delta_D$ is the Dirichlet Laplacian with
normalised eigenfunctions $\phi_k$ and eigenvalues $\Lambda_k$. By Green's formula and
$\int_T\nabla h\cdot\nabla v=0$,
\[
  -\alpha\int_{\partial T}g^2=\int_T|\nabla u|^2-\lambda\int_Tu^2
  =\int_T|\nabla h|^2-\lambda\|h\|^2-\lambda\langle h,v\rangle
  \ \ge\ -\lambda\|h\|^2-\lambda^2\sum_{k\ge1}\frac{\langle h,\phi_k\rangle^2}{\Lambda_k-\lambda},
\]
where norms and inner products without subscript are in $L^2(T)$. Bounding
$\Lambda_k-\lambda\ge\Lambda_2-\lambda$ for $k\ge2$ gives
\begin{equation}\label{p-eq:pf1}
  \alpha\int_{\partial T}g^2\ \le\ \lambda\|h\|^2\frac{\Lambda_2}{\Lambda_2-\lambda}
  +\lambda^2\frac{\langle h,\phi_1\rangle^2}{\Lambda_1-\lambda}.
\end{equation}
Since $h$ is harmonic and $\phi_1$ vanishes on $\partial T$,
$\Lambda_1\langle h,\phi_1\rangle=-\int_Th\,\Delta\phi_1=-\int_{\partial T}g\,\partial_\nu\phi_1$.
The Rellich identity $\int_{\partial T}(\partial_\nu\phi_1)^2\,(x-I)\cdot\nu\dd s=2\Lambda_1$
(valid on convex polygons, where $\phi_1\in H^2$) and $(x-I)\cdot\nu=r$ give
$\int_{\partial T}(\partial_\nu\phi_1)^2=2\Lambda_1/r$, so
$\langle h,\phi_1\rangle^2\le 2\,\|g\|_{L^2(\partial T)}^2/(r\Lambda_1)$. Inserting this and
Lemma~\ref{p-lem:harm} into \eqref{p-eq:pf1} and dividing by $\|g\|^2_{L^2(\partial T)}>0$ gives
\eqref{p-eq:large}.

For \eqref{p-eq:large2} use $\lambda<\Lambda_1$: the second term of \eqref{p-eq:large} is at most
$\alpha_*$ and the first is at most $2\Lambda_1/(r(\Lambda_1-\lambda))$.
\end{proof}

\begin{remark}\label{p-rem:sharp}
The first-order expansion $\lambda_1(\Omega,\alpha)=\Lambda_1-\alpha^{-1}\int_{\partial\Omega}
(\partial_\nu\phi_1)^2+o(\alpha^{-1})$ as $\alpha\to\infty$ is classical for smooth domains. If it
holds on $T$ (we have not located a reference for polygons and do not use it), then by the same
Rellich identity $\lambda_1(T,\alpha)=\Lambda_1(1-2/(r\alpha))+o(\alpha^{-1})$, so that
\eqref{p-eq:large2} is sharp to first order in $1/\alpha$.
\end{remark}
\section*{Declaration of AI-generated research and authorship}
This manuscript is the product of an AI-driven research process. The human author initiated a broader
investigation of open problems from the AIM problem lists, and AI systems worked on this problem within
it. AI systems produced all of the following: the results and proofs, the computations and the code for
the certificates, and the organisation and writing of the text. The systems were Claude Opus~5.5 and
Claude Fable~5.2 (Anthropic), used through Claude Code. The human author directed the AI systems but
wrote no part of the manuscript.

Theorem~\ref{thm:main} is computer-assisted and its proof is reproduced by \texttt{scripts/verify.py}.
At the time of writing, the mathematics and the code have been checked only by AI systems and by
computer verification. These checks include the certificates described above; adversarial reviews by
independent AI referees, of Part~I, of Part~II, and of the merged manuscript; a second certificate of
Lemma~\ref{lem:certC}, written by the referee of Part~I; an independent verification program for the
certificates of Part~II, written by a separate AI instance from the text and the certificate format
only; and an external AI review of the code (ChatGPT Pro), whose findings are corrected
(\S\ref{sec:verif}). The comparison with Ma's preprint~\cite{Ma} was also made by an AI referee. No independent human
expert has reviewed the manuscript. The human author has not read this manuscript or checked its
arguments. The author takes responsibility for the accuracy of this disclosure.

\raggedright
\begin{thebibliography}{99}
\bibitem{AimPL} American Institute of Mathematics, \emph{AIM Problem Lists: Shape optimization with
surface interactions}, Section~1, Conjecture~1.3 (V.~Lotoreichik, Z.~Lu),
\url{http://aimpl.org/shapesurface/1/}.
\bibitem{AFK} P.~R.~S.~Antunes, P.~Freitas, D.~Krej\v{c}i\v{r}\'{i}k, \emph{Bounds and extremal domains
for Robin eigenvalues with negative boundary parameter}, Adv.\ Calc.\ Var.\ \textbf{10} (2017), 357--379.
\href{https://doi.org/10.1515/acv-2015-0045}{doi:10.1515/acv-2015-0045}.
\bibitem{Bareket} M.~Bareket, \emph{On an isoperimetric inequality for the first eigenvalue of a boundary
value problem}, SIAM J.\ Math.\ Anal.\ \textbf{8} (1977), 280--287.
\href{https://doi.org/10.1137/0508020}{doi:10.1137/0508020}.
\bibitem{Bebendorf} M.~Bebendorf, \emph{A note on the Poincar\'e inequality for convex domains},
Z.\ Anal.\ Anwend.\ \textbf{22} (2003), 751--756.
\href{https://doi.org/10.4171/ZAA/1170}{doi:10.4171/ZAA/1170}.
\bibitem{Bossel} M.-H.~Bossel, \emph{Membranes \'elastiquement li\'ees: extension du th\'eor\`eme
de Rayleigh--Faber--Krahn et de l'in\'egalit\'e de Cheeger}, C.\ R.\ Acad.\ Sci.\ Paris S\'er.\ I Math.\
\textbf{302} (1986), 47--50.
\bibitem{BP} V.~Bruneau, N.~Popoff, \emph{On the negative spectrum of the Robin Laplacian in corner
domains}, Anal.\ PDE \textbf{9} (2016), 1259--1283.
\href{https://doi.org/10.2140/apde.2016.9.1259}{doi:10.2140/apde.2016.9.1259}.
\bibitem{BFK} D.~Bucur, P.~Freitas, J.~Kennedy, \emph{The Robin problem}, in: A.~Henrot (ed.), Shape
optimization and spectral theory, De Gruyter, 2017, 78--119.
\href{https://doi.org/10.1515/9783110550887-004}{doi:10.1515/9783110550887-004}.
\bibitem{CCP} A.~Carbotti, S.~Cito, D.~Pallara, \emph{On the optimization of the Robin eigenvalues in some
classes of polygons}, preprint (2025), arXiv:2508.03502.
\href{https://doi.org/10.48550/arXiv.2508.03502}{doi:10.48550/arXiv.2508.03502}.
\bibitem{CGal} C.~Carstensen, D.~Gallistl, \emph{Guaranteed lower eigenvalue bounds for the biharmonic
equation}, Numer.\ Math.\ \textbf{126} (2014), 33--51.
\href{https://doi.org/10.1007/s00211-013-0559-z}{doi:10.1007/s00211-013-0559-z}.
\bibitem{CG} C.~Carstensen, J.~Gedicke, \emph{Guaranteed lower bounds for eigenvalues}, Math.\ Comp.\
\textbf{83} (2014), 2605--2629.
\href{https://doi.org/10.1090/S0025-5718-2014-02833-0}{doi:10.1090/S0025-5718-2014-02833-0}.
\bibitem{CR} M.~Crouzeix, P.-A.~Raviart, \emph{Conforming and nonconforming finite element methods for
solving the stationary Stokes equations I}, Rev.\ Fran\c{c}aise Automat.\ Informat.\ Recherche
Op\'erationnelle S\'er.\ Rouge \textbf{7} (1973), no.~R-3, 33--75.
\href{https://doi.org/10.1051/m2an/197307R300331}{doi:10.1051/m2an/197307R300331}.
\bibitem{Daners} D.~Daners, \emph{A Faber--Krahn inequality for Robin problems in any space dimension},
Math.\ Ann.\ \textbf{335} (2006), 767--785.
\href{https://doi.org/10.1007/s00208-006-0753-8}{doi:10.1007/s00208-006-0753-8}.
\bibitem{ELM} R.~Endo, X.~Liu, P.~Mariano, \emph{Sharp Dirichlet eigenvalue inequalities on triangles},
preprint (2026), arXiv:2605.04331.
\href{https://doi.org/10.48550/arXiv.2605.04331}{doi:10.48550/arXiv.2605.04331}.
\bibitem{Ferudun} A.~Ferudun, \emph{Uniform weak Robin comparisons on triangles and a cubic moment
formula}, preprint, Zenodo (2026).
\href{https://doi.org/10.5281/zenodo.23404478}{doi:10.5281/zenodo.23404478}.
\bibitem{GS} T.~Giorgi, R.~Smits, \emph{Eigenvalue estimates and critical temperature in zero fields for
enhanced surface superconductivity}, Z.\ Angew.\ Math.\ Phys.\ \textbf{58} (2007), 224--245.
\href{https://doi.org/10.1007/s00033-005-0049-y}{doi:10.1007/s00033-005-0049-y}.
\bibitem{GHLY} C.~Gui, Y.~Hu, Q.~Li, H.~Yang, \emph{A quantitative Robin P\'olya--Szeg\H{o} theorem for
tangential polygons}, preprint (2026), arXiv:2610.21417.
\href{https://doi.org/10.48550/arXiv.2610.21417}{doi:10.48550/arXiv.2610.21417}.
\bibitem{HP} B.~Helffer, K.~Pankrashkin, \emph{Tunneling between corners for Robin Laplacians},
J.\ London Math.\ Soc.\ \textbf{91} (2015), 225--248.
\href{https://doi.org/10.1112/jlms/jdu070}{doi:10.1112/jlms/jdu070}.
\bibitem{Arb} F.~Johansson, \emph{Arb: efficient arbitrary-precision midpoint-radius interval
arithmetic}, IEEE Trans.\ Comput.\ \textbf{66} (2017), 1281--1292.
\href{https://doi.org/10.1109/TC.2017.2690633}{doi:10.1109/TC.2017.2690633}.
\bibitem{Khalile} M.~Khalile, \emph{Spectral asymptotics for Robin Laplacians on polygonal domains},
J.\ Math.\ Anal.\ Appl.\ \textbf{461} (2018), 1498--1543.
\href{https://doi.org/10.1016/j.jmaa.2018.01.062}{doi:10.1016/j.jmaa.2018.01.062}.
\bibitem{KP} M.~Khalile, K.~Pankrashkin, \emph{Eigenvalues of Robin Laplacians in infinite sectors},
Math.\ Nachr.\ \textbf{291} (2018), 928--965.
\href{https://doi.org/10.1002/mana.201600314}{doi:10.1002/mana.201600314}.
\bibitem{KLV} D.~Krej\v{c}i\v{r}\'{i}k, V.~Lotoreichik, T.~Vu, \emph{Reverse isoperimetric inequality for
the lowest Robin eigenvalue of a triangle}, Appl.\ Math.\ Optim.\ \textbf{88} (2023), article~63.
\href{https://doi.org/10.1007/s00245-023-10033-1}{doi:10.1007/s00245-023-10033-1}.
\bibitem{LOS} A.~A.~Lacey, J.~R.~Ockendon, J.~Sabina, \emph{Multidimensional reaction diffusion equations
with nonlinear boundary conditions}, SIAM J.\ Appl.\ Math.\ \textbf{58} (1998), 1622--1647.
\href{https://doi.org/10.1137/S0036139996308121}{doi:10.1137/S0036139996308121}.
\bibitem{LaSiPoincare} R.~S.~Laugesen, B.~A.~Siudeja, \emph{Minimizing Neumann fundamental tones of
triangles: an optimal Poincar\'e inequality}, J.\ Differential Equations \textbf{249} (2010), 118--135.
\href{https://doi.org/10.1016/j.jde.2010.02.020}{doi:10.1016/j.jde.2010.02.020}.
\bibitem{LaSi} R.~S.~Laugesen, B.~A.~Siudeja, \emph{Triangles and other special domains}, in:
A.~Henrot (ed.), Shape optimization and spectral theory, De Gruyter, 2017, 149--200.
\href{https://doi.org/10.1515/9783110550887-006}{doi:10.1515/9783110550887-006}.
\bibitem{LP} M.~Levitin, L.~Parnovski, \emph{On the principal eigenvalue of a Robin problem with a large
parameter}, Math.\ Nachr.\ \textbf{281} (2008), 272--281.
\href{https://doi.org/10.1002/mana.200510600}{doi:10.1002/mana.200510600}.
\bibitem{Liu} X.~Liu, \emph{A framework of verified eigenvalue bounds for self-adjoint differential
operators}, Appl.\ Math.\ Comput.\ \textbf{267} (2015), 341--355.
\href{https://doi.org/10.1016/j.amc.2015.03.048}{doi:10.1016/j.amc.2015.03.048}.
\bibitem{Ma} Y.~Ma, \emph{A proof of the extremal conjecture for the first Robin eigenvalue of
triangles}, Zenodo (deposited 6 October 2026; the PDF is titled \emph{A three-function comparison for
all triangles and all attractive Robin couplings}).
\href{https://doi.org/10.5281/zenodo.23191831}{doi:10.5281/zenodo.23191831}.
\bibitem{McCartin2} B.~J.~McCartin, \emph{Eigenstructure of the equilateral triangle, Part II: The
Neumann problem}, Math.\ Probl.\ Eng.\ \textbf{8} (2002), 517--539.
\href{https://doi.org/10.1080/1024123021000053664}{doi:10.1080/1024123021000053664}.
\bibitem{McCartin} B.~J.~McCartin, \emph{Eigenstructure of the equilateral triangle. Part III. The Robin
problem}, Int.\ J.\ Math.\ Math.\ Sci.\ \textbf{2004} (2004), 807--825.
\href{https://doi.org/10.1155/S0161171204306125}{doi:10.1155/S0161171204306125}.
\bibitem{mpmath} The mpmath development team, \emph{mpmath: a Python library for arbitrary-precision
floating-point arithmetic}, version~1.3.0 (2023), \url{https://mpmath.org}.
\bibitem{PW} L.~E.~Payne, H.~F.~Weinberger, \emph{An optimal Poincar\'e inequality for convex domains},
Arch.\ Rational Mech.\ Anal.\ \textbf{5} (1960), 286--292.
\href{https://doi.org/10.1007/BF00252910}{doi:10.1007/BF00252910}.
\bibitem{PS} G.~P\'olya, G.~Szeg\H{o}, \emph{Isoperimetric Inequalities in Mathematical Physics}, Annals
of Mathematics Studies~27, Princeton University Press, 1951.
\href{https://doi.org/10.1515/9781400882663}{doi:10.1515/9781400882663}.
\end{thebibliography}

\end{document}
